% 高一40_大同中学_周末作业4.2.tex
%   —— 高一上数学“2026-2027 年大同中学高一上周末作业 4.2”（试卷版/答案版）
% 试卷版：xelatex 高一40_大同中学_周末作业4.2.tex
% 答案版：xelatex "\def\isanswer{1}% 高一40_大同中学_周末作业4.2.tex
%   —— 高一上数学“2026-2027 年大同中学高一上周末作业 4.2”（试卷版/答案版）
% 试卷版：xelatex 高一40_大同中学_周末作业4.2.tex
% 答案版：xelatex "\def\isanswer{1}% 高一40_大同中学_周末作业4.2.tex
%   —— 高一上数学“2026-2027 年大同中学高一上周末作业 4.2”（试卷版/答案版）
% 试卷版：xelatex 高一40_大同中学_周末作业4.2.tex
% 答案版：xelatex "\def\isanswer{1}% 高一40_大同中学_周末作业4.2.tex
%   —— 高一上数学“2026-2027 年大同中学高一上周末作业 4.2”（试卷版/答案版）
% 试卷版：xelatex 高一40_大同中学_周末作业4.2.tex
% 答案版：xelatex "\def\isanswer{1}\input{高一40_大同中学_周末作业4.2.tex}"
% 紧凑版：xelatex "\def\iscompact{1}\input{高一40_大同中学_周末作业4.2.tex}"
%
% 原始材料是微信公众号图文（公众号“魔都数学加油站”连载“27学年高一 40”，署名“破晓”，
% 2026-09-26 21:29 发布，原文标题“27学年高一40——2026-2027年上海市大同中学高一上周末作业4.2”，
% 原文链接 https://mp.weixin.qq.com/s/RrBSqDF_nJVbbx-bmmdtkQ ），正文为 4 张 PNG 页（4762×6735）。原文本身就是一份**讲评版**——题下直接印【答案】或【解析】，
% 第 11、13、14 题的答案还直接印在题干末尾的括号里。试题文字、答案与解析均照原卷录入，
% 未改内容；卷面的粉色斜水印（“魔都数学加油站”字样）属水印，未录入。
%
% 卷首标题：原卷页首印作“2026-2027 年大同中学高一上周末作业 4.2”（**卷面标题行即带校名**，
% 与高一16／高一18 等拍照卷须由本稿补校名的情形不同），本稿照卷面，年份区间按全稿惯例
% 排 en dash，前缀照原卷保留“年”。
%
% 篇幅与题号：本篇**只有 4 页**，卷面自“一、填空题”的**第 3 题**起编，终于“三、解答题”的
% **第 19 题**，共 17 题（填空 3—10、选择 11—14、解答 15—19）。第 1、2 题未在该文中发布——
% 这是该公众号连载讲评版的通例（高一02—高一09 与 高一36、高一38、高一39 等均自第 3 题起编，
% README“说明”已记），**不是截断**，故本稿题号亦自 3 起。它与 高一39（大同中学 周末作业4.1）
% **不构成同一份试卷的两半**：4.1 是另一次作业（11 页，题号**也自第 3 题**起、一直编到**第 24 题**，
% 以集合／常用逻辑用语为主、兼及不等式），两份各自独立起编（都从第 3 题开始），只是同为“大同中学高一上
% 周末作业 4”系列。
%
% 与原卷的排版差异（均已在 README“说明”中交代）：
%   ① 原文自**第 3 题**起（第 1、2 题未在该文中发布），故本稿题号亦自 3 起；
%   ② 原卷本就印有“一、填空题”“二、选择题”“三、解答题”三节标题，本稿照原卷分节，
%      只把标题改排成全稿统一的“大节标题”样式；
%   ③ 原卷为讲评版：第 4、5、6、7、8、17、18 题只印【答案】行，第 3、9、10、12、15、16、19 题
%      印【解析】（无单独的【答案】行），第 11、13、14 题的答案直接印在题干末尾的括号内；
%      本稿统一改为试卷版隐去、答案版以 \ans{}／\ifanswer 块给出，两版彻底分离（照 高一36 口径：
%      只印【解析】的题不另给 \ans{}，只印【答案】的题仅给 \ans{}）；
%   ④ 数集记号原卷作斜体的 $R$（第 5、10、13 题），本稿按全稿惯例统一写作 $\R$；
%   ⑤ 补集符号原卷作上横线（第 5 题的 $\overline{M}$），本稿照录，未改为 $\complement_\R$；
%   ⑥ 包含符号原卷两套并用：第 4 题作**无下划线**的 $\subset$（真包含），第 18 题作**带下划线**的
%      $\sub$（即 $\subseteq$），本稿照录、不归一；
%   ⑦ 原卷填空的横线约两个字宽，本稿按全稿惯例统一排 \blankline[4.4em]；
%   ⑧ 原卷未印副标题，本稿按**本卷实际内容**补出副标题“不等式”——全卷 17 题里约 15 题为
%      不等式（解一元二次／无理／分式不等式、恒成立与根的分布、含参不等式、应用题），
%      集合仅第 4、5 两题（第 18 题只是用集合关系表述解集），常用逻辑用语 0 题；
%   ⑨ 本卷**无插图**（全卷检索确认无“如图”）；
%   ⑩ 标点与单位：第 6 题的“$a,b$”本稿按全稿惯例排顿号“$a$、$b$”，第 10 题的三例
%      “$[\pi]=3,[-1.2]=-2,[\tfrac12]=0$”之间排顿号，第 5 题原卷“全集 $U=R$ 集合 $M$”之间
%      漏排标点、本稿补逗号；车速单位原卷印作斜体 $km/h$（第 19 题），本稿照录；半角括号、逗号、
%      问号一律排全角（第 17、19 题原卷的两个问号本就作全角，本稿照录）。
%
% 原卷存疑两处（照抄原卷、未改，见源码中该处上方注释）：
%   ① 第 13 题的 D 选项与 A 选项文字完全重复（均作“$a<0$，且 $b^2-4ac>0$”），原卷如此。
%   ② 第 10 题题干末作“…的取值范围____”，**漏排“是”**（本卷其余填空均作“…是____”）——
%      本稿照录、未补字，见源码该题上方 `%` 注。

\documentclass[a4paper,11pt]{article}
\usepackage{common}

\begin{document}

\examtitlest[3]{2026--2027 年大同中学高一上周末作业 4.2}{不等式}

\bigsec{一、填空题}

\begin{enumerate}
\setcounter{enumi}{2}

\item 在平面直角坐标系中，若点 $(2,b)$ 到原点的距离不小于 $5$，则 $b$ 的取值范围是 \blankline[4.4em].

\ifanswer
\solbegin
\step{因为点 $(2,b)$ 到原点的距离不小于 $5$，所以 $\sqrt{(2-0)^2+(b-0)^2}\geqslant5$，}
\step{所以 $4+b^2\geqslant25$，得 $b\geqslant\sqrt{21}$ 或 $b\leqslant-\sqrt{21}$，}
\step{则 $b$ 的取值范围是 $(-\infty,-\sqrt{21}]\cup[\sqrt{21},+\infty)$.}
\solend

\fi

\item 若集合 $A=\{x\mid|x|>3\}$，集合 $B=\{x\mid x\geqslant a\}$，且 $B\subset A$，
则实数 $a$ 的取值范围是 \blankline[4.4em].

\ans{$(3,+\infty)$}

% 注：原卷“全集 $U=R$”与“集合 $M$”之间**漏排标点**（只有空格），本稿补逗号——
%     照 高一b 第 13 题“原卷漏排、本稿补上”的先例（高一10、高一24 同例），见文件头⑩。
\item 若全集 $U=\R$，集合 $M=\{x\mid y=\sqrt{9-x^2}\}$，$S=\{x\mid x^2-3x+2\geqslant0\}$，
则 $\overline{M}\cup S=$ \blankline[4.4em].

\ans{$(-\infty,1]\cup[2,+\infty)$}

\item 设 $a$、$b$ 为实数，若不等式 $x^2-ax-b<0$ 的解集为 $\{x\mid2<x<8\}$，
则 $bx^2-ax-1>0$ 的解集为 \blankline[4.4em].

\ans{$\left(-\dfrac12,-\dfrac18\right)$}

\item 不等式 $4x-13+\sqrt{x+5}\geqslant8x-1+\sqrt{x+5}$ 的解集是 \blankline[4.4em].

\ans{$[-5,-3]$}

\item 若关于 $x$ 的不等式 $ax^2-2ax+3\leqslant0$ 的解集为 $\varnothing$，
则实数 $a$ 的取值范围是 \blankline[4.4em].

\ans{$[0,3)$}

\item 若不等式 $x^2-ax-2<0$ 在 $x\in(1,2)$ 内恒成立，则实数 $a$ 的取值范围是 \blankline[4.4em].

\ifanswer
\solbegin
\step{由不等式 $x^2-ax-2<0$ 在 $x\in(1,2)$ 内恒成立，}
\step{得 $ax>x^2-2$，即 $a>x-\dfrac2x$ 在 $x\in(1,2)$ 内恒成立，}
\step{令 $t(x)=x-\dfrac2x$，$x\in(1,2)$，该函数为严格增函数，则 $t(x)<t(2)=1$，得 $a\geqslant1$，}
\step{所以 $a$ 的取值范围是 $[1,+\infty)$.}
\solend

\fi

% 注：原卷此处作“的取值范围\_\_\_\_”，**漏排“是”**（本卷其余填空均作“…是\_\_\_\_”），
%     本稿照录、未补字，见文件头“原卷存疑”②。
\item 设 $x\in\R$，$[x]$ 表示不大于 $x$ 的最大整数，如 $[\pi]=3$、$[-1.2]=-2$、$\left[\dfrac12\right]=0$，
则使 $[|x-1|]=3$ 成立的 $x$ 的取值范围 \blankline[4.4em].

\ifanswer
\solbegin
\step{由题意 $[|x-1|]=3$，则 $3\leqslant|x-1|<4$，所以 $3\leqslant x-1<4$ 或 $-4\leqslant x-1<-3$，}
\step{解得 $4\leqslant x<5$ 或 $-3<x\leqslant-2$，}
\step{所以使 $[|x-1|]=3$ 成立的 $x$ 的取值范围是 $(-3,-2]\cup[4,5)$.}
\solend

\fi

\end{enumerate}

\bigsec{二、选择题}

\begin{enumerate}
\setcounter{enumi}{10}

\item 若关于 $x$ 的方程 $k(x+2)-3(k-1)=x+1$ 有负数解，则实数 $k$ 的取值范围是\mbox{（\quad）}.

\keepwithprev
A. $(1,2)$\qquad B. $(-\infty,1)\cup(2,+\infty)$\qquad C. $(-\infty,2)$\qquad D. $(-\infty,1)$

\ans{$A$}

\item 与 $x^2>4$ 同解的不等式是\mbox{（\quad）}.

\keepwithprev
A. $x^2+\dfrac{1}{x-4}>4+\dfrac{1}{x-4}$\qquad B. $x^2+\dfrac1x>4+\dfrac1x$

C. $x^2+\sqrt{x-2}>4+\sqrt{x-2}$\qquad D. $\dfrac{1}{x^2-4}<0$

\ifanswer
\solbegin
\step{不等式 $x^2+\dfrac{1}{x-4}>4+\dfrac{1}{x-4}$ 等价于 $\begin{cases}x^2>4\\ x\neq4\end{cases}$，
与 $x^2>4$ 不同解，故选项 $A$ 错误；}
\step{不等式 $x^2+\dfrac1x>4+\dfrac1x$ 等价于 $\begin{cases}x^2>4\\ x\neq0\end{cases}$，
即 $x^2>4$，所以与 $x^2>4$ 同解，}
\step{故选项 $B$ 正确；}
\step{不等式 $x^2+\sqrt{x-2}>4+\sqrt{x-2}$ 等价于 $\begin{cases}x^2>4\\ x-2\geqslant0\end{cases}$，
与 $x^2>4$ 不同解，}
\step{故选项 $C$ 错误；}
\step{不等式 $\dfrac{1}{x^2-4}<0$ 等价于 $x^2<4$，与 $x^2>4$ 不同解，故选项 $D$ 错误.}
\step{故选 $B$.}
\solend

\fi

\item 若关于 $x$ 的二次不等式 $ax^2+bx+c\geqslant0(a\neq0)$ 的解集是 $\R$，那么\mbox{（\quad）}.

\keepwithprev
A. $a<0$，且 $b^2-4ac>0$\qquad B. $a<0$，且 $b^2-4ac\leqslant0$

% 注：原卷本题 D 选项与 A 选项文字完全重复（均作“$a<0$，且 $b^2-4ac>0$”），
%     照录未改，见文件头“原卷存疑”①。
C. $a>0$，且 $b^2-4ac\leqslant0$\qquad D. $a<0$，且 $b^2-4ac>0$

\ans{$C$}

\item 若不等式 $x^2+mx+1>2x+m$ 对满足 $|m|<2$ 的所有实数 $m$ 恒成立，
则实数 $x$ 的取值范围是\mbox{（\quad）}.

\keepwithprev
A. $-2<x<2$\qquad B. $x\geqslant3$\qquad C. $x\leqslant1$\qquad D. $x\leqslant-1$ 或 $x\geqslant3$

\ans{$D$}

\end{enumerate}

% “三、解答题”原本落在试卷版第 1 页页末、第 15 题翻到次页（切页审计 B 类）。
% 节标题自带的守卫只有 3\baselineskip，够放标题不够放标题＋首题，故在标题前再量一次页高。
\needvspace{4cm}
\bigsec{三、解答题}

\begin{enumerate}
\setcounter{enumi}{14}

\item 设关于 $x$ 的二次方程 $(m-1)x^2+(2m-4)x+m=0$ 有两个正根，求实数 $m$ 的取值范围.

\ifanswer
\solbegin
\step{$\begin{cases}\Delta=(2m-4)^2-4m(m-1)\geqslant0\\ -\dfrac{2m-4}{m-1}>0\\ \dfrac{m}{m-1}>0\end{cases}$，
解得 $1<m\leqslant\dfrac43$.}
\solend

\fi

\ansspace[6]

\item 已知关于 $x$ 的不等式 $2x^2-9x+a<0$ 的解集非空，对于其解集内的每一个 $x$ 的值，
至少能使不等式 $x^2-4x+3<0$ 或 $x^2-6x+8<0$ 中的一个成立，求实数 $a$ 的取值范围.

\ifanswer
\solbegin
\step{不等式 $x^2-4x+3<0$ 的解集为 $(1,3)$；$x^2-6x+8<0$ 的解集为 $(2,4)$.}
\step{关于 $x$ 的不等式 $2x^2-9x+a<0$ 的解集非空，所以 $\Delta=81-8a>0$，解得 $a<\dfrac{81}{8}$.}
\step{所以 $\dfrac{9-\sqrt{81-8a}}{4}<x<\dfrac{9+\sqrt{81-8a}}{4}$.}
\step{由于关于 $x$ 的不等式 $2x^2-9x+a<0$（解集非空）的每一个 $x$ 的值}
\step{至少满足不等式 $x^2-4x+3<0$ 和 $x^2-6x+8<0$ 中的一个，}
\step{所以 $\begin{cases}\dfrac{9-\sqrt{81-8a}}{4}\geqslant1\\ \dfrac{9+\sqrt{81-8a}}{4}\leqslant4\\ a<\dfrac{81}{8}\end{cases}$，
解得 $7\leqslant a<\dfrac{81}{8}$，所以实数 $a$ 的取值范围是 $\left[7,\dfrac{81}{8}\right)$.}
\solend

\fi

\ansspace[8]

\item 是否存在实数 $k$，使 $x=3$ 在不等式 $\dfrac{x+2}{k}>1+\dfrac{x^2-3}{k^2}$ 的解集中？
若存在，求 $k$ 的取值范围；若不存在，说明理由.

\ans{$(2,3)$}

\ansspace[6]

\item 关于 $x$ 的不等式 $x^2-(a+1)x+a<0$ 的解集为 $A$，$x^2-5x+4<0$ 的解集为 $B$，
若 $A\sub B$，求实数 $a$ 的取值范围.

\ans{$[1,4]$}

\ansspace[6]

\item 汽车在行驶中，由于惯性作用，刹车后还要继续向前滑行一段距离才能停住，我们称这段
距离为“刹车距离”. 刹车距离是分析事故的一个重要因素. 在一个限速 $40km/h$ 以内的弯道上，
甲、乙两辆汽车相向而行，发现情况不对，同时刹车，但还是相撞了，事后现场测得甲车的刹车
距离略超过 $12m$，乙车的刹车距离略超过 $10m$，又知甲、乙两种车型的刹车距离 $s(m)$ 与
车速 $x(km/h)$ 之间有如下关系：$s_{\text{甲}}=0.1x+0.01x^2$，$s_{\text{乙}}=0.05x+0.005x^2$. 问：
两车相撞的主要责任是谁？

\ifanswer
\solbegin
\step{由题意列出不等式组 $\begin{cases}0.1x+0.01x^2>12\\ 0.05x+0.005x^2>10\end{cases}$，
得 $\begin{cases}x<-40\text{ 或 }x>30\\ x<-50\text{ 或 }x>40\end{cases}$，}
\step{由于 $x>0$，从而得 $x_{\text{甲}}>30km/h$，$x_{\text{乙}}>40km/h$.}
\step{经比较得乙车超过限速，应负主要责任.}
\solend

\fi

\ansspace*[8]

\end{enumerate}

\end{document}
"
% 紧凑版：xelatex "\def\iscompact{1}% 高一40_大同中学_周末作业4.2.tex
%   —— 高一上数学“2026-2027 年大同中学高一上周末作业 4.2”（试卷版/答案版）
% 试卷版：xelatex 高一40_大同中学_周末作业4.2.tex
% 答案版：xelatex "\def\isanswer{1}\input{高一40_大同中学_周末作业4.2.tex}"
% 紧凑版：xelatex "\def\iscompact{1}\input{高一40_大同中学_周末作业4.2.tex}"
%
% 原始材料是微信公众号图文（公众号“魔都数学加油站”连载“27学年高一 40”，署名“破晓”，
% 2026-09-26 21:29 发布，原文标题“27学年高一40——2026-2027年上海市大同中学高一上周末作业4.2”，
% 原文链接 https://mp.weixin.qq.com/s/RrBSqDF_nJVbbx-bmmdtkQ ），正文为 4 张 PNG 页（4762×6735）。原文本身就是一份**讲评版**——题下直接印【答案】或【解析】，
% 第 11、13、14 题的答案还直接印在题干末尾的括号里。试题文字、答案与解析均照原卷录入，
% 未改内容；卷面的粉色斜水印（“魔都数学加油站”字样）属水印，未录入。
%
% 卷首标题：原卷页首印作“2026-2027 年大同中学高一上周末作业 4.2”（**卷面标题行即带校名**，
% 与高一16／高一18 等拍照卷须由本稿补校名的情形不同），本稿照卷面，年份区间按全稿惯例
% 排 en dash，前缀照原卷保留“年”。
%
% 篇幅与题号：本篇**只有 4 页**，卷面自“一、填空题”的**第 3 题**起编，终于“三、解答题”的
% **第 19 题**，共 17 题（填空 3—10、选择 11—14、解答 15—19）。第 1、2 题未在该文中发布——
% 这是该公众号连载讲评版的通例（高一02—高一09 与 高一36、高一38、高一39 等均自第 3 题起编，
% README“说明”已记），**不是截断**，故本稿题号亦自 3 起。它与 高一39（大同中学 周末作业4.1）
% **不构成同一份试卷的两半**：4.1 是另一次作业（11 页，题号**也自第 3 题**起、一直编到**第 24 题**，
% 以集合／常用逻辑用语为主、兼及不等式），两份各自独立起编（都从第 3 题开始），只是同为“大同中学高一上
% 周末作业 4”系列。
%
% 与原卷的排版差异（均已在 README“说明”中交代）：
%   ① 原文自**第 3 题**起（第 1、2 题未在该文中发布），故本稿题号亦自 3 起；
%   ② 原卷本就印有“一、填空题”“二、选择题”“三、解答题”三节标题，本稿照原卷分节，
%      只把标题改排成全稿统一的“大节标题”样式；
%   ③ 原卷为讲评版：第 4、5、6、7、8、17、18 题只印【答案】行，第 3、9、10、12、15、16、19 题
%      印【解析】（无单独的【答案】行），第 11、13、14 题的答案直接印在题干末尾的括号内；
%      本稿统一改为试卷版隐去、答案版以 \ans{}／\ifanswer 块给出，两版彻底分离（照 高一36 口径：
%      只印【解析】的题不另给 \ans{}，只印【答案】的题仅给 \ans{}）；
%   ④ 数集记号原卷作斜体的 $R$（第 5、10、13 题），本稿按全稿惯例统一写作 $\R$；
%   ⑤ 补集符号原卷作上横线（第 5 题的 $\overline{M}$），本稿照录，未改为 $\complement_\R$；
%   ⑥ 包含符号原卷两套并用：第 4 题作**无下划线**的 $\subset$（真包含），第 18 题作**带下划线**的
%      $\sub$（即 $\subseteq$），本稿照录、不归一；
%   ⑦ 原卷填空的横线约两个字宽，本稿按全稿惯例统一排 \blankline[4.4em]；
%   ⑧ 原卷未印副标题，本稿按**本卷实际内容**补出副标题“不等式”——全卷 17 题里约 15 题为
%      不等式（解一元二次／无理／分式不等式、恒成立与根的分布、含参不等式、应用题），
%      集合仅第 4、5 两题（第 18 题只是用集合关系表述解集），常用逻辑用语 0 题；
%   ⑨ 本卷**无插图**（全卷检索确认无“如图”）；
%   ⑩ 标点与单位：第 6 题的“$a,b$”本稿按全稿惯例排顿号“$a$、$b$”，第 10 题的三例
%      “$[\pi]=3,[-1.2]=-2,[\tfrac12]=0$”之间排顿号，第 5 题原卷“全集 $U=R$ 集合 $M$”之间
%      漏排标点、本稿补逗号；车速单位原卷印作斜体 $km/h$（第 19 题），本稿照录；半角括号、逗号、
%      问号一律排全角（第 17、19 题原卷的两个问号本就作全角，本稿照录）。
%
% 原卷存疑两处（照抄原卷、未改，见源码中该处上方注释）：
%   ① 第 13 题的 D 选项与 A 选项文字完全重复（均作“$a<0$，且 $b^2-4ac>0$”），原卷如此。
%   ② 第 10 题题干末作“…的取值范围____”，**漏排“是”**（本卷其余填空均作“…是____”）——
%      本稿照录、未补字，见源码该题上方 `%` 注。

\documentclass[a4paper,11pt]{article}
\usepackage{common}

\begin{document}

\examtitlest[3]{2026--2027 年大同中学高一上周末作业 4.2}{不等式}

\bigsec{一、填空题}

\begin{enumerate}
\setcounter{enumi}{2}

\item 在平面直角坐标系中，若点 $(2,b)$ 到原点的距离不小于 $5$，则 $b$ 的取值范围是 \blankline[4.4em].

\ifanswer
\solbegin
\step{因为点 $(2,b)$ 到原点的距离不小于 $5$，所以 $\sqrt{(2-0)^2+(b-0)^2}\geqslant5$，}
\step{所以 $4+b^2\geqslant25$，得 $b\geqslant\sqrt{21}$ 或 $b\leqslant-\sqrt{21}$，}
\step{则 $b$ 的取值范围是 $(-\infty,-\sqrt{21}]\cup[\sqrt{21},+\infty)$.}
\solend

\fi

\item 若集合 $A=\{x\mid|x|>3\}$，集合 $B=\{x\mid x\geqslant a\}$，且 $B\subset A$，
则实数 $a$ 的取值范围是 \blankline[4.4em].

\ans{$(3,+\infty)$}

% 注：原卷“全集 $U=R$”与“集合 $M$”之间**漏排标点**（只有空格），本稿补逗号——
%     照 高一b 第 13 题“原卷漏排、本稿补上”的先例（高一10、高一24 同例），见文件头⑩。
\item 若全集 $U=\R$，集合 $M=\{x\mid y=\sqrt{9-x^2}\}$，$S=\{x\mid x^2-3x+2\geqslant0\}$，
则 $\overline{M}\cup S=$ \blankline[4.4em].

\ans{$(-\infty,1]\cup[2,+\infty)$}

\item 设 $a$、$b$ 为实数，若不等式 $x^2-ax-b<0$ 的解集为 $\{x\mid2<x<8\}$，
则 $bx^2-ax-1>0$ 的解集为 \blankline[4.4em].

\ans{$\left(-\dfrac12,-\dfrac18\right)$}

\item 不等式 $4x-13+\sqrt{x+5}\geqslant8x-1+\sqrt{x+5}$ 的解集是 \blankline[4.4em].

\ans{$[-5,-3]$}

\item 若关于 $x$ 的不等式 $ax^2-2ax+3\leqslant0$ 的解集为 $\varnothing$，
则实数 $a$ 的取值范围是 \blankline[4.4em].

\ans{$[0,3)$}

\item 若不等式 $x^2-ax-2<0$ 在 $x\in(1,2)$ 内恒成立，则实数 $a$ 的取值范围是 \blankline[4.4em].

\ifanswer
\solbegin
\step{由不等式 $x^2-ax-2<0$ 在 $x\in(1,2)$ 内恒成立，}
\step{得 $ax>x^2-2$，即 $a>x-\dfrac2x$ 在 $x\in(1,2)$ 内恒成立，}
\step{令 $t(x)=x-\dfrac2x$，$x\in(1,2)$，该函数为严格增函数，则 $t(x)<t(2)=1$，得 $a\geqslant1$，}
\step{所以 $a$ 的取值范围是 $[1,+\infty)$.}
\solend

\fi

% 注：原卷此处作“的取值范围\_\_\_\_”，**漏排“是”**（本卷其余填空均作“…是\_\_\_\_”），
%     本稿照录、未补字，见文件头“原卷存疑”②。
\item 设 $x\in\R$，$[x]$ 表示不大于 $x$ 的最大整数，如 $[\pi]=3$、$[-1.2]=-2$、$\left[\dfrac12\right]=0$，
则使 $[|x-1|]=3$ 成立的 $x$ 的取值范围 \blankline[4.4em].

\ifanswer
\solbegin
\step{由题意 $[|x-1|]=3$，则 $3\leqslant|x-1|<4$，所以 $3\leqslant x-1<4$ 或 $-4\leqslant x-1<-3$，}
\step{解得 $4\leqslant x<5$ 或 $-3<x\leqslant-2$，}
\step{所以使 $[|x-1|]=3$ 成立的 $x$ 的取值范围是 $(-3,-2]\cup[4,5)$.}
\solend

\fi

\end{enumerate}

\bigsec{二、选择题}

\begin{enumerate}
\setcounter{enumi}{10}

\item 若关于 $x$ 的方程 $k(x+2)-3(k-1)=x+1$ 有负数解，则实数 $k$ 的取值范围是\mbox{（\quad）}.

\keepwithprev
A. $(1,2)$\qquad B. $(-\infty,1)\cup(2,+\infty)$\qquad C. $(-\infty,2)$\qquad D. $(-\infty,1)$

\ans{$A$}

\item 与 $x^2>4$ 同解的不等式是\mbox{（\quad）}.

\keepwithprev
A. $x^2+\dfrac{1}{x-4}>4+\dfrac{1}{x-4}$\qquad B. $x^2+\dfrac1x>4+\dfrac1x$

C. $x^2+\sqrt{x-2}>4+\sqrt{x-2}$\qquad D. $\dfrac{1}{x^2-4}<0$

\ifanswer
\solbegin
\step{不等式 $x^2+\dfrac{1}{x-4}>4+\dfrac{1}{x-4}$ 等价于 $\begin{cases}x^2>4\\ x\neq4\end{cases}$，
与 $x^2>4$ 不同解，故选项 $A$ 错误；}
\step{不等式 $x^2+\dfrac1x>4+\dfrac1x$ 等价于 $\begin{cases}x^2>4\\ x\neq0\end{cases}$，
即 $x^2>4$，所以与 $x^2>4$ 同解，}
\step{故选项 $B$ 正确；}
\step{不等式 $x^2+\sqrt{x-2}>4+\sqrt{x-2}$ 等价于 $\begin{cases}x^2>4\\ x-2\geqslant0\end{cases}$，
与 $x^2>4$ 不同解，}
\step{故选项 $C$ 错误；}
\step{不等式 $\dfrac{1}{x^2-4}<0$ 等价于 $x^2<4$，与 $x^2>4$ 不同解，故选项 $D$ 错误.}
\step{故选 $B$.}
\solend

\fi

\item 若关于 $x$ 的二次不等式 $ax^2+bx+c\geqslant0(a\neq0)$ 的解集是 $\R$，那么\mbox{（\quad）}.

\keepwithprev
A. $a<0$，且 $b^2-4ac>0$\qquad B. $a<0$，且 $b^2-4ac\leqslant0$

% 注：原卷本题 D 选项与 A 选项文字完全重复（均作“$a<0$，且 $b^2-4ac>0$”），
%     照录未改，见文件头“原卷存疑”①。
C. $a>0$，且 $b^2-4ac\leqslant0$\qquad D. $a<0$，且 $b^2-4ac>0$

\ans{$C$}

\item 若不等式 $x^2+mx+1>2x+m$ 对满足 $|m|<2$ 的所有实数 $m$ 恒成立，
则实数 $x$ 的取值范围是\mbox{（\quad）}.

\keepwithprev
A. $-2<x<2$\qquad B. $x\geqslant3$\qquad C. $x\leqslant1$\qquad D. $x\leqslant-1$ 或 $x\geqslant3$

\ans{$D$}

\end{enumerate}

% “三、解答题”原本落在试卷版第 1 页页末、第 15 题翻到次页（切页审计 B 类）。
% 节标题自带的守卫只有 3\baselineskip，够放标题不够放标题＋首题，故在标题前再量一次页高。
\needvspace{4cm}
\bigsec{三、解答题}

\begin{enumerate}
\setcounter{enumi}{14}

\item 设关于 $x$ 的二次方程 $(m-1)x^2+(2m-4)x+m=0$ 有两个正根，求实数 $m$ 的取值范围.

\ifanswer
\solbegin
\step{$\begin{cases}\Delta=(2m-4)^2-4m(m-1)\geqslant0\\ -\dfrac{2m-4}{m-1}>0\\ \dfrac{m}{m-1}>0\end{cases}$，
解得 $1<m\leqslant\dfrac43$.}
\solend

\fi

\ansspace[6]

\item 已知关于 $x$ 的不等式 $2x^2-9x+a<0$ 的解集非空，对于其解集内的每一个 $x$ 的值，
至少能使不等式 $x^2-4x+3<0$ 或 $x^2-6x+8<0$ 中的一个成立，求实数 $a$ 的取值范围.

\ifanswer
\solbegin
\step{不等式 $x^2-4x+3<0$ 的解集为 $(1,3)$；$x^2-6x+8<0$ 的解集为 $(2,4)$.}
\step{关于 $x$ 的不等式 $2x^2-9x+a<0$ 的解集非空，所以 $\Delta=81-8a>0$，解得 $a<\dfrac{81}{8}$.}
\step{所以 $\dfrac{9-\sqrt{81-8a}}{4}<x<\dfrac{9+\sqrt{81-8a}}{4}$.}
\step{由于关于 $x$ 的不等式 $2x^2-9x+a<0$（解集非空）的每一个 $x$ 的值}
\step{至少满足不等式 $x^2-4x+3<0$ 和 $x^2-6x+8<0$ 中的一个，}
\step{所以 $\begin{cases}\dfrac{9-\sqrt{81-8a}}{4}\geqslant1\\ \dfrac{9+\sqrt{81-8a}}{4}\leqslant4\\ a<\dfrac{81}{8}\end{cases}$，
解得 $7\leqslant a<\dfrac{81}{8}$，所以实数 $a$ 的取值范围是 $\left[7,\dfrac{81}{8}\right)$.}
\solend

\fi

\ansspace[8]

\item 是否存在实数 $k$，使 $x=3$ 在不等式 $\dfrac{x+2}{k}>1+\dfrac{x^2-3}{k^2}$ 的解集中？
若存在，求 $k$ 的取值范围；若不存在，说明理由.

\ans{$(2,3)$}

\ansspace[6]

\item 关于 $x$ 的不等式 $x^2-(a+1)x+a<0$ 的解集为 $A$，$x^2-5x+4<0$ 的解集为 $B$，
若 $A\sub B$，求实数 $a$ 的取值范围.

\ans{$[1,4]$}

\ansspace[6]

\item 汽车在行驶中，由于惯性作用，刹车后还要继续向前滑行一段距离才能停住，我们称这段
距离为“刹车距离”. 刹车距离是分析事故的一个重要因素. 在一个限速 $40km/h$ 以内的弯道上，
甲、乙两辆汽车相向而行，发现情况不对，同时刹车，但还是相撞了，事后现场测得甲车的刹车
距离略超过 $12m$，乙车的刹车距离略超过 $10m$，又知甲、乙两种车型的刹车距离 $s(m)$ 与
车速 $x(km/h)$ 之间有如下关系：$s_{\text{甲}}=0.1x+0.01x^2$，$s_{\text{乙}}=0.05x+0.005x^2$. 问：
两车相撞的主要责任是谁？

\ifanswer
\solbegin
\step{由题意列出不等式组 $\begin{cases}0.1x+0.01x^2>12\\ 0.05x+0.005x^2>10\end{cases}$，
得 $\begin{cases}x<-40\text{ 或 }x>30\\ x<-50\text{ 或 }x>40\end{cases}$，}
\step{由于 $x>0$，从而得 $x_{\text{甲}}>30km/h$，$x_{\text{乙}}>40km/h$.}
\step{经比较得乙车超过限速，应负主要责任.}
\solend

\fi

\ansspace*[8]

\end{enumerate}

\end{document}
"
%
% 原始材料是微信公众号图文（公众号“魔都数学加油站”连载“27学年高一 40”，署名“破晓”，
% 2026-09-26 21:29 发布，原文标题“27学年高一40——2026-2027年上海市大同中学高一上周末作业4.2”，
% 原文链接 https://mp.weixin.qq.com/s/RrBSqDF_nJVbbx-bmmdtkQ ），正文为 4 张 PNG 页（4762×6735）。原文本身就是一份**讲评版**——题下直接印【答案】或【解析】，
% 第 11、13、14 题的答案还直接印在题干末尾的括号里。试题文字、答案与解析均照原卷录入，
% 未改内容；卷面的粉色斜水印（“魔都数学加油站”字样）属水印，未录入。
%
% 卷首标题：原卷页首印作“2026-2027 年大同中学高一上周末作业 4.2”（**卷面标题行即带校名**，
% 与高一16／高一18 等拍照卷须由本稿补校名的情形不同），本稿照卷面，年份区间按全稿惯例
% 排 en dash，前缀照原卷保留“年”。
%
% 篇幅与题号：本篇**只有 4 页**，卷面自“一、填空题”的**第 3 题**起编，终于“三、解答题”的
% **第 19 题**，共 17 题（填空 3—10、选择 11—14、解答 15—19）。第 1、2 题未在该文中发布——
% 这是该公众号连载讲评版的通例（高一02—高一09 与 高一36、高一38、高一39 等均自第 3 题起编，
% README“说明”已记），**不是截断**，故本稿题号亦自 3 起。它与 高一39（大同中学 周末作业4.1）
% **不构成同一份试卷的两半**：4.1 是另一次作业（11 页，题号**也自第 3 题**起、一直编到**第 24 题**，
% 以集合／常用逻辑用语为主、兼及不等式），两份各自独立起编（都从第 3 题开始），只是同为“大同中学高一上
% 周末作业 4”系列。
%
% 与原卷的排版差异（均已在 README“说明”中交代）：
%   ① 原文自**第 3 题**起（第 1、2 题未在该文中发布），故本稿题号亦自 3 起；
%   ② 原卷本就印有“一、填空题”“二、选择题”“三、解答题”三节标题，本稿照原卷分节，
%      只把标题改排成全稿统一的“大节标题”样式；
%   ③ 原卷为讲评版：第 4、5、6、7、8、17、18 题只印【答案】行，第 3、9、10、12、15、16、19 题
%      印【解析】（无单独的【答案】行），第 11、13、14 题的答案直接印在题干末尾的括号内；
%      本稿统一改为试卷版隐去、答案版以 \ans{}／\ifanswer 块给出，两版彻底分离（照 高一36 口径：
%      只印【解析】的题不另给 \ans{}，只印【答案】的题仅给 \ans{}）；
%   ④ 数集记号原卷作斜体的 $R$（第 5、10、13 题），本稿按全稿惯例统一写作 $\R$；
%   ⑤ 补集符号原卷作上横线（第 5 题的 $\overline{M}$），本稿照录，未改为 $\complement_\R$；
%   ⑥ 包含符号原卷两套并用：第 4 题作**无下划线**的 $\subset$（真包含），第 18 题作**带下划线**的
%      $\sub$（即 $\subseteq$），本稿照录、不归一；
%   ⑦ 原卷填空的横线约两个字宽，本稿按全稿惯例统一排 \blankline[4.4em]；
%   ⑧ 原卷未印副标题，本稿按**本卷实际内容**补出副标题“不等式”——全卷 17 题里约 15 题为
%      不等式（解一元二次／无理／分式不等式、恒成立与根的分布、含参不等式、应用题），
%      集合仅第 4、5 两题（第 18 题只是用集合关系表述解集），常用逻辑用语 0 题；
%   ⑨ 本卷**无插图**（全卷检索确认无“如图”）；
%   ⑩ 标点与单位：第 6 题的“$a,b$”本稿按全稿惯例排顿号“$a$、$b$”，第 10 题的三例
%      “$[\pi]=3,[-1.2]=-2,[\tfrac12]=0$”之间排顿号，第 5 题原卷“全集 $U=R$ 集合 $M$”之间
%      漏排标点、本稿补逗号；车速单位原卷印作斜体 $km/h$（第 19 题），本稿照录；半角括号、逗号、
%      问号一律排全角（第 17、19 题原卷的两个问号本就作全角，本稿照录）。
%
% 原卷存疑两处（照抄原卷、未改，见源码中该处上方注释）：
%   ① 第 13 题的 D 选项与 A 选项文字完全重复（均作“$a<0$，且 $b^2-4ac>0$”），原卷如此。
%   ② 第 10 题题干末作“…的取值范围____”，**漏排“是”**（本卷其余填空均作“…是____”）——
%      本稿照录、未补字，见源码该题上方 `%` 注。

\documentclass[a4paper,11pt]{article}
\usepackage{common}

\begin{document}

\examtitlest[3]{2026--2027 年大同中学高一上周末作业 4.2}{不等式}

\bigsec{一、填空题}

\begin{enumerate}
\setcounter{enumi}{2}

\item 在平面直角坐标系中，若点 $(2,b)$ 到原点的距离不小于 $5$，则 $b$ 的取值范围是 \blankline[4.4em].

\ifanswer
\solbegin
\step{因为点 $(2,b)$ 到原点的距离不小于 $5$，所以 $\sqrt{(2-0)^2+(b-0)^2}\geqslant5$，}
\step{所以 $4+b^2\geqslant25$，得 $b\geqslant\sqrt{21}$ 或 $b\leqslant-\sqrt{21}$，}
\step{则 $b$ 的取值范围是 $(-\infty,-\sqrt{21}]\cup[\sqrt{21},+\infty)$.}
\solend

\fi

\item 若集合 $A=\{x\mid|x|>3\}$，集合 $B=\{x\mid x\geqslant a\}$，且 $B\subset A$，
则实数 $a$ 的取值范围是 \blankline[4.4em].

\ans{$(3,+\infty)$}

% 注：原卷“全集 $U=R$”与“集合 $M$”之间**漏排标点**（只有空格），本稿补逗号——
%     照 高一b 第 13 题“原卷漏排、本稿补上”的先例（高一10、高一24 同例），见文件头⑩。
\item 若全集 $U=\R$，集合 $M=\{x\mid y=\sqrt{9-x^2}\}$，$S=\{x\mid x^2-3x+2\geqslant0\}$，
则 $\overline{M}\cup S=$ \blankline[4.4em].

\ans{$(-\infty,1]\cup[2,+\infty)$}

\item 设 $a$、$b$ 为实数，若不等式 $x^2-ax-b<0$ 的解集为 $\{x\mid2<x<8\}$，
则 $bx^2-ax-1>0$ 的解集为 \blankline[4.4em].

\ans{$\left(-\dfrac12,-\dfrac18\right)$}

\item 不等式 $4x-13+\sqrt{x+5}\geqslant8x-1+\sqrt{x+5}$ 的解集是 \blankline[4.4em].

\ans{$[-5,-3]$}

\item 若关于 $x$ 的不等式 $ax^2-2ax+3\leqslant0$ 的解集为 $\varnothing$，
则实数 $a$ 的取值范围是 \blankline[4.4em].

\ans{$[0,3)$}

\item 若不等式 $x^2-ax-2<0$ 在 $x\in(1,2)$ 内恒成立，则实数 $a$ 的取值范围是 \blankline[4.4em].

\ifanswer
\solbegin
\step{由不等式 $x^2-ax-2<0$ 在 $x\in(1,2)$ 内恒成立，}
\step{得 $ax>x^2-2$，即 $a>x-\dfrac2x$ 在 $x\in(1,2)$ 内恒成立，}
\step{令 $t(x)=x-\dfrac2x$，$x\in(1,2)$，该函数为严格增函数，则 $t(x)<t(2)=1$，得 $a\geqslant1$，}
\step{所以 $a$ 的取值范围是 $[1,+\infty)$.}
\solend

\fi

% 注：原卷此处作“的取值范围\_\_\_\_”，**漏排“是”**（本卷其余填空均作“…是\_\_\_\_”），
%     本稿照录、未补字，见文件头“原卷存疑”②。
\item 设 $x\in\R$，$[x]$ 表示不大于 $x$ 的最大整数，如 $[\pi]=3$、$[-1.2]=-2$、$\left[\dfrac12\right]=0$，
则使 $[|x-1|]=3$ 成立的 $x$ 的取值范围 \blankline[4.4em].

\ifanswer
\solbegin
\step{由题意 $[|x-1|]=3$，则 $3\leqslant|x-1|<4$，所以 $3\leqslant x-1<4$ 或 $-4\leqslant x-1<-3$，}
\step{解得 $4\leqslant x<5$ 或 $-3<x\leqslant-2$，}
\step{所以使 $[|x-1|]=3$ 成立的 $x$ 的取值范围是 $(-3,-2]\cup[4,5)$.}
\solend

\fi

\end{enumerate}

\bigsec{二、选择题}

\begin{enumerate}
\setcounter{enumi}{10}

\item 若关于 $x$ 的方程 $k(x+2)-3(k-1)=x+1$ 有负数解，则实数 $k$ 的取值范围是\mbox{（\quad）}.

\keepwithprev
A. $(1,2)$\qquad B. $(-\infty,1)\cup(2,+\infty)$\qquad C. $(-\infty,2)$\qquad D. $(-\infty,1)$

\ans{$A$}

\item 与 $x^2>4$ 同解的不等式是\mbox{（\quad）}.

\keepwithprev
A. $x^2+\dfrac{1}{x-4}>4+\dfrac{1}{x-4}$\qquad B. $x^2+\dfrac1x>4+\dfrac1x$

C. $x^2+\sqrt{x-2}>4+\sqrt{x-2}$\qquad D. $\dfrac{1}{x^2-4}<0$

\ifanswer
\solbegin
\step{不等式 $x^2+\dfrac{1}{x-4}>4+\dfrac{1}{x-4}$ 等价于 $\begin{cases}x^2>4\\ x\neq4\end{cases}$，
与 $x^2>4$ 不同解，故选项 $A$ 错误；}
\step{不等式 $x^2+\dfrac1x>4+\dfrac1x$ 等价于 $\begin{cases}x^2>4\\ x\neq0\end{cases}$，
即 $x^2>4$，所以与 $x^2>4$ 同解，}
\step{故选项 $B$ 正确；}
\step{不等式 $x^2+\sqrt{x-2}>4+\sqrt{x-2}$ 等价于 $\begin{cases}x^2>4\\ x-2\geqslant0\end{cases}$，
与 $x^2>4$ 不同解，}
\step{故选项 $C$ 错误；}
\step{不等式 $\dfrac{1}{x^2-4}<0$ 等价于 $x^2<4$，与 $x^2>4$ 不同解，故选项 $D$ 错误.}
\step{故选 $B$.}
\solend

\fi

\item 若关于 $x$ 的二次不等式 $ax^2+bx+c\geqslant0(a\neq0)$ 的解集是 $\R$，那么\mbox{（\quad）}.

\keepwithprev
A. $a<0$，且 $b^2-4ac>0$\qquad B. $a<0$，且 $b^2-4ac\leqslant0$

% 注：原卷本题 D 选项与 A 选项文字完全重复（均作“$a<0$，且 $b^2-4ac>0$”），
%     照录未改，见文件头“原卷存疑”①。
C. $a>0$，且 $b^2-4ac\leqslant0$\qquad D. $a<0$，且 $b^2-4ac>0$

\ans{$C$}

\item 若不等式 $x^2+mx+1>2x+m$ 对满足 $|m|<2$ 的所有实数 $m$ 恒成立，
则实数 $x$ 的取值范围是\mbox{（\quad）}.

\keepwithprev
A. $-2<x<2$\qquad B. $x\geqslant3$\qquad C. $x\leqslant1$\qquad D. $x\leqslant-1$ 或 $x\geqslant3$

\ans{$D$}

\end{enumerate}

% “三、解答题”原本落在试卷版第 1 页页末、第 15 题翻到次页（切页审计 B 类）。
% 节标题自带的守卫只有 3\baselineskip，够放标题不够放标题＋首题，故在标题前再量一次页高。
\needvspace{4cm}
\bigsec{三、解答题}

\begin{enumerate}
\setcounter{enumi}{14}

\item 设关于 $x$ 的二次方程 $(m-1)x^2+(2m-4)x+m=0$ 有两个正根，求实数 $m$ 的取值范围.

\ifanswer
\solbegin
\step{$\begin{cases}\Delta=(2m-4)^2-4m(m-1)\geqslant0\\ -\dfrac{2m-4}{m-1}>0\\ \dfrac{m}{m-1}>0\end{cases}$，
解得 $1<m\leqslant\dfrac43$.}
\solend

\fi

\ansspace[6]

\item 已知关于 $x$ 的不等式 $2x^2-9x+a<0$ 的解集非空，对于其解集内的每一个 $x$ 的值，
至少能使不等式 $x^2-4x+3<0$ 或 $x^2-6x+8<0$ 中的一个成立，求实数 $a$ 的取值范围.

\ifanswer
\solbegin
\step{不等式 $x^2-4x+3<0$ 的解集为 $(1,3)$；$x^2-6x+8<0$ 的解集为 $(2,4)$.}
\step{关于 $x$ 的不等式 $2x^2-9x+a<0$ 的解集非空，所以 $\Delta=81-8a>0$，解得 $a<\dfrac{81}{8}$.}
\step{所以 $\dfrac{9-\sqrt{81-8a}}{4}<x<\dfrac{9+\sqrt{81-8a}}{4}$.}
\step{由于关于 $x$ 的不等式 $2x^2-9x+a<0$（解集非空）的每一个 $x$ 的值}
\step{至少满足不等式 $x^2-4x+3<0$ 和 $x^2-6x+8<0$ 中的一个，}
\step{所以 $\begin{cases}\dfrac{9-\sqrt{81-8a}}{4}\geqslant1\\ \dfrac{9+\sqrt{81-8a}}{4}\leqslant4\\ a<\dfrac{81}{8}\end{cases}$，
解得 $7\leqslant a<\dfrac{81}{8}$，所以实数 $a$ 的取值范围是 $\left[7,\dfrac{81}{8}\right)$.}
\solend

\fi

\ansspace[8]

\item 是否存在实数 $k$，使 $x=3$ 在不等式 $\dfrac{x+2}{k}>1+\dfrac{x^2-3}{k^2}$ 的解集中？
若存在，求 $k$ 的取值范围；若不存在，说明理由.

\ans{$(2,3)$}

\ansspace[6]

\item 关于 $x$ 的不等式 $x^2-(a+1)x+a<0$ 的解集为 $A$，$x^2-5x+4<0$ 的解集为 $B$，
若 $A\sub B$，求实数 $a$ 的取值范围.

\ans{$[1,4]$}

\ansspace[6]

\item 汽车在行驶中，由于惯性作用，刹车后还要继续向前滑行一段距离才能停住，我们称这段
距离为“刹车距离”. 刹车距离是分析事故的一个重要因素. 在一个限速 $40km/h$ 以内的弯道上，
甲、乙两辆汽车相向而行，发现情况不对，同时刹车，但还是相撞了，事后现场测得甲车的刹车
距离略超过 $12m$，乙车的刹车距离略超过 $10m$，又知甲、乙两种车型的刹车距离 $s(m)$ 与
车速 $x(km/h)$ 之间有如下关系：$s_{\text{甲}}=0.1x+0.01x^2$，$s_{\text{乙}}=0.05x+0.005x^2$. 问：
两车相撞的主要责任是谁？

\ifanswer
\solbegin
\step{由题意列出不等式组 $\begin{cases}0.1x+0.01x^2>12\\ 0.05x+0.005x^2>10\end{cases}$，
得 $\begin{cases}x<-40\text{ 或 }x>30\\ x<-50\text{ 或 }x>40\end{cases}$，}
\step{由于 $x>0$，从而得 $x_{\text{甲}}>30km/h$，$x_{\text{乙}}>40km/h$.}
\step{经比较得乙车超过限速，应负主要责任.}
\solend

\fi

\ansspace*[8]

\end{enumerate}

\end{document}
"
% 紧凑版：xelatex "\def\iscompact{1}% 高一40_大同中学_周末作业4.2.tex
%   —— 高一上数学“2026-2027 年大同中学高一上周末作业 4.2”（试卷版/答案版）
% 试卷版：xelatex 高一40_大同中学_周末作业4.2.tex
% 答案版：xelatex "\def\isanswer{1}% 高一40_大同中学_周末作业4.2.tex
%   —— 高一上数学“2026-2027 年大同中学高一上周末作业 4.2”（试卷版/答案版）
% 试卷版：xelatex 高一40_大同中学_周末作业4.2.tex
% 答案版：xelatex "\def\isanswer{1}\input{高一40_大同中学_周末作业4.2.tex}"
% 紧凑版：xelatex "\def\iscompact{1}\input{高一40_大同中学_周末作业4.2.tex}"
%
% 原始材料是微信公众号图文（公众号“魔都数学加油站”连载“27学年高一 40”，署名“破晓”，
% 2026-09-26 21:29 发布，原文标题“27学年高一40——2026-2027年上海市大同中学高一上周末作业4.2”，
% 原文链接 https://mp.weixin.qq.com/s/RrBSqDF_nJVbbx-bmmdtkQ ），正文为 4 张 PNG 页（4762×6735）。原文本身就是一份**讲评版**——题下直接印【答案】或【解析】，
% 第 11、13、14 题的答案还直接印在题干末尾的括号里。试题文字、答案与解析均照原卷录入，
% 未改内容；卷面的粉色斜水印（“魔都数学加油站”字样）属水印，未录入。
%
% 卷首标题：原卷页首印作“2026-2027 年大同中学高一上周末作业 4.2”（**卷面标题行即带校名**，
% 与高一16／高一18 等拍照卷须由本稿补校名的情形不同），本稿照卷面，年份区间按全稿惯例
% 排 en dash，前缀照原卷保留“年”。
%
% 篇幅与题号：本篇**只有 4 页**，卷面自“一、填空题”的**第 3 题**起编，终于“三、解答题”的
% **第 19 题**，共 17 题（填空 3—10、选择 11—14、解答 15—19）。第 1、2 题未在该文中发布——
% 这是该公众号连载讲评版的通例（高一02—高一09 与 高一36、高一38、高一39 等均自第 3 题起编，
% README“说明”已记），**不是截断**，故本稿题号亦自 3 起。它与 高一39（大同中学 周末作业4.1）
% **不构成同一份试卷的两半**：4.1 是另一次作业（11 页，题号**也自第 3 题**起、一直编到**第 24 题**，
% 以集合／常用逻辑用语为主、兼及不等式），两份各自独立起编（都从第 3 题开始），只是同为“大同中学高一上
% 周末作业 4”系列。
%
% 与原卷的排版差异（均已在 README“说明”中交代）：
%   ① 原文自**第 3 题**起（第 1、2 题未在该文中发布），故本稿题号亦自 3 起；
%   ② 原卷本就印有“一、填空题”“二、选择题”“三、解答题”三节标题，本稿照原卷分节，
%      只把标题改排成全稿统一的“大节标题”样式；
%   ③ 原卷为讲评版：第 4、5、6、7、8、17、18 题只印【答案】行，第 3、9、10、12、15、16、19 题
%      印【解析】（无单独的【答案】行），第 11、13、14 题的答案直接印在题干末尾的括号内；
%      本稿统一改为试卷版隐去、答案版以 \ans{}／\ifanswer 块给出，两版彻底分离（照 高一36 口径：
%      只印【解析】的题不另给 \ans{}，只印【答案】的题仅给 \ans{}）；
%   ④ 数集记号原卷作斜体的 $R$（第 5、10、13 题），本稿按全稿惯例统一写作 $\R$；
%   ⑤ 补集符号原卷作上横线（第 5 题的 $\overline{M}$），本稿照录，未改为 $\complement_\R$；
%   ⑥ 包含符号原卷两套并用：第 4 题作**无下划线**的 $\subset$（真包含），第 18 题作**带下划线**的
%      $\sub$（即 $\subseteq$），本稿照录、不归一；
%   ⑦ 原卷填空的横线约两个字宽，本稿按全稿惯例统一排 \blankline[4.4em]；
%   ⑧ 原卷未印副标题，本稿按**本卷实际内容**补出副标题“不等式”——全卷 17 题里约 15 题为
%      不等式（解一元二次／无理／分式不等式、恒成立与根的分布、含参不等式、应用题），
%      集合仅第 4、5 两题（第 18 题只是用集合关系表述解集），常用逻辑用语 0 题；
%   ⑨ 本卷**无插图**（全卷检索确认无“如图”）；
%   ⑩ 标点与单位：第 6 题的“$a,b$”本稿按全稿惯例排顿号“$a$、$b$”，第 10 题的三例
%      “$[\pi]=3,[-1.2]=-2,[\tfrac12]=0$”之间排顿号，第 5 题原卷“全集 $U=R$ 集合 $M$”之间
%      漏排标点、本稿补逗号；车速单位原卷印作斜体 $km/h$（第 19 题），本稿照录；半角括号、逗号、
%      问号一律排全角（第 17、19 题原卷的两个问号本就作全角，本稿照录）。
%
% 原卷存疑两处（照抄原卷、未改，见源码中该处上方注释）：
%   ① 第 13 题的 D 选项与 A 选项文字完全重复（均作“$a<0$，且 $b^2-4ac>0$”），原卷如此。
%   ② 第 10 题题干末作“…的取值范围____”，**漏排“是”**（本卷其余填空均作“…是____”）——
%      本稿照录、未补字，见源码该题上方 `%` 注。

\documentclass[a4paper,11pt]{article}
\usepackage{common}

\begin{document}

\examtitlest[3]{2026--2027 年大同中学高一上周末作业 4.2}{不等式}

\bigsec{一、填空题}

\begin{enumerate}
\setcounter{enumi}{2}

\item 在平面直角坐标系中，若点 $(2,b)$ 到原点的距离不小于 $5$，则 $b$ 的取值范围是 \blankline[4.4em].

\ifanswer
\solbegin
\step{因为点 $(2,b)$ 到原点的距离不小于 $5$，所以 $\sqrt{(2-0)^2+(b-0)^2}\geqslant5$，}
\step{所以 $4+b^2\geqslant25$，得 $b\geqslant\sqrt{21}$ 或 $b\leqslant-\sqrt{21}$，}
\step{则 $b$ 的取值范围是 $(-\infty,-\sqrt{21}]\cup[\sqrt{21},+\infty)$.}
\solend

\fi

\item 若集合 $A=\{x\mid|x|>3\}$，集合 $B=\{x\mid x\geqslant a\}$，且 $B\subset A$，
则实数 $a$ 的取值范围是 \blankline[4.4em].

\ans{$(3,+\infty)$}

% 注：原卷“全集 $U=R$”与“集合 $M$”之间**漏排标点**（只有空格），本稿补逗号——
%     照 高一b 第 13 题“原卷漏排、本稿补上”的先例（高一10、高一24 同例），见文件头⑩。
\item 若全集 $U=\R$，集合 $M=\{x\mid y=\sqrt{9-x^2}\}$，$S=\{x\mid x^2-3x+2\geqslant0\}$，
则 $\overline{M}\cup S=$ \blankline[4.4em].

\ans{$(-\infty,1]\cup[2,+\infty)$}

\item 设 $a$、$b$ 为实数，若不等式 $x^2-ax-b<0$ 的解集为 $\{x\mid2<x<8\}$，
则 $bx^2-ax-1>0$ 的解集为 \blankline[4.4em].

\ans{$\left(-\dfrac12,-\dfrac18\right)$}

\item 不等式 $4x-13+\sqrt{x+5}\geqslant8x-1+\sqrt{x+5}$ 的解集是 \blankline[4.4em].

\ans{$[-5,-3]$}

\item 若关于 $x$ 的不等式 $ax^2-2ax+3\leqslant0$ 的解集为 $\varnothing$，
则实数 $a$ 的取值范围是 \blankline[4.4em].

\ans{$[0,3)$}

\item 若不等式 $x^2-ax-2<0$ 在 $x\in(1,2)$ 内恒成立，则实数 $a$ 的取值范围是 \blankline[4.4em].

\ifanswer
\solbegin
\step{由不等式 $x^2-ax-2<0$ 在 $x\in(1,2)$ 内恒成立，}
\step{得 $ax>x^2-2$，即 $a>x-\dfrac2x$ 在 $x\in(1,2)$ 内恒成立，}
\step{令 $t(x)=x-\dfrac2x$，$x\in(1,2)$，该函数为严格增函数，则 $t(x)<t(2)=1$，得 $a\geqslant1$，}
\step{所以 $a$ 的取值范围是 $[1,+\infty)$.}
\solend

\fi

% 注：原卷此处作“的取值范围\_\_\_\_”，**漏排“是”**（本卷其余填空均作“…是\_\_\_\_”），
%     本稿照录、未补字，见文件头“原卷存疑”②。
\item 设 $x\in\R$，$[x]$ 表示不大于 $x$ 的最大整数，如 $[\pi]=3$、$[-1.2]=-2$、$\left[\dfrac12\right]=0$，
则使 $[|x-1|]=3$ 成立的 $x$ 的取值范围 \blankline[4.4em].

\ifanswer
\solbegin
\step{由题意 $[|x-1|]=3$，则 $3\leqslant|x-1|<4$，所以 $3\leqslant x-1<4$ 或 $-4\leqslant x-1<-3$，}
\step{解得 $4\leqslant x<5$ 或 $-3<x\leqslant-2$，}
\step{所以使 $[|x-1|]=3$ 成立的 $x$ 的取值范围是 $(-3,-2]\cup[4,5)$.}
\solend

\fi

\end{enumerate}

\bigsec{二、选择题}

\begin{enumerate}
\setcounter{enumi}{10}

\item 若关于 $x$ 的方程 $k(x+2)-3(k-1)=x+1$ 有负数解，则实数 $k$ 的取值范围是\mbox{（\quad）}.

\keepwithprev
A. $(1,2)$\qquad B. $(-\infty,1)\cup(2,+\infty)$\qquad C. $(-\infty,2)$\qquad D. $(-\infty,1)$

\ans{$A$}

\item 与 $x^2>4$ 同解的不等式是\mbox{（\quad）}.

\keepwithprev
A. $x^2+\dfrac{1}{x-4}>4+\dfrac{1}{x-4}$\qquad B. $x^2+\dfrac1x>4+\dfrac1x$

C. $x^2+\sqrt{x-2}>4+\sqrt{x-2}$\qquad D. $\dfrac{1}{x^2-4}<0$

\ifanswer
\solbegin
\step{不等式 $x^2+\dfrac{1}{x-4}>4+\dfrac{1}{x-4}$ 等价于 $\begin{cases}x^2>4\\ x\neq4\end{cases}$，
与 $x^2>4$ 不同解，故选项 $A$ 错误；}
\step{不等式 $x^2+\dfrac1x>4+\dfrac1x$ 等价于 $\begin{cases}x^2>4\\ x\neq0\end{cases}$，
即 $x^2>4$，所以与 $x^2>4$ 同解，}
\step{故选项 $B$ 正确；}
\step{不等式 $x^2+\sqrt{x-2}>4+\sqrt{x-2}$ 等价于 $\begin{cases}x^2>4\\ x-2\geqslant0\end{cases}$，
与 $x^2>4$ 不同解，}
\step{故选项 $C$ 错误；}
\step{不等式 $\dfrac{1}{x^2-4}<0$ 等价于 $x^2<4$，与 $x^2>4$ 不同解，故选项 $D$ 错误.}
\step{故选 $B$.}
\solend

\fi

\item 若关于 $x$ 的二次不等式 $ax^2+bx+c\geqslant0(a\neq0)$ 的解集是 $\R$，那么\mbox{（\quad）}.

\keepwithprev
A. $a<0$，且 $b^2-4ac>0$\qquad B. $a<0$，且 $b^2-4ac\leqslant0$

% 注：原卷本题 D 选项与 A 选项文字完全重复（均作“$a<0$，且 $b^2-4ac>0$”），
%     照录未改，见文件头“原卷存疑”①。
C. $a>0$，且 $b^2-4ac\leqslant0$\qquad D. $a<0$，且 $b^2-4ac>0$

\ans{$C$}

\item 若不等式 $x^2+mx+1>2x+m$ 对满足 $|m|<2$ 的所有实数 $m$ 恒成立，
则实数 $x$ 的取值范围是\mbox{（\quad）}.

\keepwithprev
A. $-2<x<2$\qquad B. $x\geqslant3$\qquad C. $x\leqslant1$\qquad D. $x\leqslant-1$ 或 $x\geqslant3$

\ans{$D$}

\end{enumerate}

% “三、解答题”原本落在试卷版第 1 页页末、第 15 题翻到次页（切页审计 B 类）。
% 节标题自带的守卫只有 3\baselineskip，够放标题不够放标题＋首题，故在标题前再量一次页高。
\needvspace{4cm}
\bigsec{三、解答题}

\begin{enumerate}
\setcounter{enumi}{14}

\item 设关于 $x$ 的二次方程 $(m-1)x^2+(2m-4)x+m=0$ 有两个正根，求实数 $m$ 的取值范围.

\ifanswer
\solbegin
\step{$\begin{cases}\Delta=(2m-4)^2-4m(m-1)\geqslant0\\ -\dfrac{2m-4}{m-1}>0\\ \dfrac{m}{m-1}>0\end{cases}$，
解得 $1<m\leqslant\dfrac43$.}
\solend

\fi

\ansspace[6]

\item 已知关于 $x$ 的不等式 $2x^2-9x+a<0$ 的解集非空，对于其解集内的每一个 $x$ 的值，
至少能使不等式 $x^2-4x+3<0$ 或 $x^2-6x+8<0$ 中的一个成立，求实数 $a$ 的取值范围.

\ifanswer
\solbegin
\step{不等式 $x^2-4x+3<0$ 的解集为 $(1,3)$；$x^2-6x+8<0$ 的解集为 $(2,4)$.}
\step{关于 $x$ 的不等式 $2x^2-9x+a<0$ 的解集非空，所以 $\Delta=81-8a>0$，解得 $a<\dfrac{81}{8}$.}
\step{所以 $\dfrac{9-\sqrt{81-8a}}{4}<x<\dfrac{9+\sqrt{81-8a}}{4}$.}
\step{由于关于 $x$ 的不等式 $2x^2-9x+a<0$（解集非空）的每一个 $x$ 的值}
\step{至少满足不等式 $x^2-4x+3<0$ 和 $x^2-6x+8<0$ 中的一个，}
\step{所以 $\begin{cases}\dfrac{9-\sqrt{81-8a}}{4}\geqslant1\\ \dfrac{9+\sqrt{81-8a}}{4}\leqslant4\\ a<\dfrac{81}{8}\end{cases}$，
解得 $7\leqslant a<\dfrac{81}{8}$，所以实数 $a$ 的取值范围是 $\left[7,\dfrac{81}{8}\right)$.}
\solend

\fi

\ansspace[8]

\item 是否存在实数 $k$，使 $x=3$ 在不等式 $\dfrac{x+2}{k}>1+\dfrac{x^2-3}{k^2}$ 的解集中？
若存在，求 $k$ 的取值范围；若不存在，说明理由.

\ans{$(2,3)$}

\ansspace[6]

\item 关于 $x$ 的不等式 $x^2-(a+1)x+a<0$ 的解集为 $A$，$x^2-5x+4<0$ 的解集为 $B$，
若 $A\sub B$，求实数 $a$ 的取值范围.

\ans{$[1,4]$}

\ansspace[6]

\item 汽车在行驶中，由于惯性作用，刹车后还要继续向前滑行一段距离才能停住，我们称这段
距离为“刹车距离”. 刹车距离是分析事故的一个重要因素. 在一个限速 $40km/h$ 以内的弯道上，
甲、乙两辆汽车相向而行，发现情况不对，同时刹车，但还是相撞了，事后现场测得甲车的刹车
距离略超过 $12m$，乙车的刹车距离略超过 $10m$，又知甲、乙两种车型的刹车距离 $s(m)$ 与
车速 $x(km/h)$ 之间有如下关系：$s_{\text{甲}}=0.1x+0.01x^2$，$s_{\text{乙}}=0.05x+0.005x^2$. 问：
两车相撞的主要责任是谁？

\ifanswer
\solbegin
\step{由题意列出不等式组 $\begin{cases}0.1x+0.01x^2>12\\ 0.05x+0.005x^2>10\end{cases}$，
得 $\begin{cases}x<-40\text{ 或 }x>30\\ x<-50\text{ 或 }x>40\end{cases}$，}
\step{由于 $x>0$，从而得 $x_{\text{甲}}>30km/h$，$x_{\text{乙}}>40km/h$.}
\step{经比较得乙车超过限速，应负主要责任.}
\solend

\fi

\ansspace*[8]

\end{enumerate}

\end{document}
"
% 紧凑版：xelatex "\def\iscompact{1}% 高一40_大同中学_周末作业4.2.tex
%   —— 高一上数学“2026-2027 年大同中学高一上周末作业 4.2”（试卷版/答案版）
% 试卷版：xelatex 高一40_大同中学_周末作业4.2.tex
% 答案版：xelatex "\def\isanswer{1}\input{高一40_大同中学_周末作业4.2.tex}"
% 紧凑版：xelatex "\def\iscompact{1}\input{高一40_大同中学_周末作业4.2.tex}"
%
% 原始材料是微信公众号图文（公众号“魔都数学加油站”连载“27学年高一 40”，署名“破晓”，
% 2026-09-26 21:29 发布，原文标题“27学年高一40——2026-2027年上海市大同中学高一上周末作业4.2”，
% 原文链接 https://mp.weixin.qq.com/s/RrBSqDF_nJVbbx-bmmdtkQ ），正文为 4 张 PNG 页（4762×6735）。原文本身就是一份**讲评版**——题下直接印【答案】或【解析】，
% 第 11、13、14 题的答案还直接印在题干末尾的括号里。试题文字、答案与解析均照原卷录入，
% 未改内容；卷面的粉色斜水印（“魔都数学加油站”字样）属水印，未录入。
%
% 卷首标题：原卷页首印作“2026-2027 年大同中学高一上周末作业 4.2”（**卷面标题行即带校名**，
% 与高一16／高一18 等拍照卷须由本稿补校名的情形不同），本稿照卷面，年份区间按全稿惯例
% 排 en dash，前缀照原卷保留“年”。
%
% 篇幅与题号：本篇**只有 4 页**，卷面自“一、填空题”的**第 3 题**起编，终于“三、解答题”的
% **第 19 题**，共 17 题（填空 3—10、选择 11—14、解答 15—19）。第 1、2 题未在该文中发布——
% 这是该公众号连载讲评版的通例（高一02—高一09 与 高一36、高一38、高一39 等均自第 3 题起编，
% README“说明”已记），**不是截断**，故本稿题号亦自 3 起。它与 高一39（大同中学 周末作业4.1）
% **不构成同一份试卷的两半**：4.1 是另一次作业（11 页，题号**也自第 3 题**起、一直编到**第 24 题**，
% 以集合／常用逻辑用语为主、兼及不等式），两份各自独立起编（都从第 3 题开始），只是同为“大同中学高一上
% 周末作业 4”系列。
%
% 与原卷的排版差异（均已在 README“说明”中交代）：
%   ① 原文自**第 3 题**起（第 1、2 题未在该文中发布），故本稿题号亦自 3 起；
%   ② 原卷本就印有“一、填空题”“二、选择题”“三、解答题”三节标题，本稿照原卷分节，
%      只把标题改排成全稿统一的“大节标题”样式；
%   ③ 原卷为讲评版：第 4、5、6、7、8、17、18 题只印【答案】行，第 3、9、10、12、15、16、19 题
%      印【解析】（无单独的【答案】行），第 11、13、14 题的答案直接印在题干末尾的括号内；
%      本稿统一改为试卷版隐去、答案版以 \ans{}／\ifanswer 块给出，两版彻底分离（照 高一36 口径：
%      只印【解析】的题不另给 \ans{}，只印【答案】的题仅给 \ans{}）；
%   ④ 数集记号原卷作斜体的 $R$（第 5、10、13 题），本稿按全稿惯例统一写作 $\R$；
%   ⑤ 补集符号原卷作上横线（第 5 题的 $\overline{M}$），本稿照录，未改为 $\complement_\R$；
%   ⑥ 包含符号原卷两套并用：第 4 题作**无下划线**的 $\subset$（真包含），第 18 题作**带下划线**的
%      $\sub$（即 $\subseteq$），本稿照录、不归一；
%   ⑦ 原卷填空的横线约两个字宽，本稿按全稿惯例统一排 \blankline[4.4em]；
%   ⑧ 原卷未印副标题，本稿按**本卷实际内容**补出副标题“不等式”——全卷 17 题里约 15 题为
%      不等式（解一元二次／无理／分式不等式、恒成立与根的分布、含参不等式、应用题），
%      集合仅第 4、5 两题（第 18 题只是用集合关系表述解集），常用逻辑用语 0 题；
%   ⑨ 本卷**无插图**（全卷检索确认无“如图”）；
%   ⑩ 标点与单位：第 6 题的“$a,b$”本稿按全稿惯例排顿号“$a$、$b$”，第 10 题的三例
%      “$[\pi]=3,[-1.2]=-2,[\tfrac12]=0$”之间排顿号，第 5 题原卷“全集 $U=R$ 集合 $M$”之间
%      漏排标点、本稿补逗号；车速单位原卷印作斜体 $km/h$（第 19 题），本稿照录；半角括号、逗号、
%      问号一律排全角（第 17、19 题原卷的两个问号本就作全角，本稿照录）。
%
% 原卷存疑两处（照抄原卷、未改，见源码中该处上方注释）：
%   ① 第 13 题的 D 选项与 A 选项文字完全重复（均作“$a<0$，且 $b^2-4ac>0$”），原卷如此。
%   ② 第 10 题题干末作“…的取值范围____”，**漏排“是”**（本卷其余填空均作“…是____”）——
%      本稿照录、未补字，见源码该题上方 `%` 注。

\documentclass[a4paper,11pt]{article}
\usepackage{common}

\begin{document}

\examtitlest[3]{2026--2027 年大同中学高一上周末作业 4.2}{不等式}

\bigsec{一、填空题}

\begin{enumerate}
\setcounter{enumi}{2}

\item 在平面直角坐标系中，若点 $(2,b)$ 到原点的距离不小于 $5$，则 $b$ 的取值范围是 \blankline[4.4em].

\ifanswer
\solbegin
\step{因为点 $(2,b)$ 到原点的距离不小于 $5$，所以 $\sqrt{(2-0)^2+(b-0)^2}\geqslant5$，}
\step{所以 $4+b^2\geqslant25$，得 $b\geqslant\sqrt{21}$ 或 $b\leqslant-\sqrt{21}$，}
\step{则 $b$ 的取值范围是 $(-\infty,-\sqrt{21}]\cup[\sqrt{21},+\infty)$.}
\solend

\fi

\item 若集合 $A=\{x\mid|x|>3\}$，集合 $B=\{x\mid x\geqslant a\}$，且 $B\subset A$，
则实数 $a$ 的取值范围是 \blankline[4.4em].

\ans{$(3,+\infty)$}

% 注：原卷“全集 $U=R$”与“集合 $M$”之间**漏排标点**（只有空格），本稿补逗号——
%     照 高一b 第 13 题“原卷漏排、本稿补上”的先例（高一10、高一24 同例），见文件头⑩。
\item 若全集 $U=\R$，集合 $M=\{x\mid y=\sqrt{9-x^2}\}$，$S=\{x\mid x^2-3x+2\geqslant0\}$，
则 $\overline{M}\cup S=$ \blankline[4.4em].

\ans{$(-\infty,1]\cup[2,+\infty)$}

\item 设 $a$、$b$ 为实数，若不等式 $x^2-ax-b<0$ 的解集为 $\{x\mid2<x<8\}$，
则 $bx^2-ax-1>0$ 的解集为 \blankline[4.4em].

\ans{$\left(-\dfrac12,-\dfrac18\right)$}

\item 不等式 $4x-13+\sqrt{x+5}\geqslant8x-1+\sqrt{x+5}$ 的解集是 \blankline[4.4em].

\ans{$[-5,-3]$}

\item 若关于 $x$ 的不等式 $ax^2-2ax+3\leqslant0$ 的解集为 $\varnothing$，
则实数 $a$ 的取值范围是 \blankline[4.4em].

\ans{$[0,3)$}

\item 若不等式 $x^2-ax-2<0$ 在 $x\in(1,2)$ 内恒成立，则实数 $a$ 的取值范围是 \blankline[4.4em].

\ifanswer
\solbegin
\step{由不等式 $x^2-ax-2<0$ 在 $x\in(1,2)$ 内恒成立，}
\step{得 $ax>x^2-2$，即 $a>x-\dfrac2x$ 在 $x\in(1,2)$ 内恒成立，}
\step{令 $t(x)=x-\dfrac2x$，$x\in(1,2)$，该函数为严格增函数，则 $t(x)<t(2)=1$，得 $a\geqslant1$，}
\step{所以 $a$ 的取值范围是 $[1,+\infty)$.}
\solend

\fi

% 注：原卷此处作“的取值范围\_\_\_\_”，**漏排“是”**（本卷其余填空均作“…是\_\_\_\_”），
%     本稿照录、未补字，见文件头“原卷存疑”②。
\item 设 $x\in\R$，$[x]$ 表示不大于 $x$ 的最大整数，如 $[\pi]=3$、$[-1.2]=-2$、$\left[\dfrac12\right]=0$，
则使 $[|x-1|]=3$ 成立的 $x$ 的取值范围 \blankline[4.4em].

\ifanswer
\solbegin
\step{由题意 $[|x-1|]=3$，则 $3\leqslant|x-1|<4$，所以 $3\leqslant x-1<4$ 或 $-4\leqslant x-1<-3$，}
\step{解得 $4\leqslant x<5$ 或 $-3<x\leqslant-2$，}
\step{所以使 $[|x-1|]=3$ 成立的 $x$ 的取值范围是 $(-3,-2]\cup[4,5)$.}
\solend

\fi

\end{enumerate}

\bigsec{二、选择题}

\begin{enumerate}
\setcounter{enumi}{10}

\item 若关于 $x$ 的方程 $k(x+2)-3(k-1)=x+1$ 有负数解，则实数 $k$ 的取值范围是\mbox{（\quad）}.

\keepwithprev
A. $(1,2)$\qquad B. $(-\infty,1)\cup(2,+\infty)$\qquad C. $(-\infty,2)$\qquad D. $(-\infty,1)$

\ans{$A$}

\item 与 $x^2>4$ 同解的不等式是\mbox{（\quad）}.

\keepwithprev
A. $x^2+\dfrac{1}{x-4}>4+\dfrac{1}{x-4}$\qquad B. $x^2+\dfrac1x>4+\dfrac1x$

C. $x^2+\sqrt{x-2}>4+\sqrt{x-2}$\qquad D. $\dfrac{1}{x^2-4}<0$

\ifanswer
\solbegin
\step{不等式 $x^2+\dfrac{1}{x-4}>4+\dfrac{1}{x-4}$ 等价于 $\begin{cases}x^2>4\\ x\neq4\end{cases}$，
与 $x^2>4$ 不同解，故选项 $A$ 错误；}
\step{不等式 $x^2+\dfrac1x>4+\dfrac1x$ 等价于 $\begin{cases}x^2>4\\ x\neq0\end{cases}$，
即 $x^2>4$，所以与 $x^2>4$ 同解，}
\step{故选项 $B$ 正确；}
\step{不等式 $x^2+\sqrt{x-2}>4+\sqrt{x-2}$ 等价于 $\begin{cases}x^2>4\\ x-2\geqslant0\end{cases}$，
与 $x^2>4$ 不同解，}
\step{故选项 $C$ 错误；}
\step{不等式 $\dfrac{1}{x^2-4}<0$ 等价于 $x^2<4$，与 $x^2>4$ 不同解，故选项 $D$ 错误.}
\step{故选 $B$.}
\solend

\fi

\item 若关于 $x$ 的二次不等式 $ax^2+bx+c\geqslant0(a\neq0)$ 的解集是 $\R$，那么\mbox{（\quad）}.

\keepwithprev
A. $a<0$，且 $b^2-4ac>0$\qquad B. $a<0$，且 $b^2-4ac\leqslant0$

% 注：原卷本题 D 选项与 A 选项文字完全重复（均作“$a<0$，且 $b^2-4ac>0$”），
%     照录未改，见文件头“原卷存疑”①。
C. $a>0$，且 $b^2-4ac\leqslant0$\qquad D. $a<0$，且 $b^2-4ac>0$

\ans{$C$}

\item 若不等式 $x^2+mx+1>2x+m$ 对满足 $|m|<2$ 的所有实数 $m$ 恒成立，
则实数 $x$ 的取值范围是\mbox{（\quad）}.

\keepwithprev
A. $-2<x<2$\qquad B. $x\geqslant3$\qquad C. $x\leqslant1$\qquad D. $x\leqslant-1$ 或 $x\geqslant3$

\ans{$D$}

\end{enumerate}

% “三、解答题”原本落在试卷版第 1 页页末、第 15 题翻到次页（切页审计 B 类）。
% 节标题自带的守卫只有 3\baselineskip，够放标题不够放标题＋首题，故在标题前再量一次页高。
\needvspace{4cm}
\bigsec{三、解答题}

\begin{enumerate}
\setcounter{enumi}{14}

\item 设关于 $x$ 的二次方程 $(m-1)x^2+(2m-4)x+m=0$ 有两个正根，求实数 $m$ 的取值范围.

\ifanswer
\solbegin
\step{$\begin{cases}\Delta=(2m-4)^2-4m(m-1)\geqslant0\\ -\dfrac{2m-4}{m-1}>0\\ \dfrac{m}{m-1}>0\end{cases}$，
解得 $1<m\leqslant\dfrac43$.}
\solend

\fi

\ansspace[6]

\item 已知关于 $x$ 的不等式 $2x^2-9x+a<0$ 的解集非空，对于其解集内的每一个 $x$ 的值，
至少能使不等式 $x^2-4x+3<0$ 或 $x^2-6x+8<0$ 中的一个成立，求实数 $a$ 的取值范围.

\ifanswer
\solbegin
\step{不等式 $x^2-4x+3<0$ 的解集为 $(1,3)$；$x^2-6x+8<0$ 的解集为 $(2,4)$.}
\step{关于 $x$ 的不等式 $2x^2-9x+a<0$ 的解集非空，所以 $\Delta=81-8a>0$，解得 $a<\dfrac{81}{8}$.}
\step{所以 $\dfrac{9-\sqrt{81-8a}}{4}<x<\dfrac{9+\sqrt{81-8a}}{4}$.}
\step{由于关于 $x$ 的不等式 $2x^2-9x+a<0$（解集非空）的每一个 $x$ 的值}
\step{至少满足不等式 $x^2-4x+3<0$ 和 $x^2-6x+8<0$ 中的一个，}
\step{所以 $\begin{cases}\dfrac{9-\sqrt{81-8a}}{4}\geqslant1\\ \dfrac{9+\sqrt{81-8a}}{4}\leqslant4\\ a<\dfrac{81}{8}\end{cases}$，
解得 $7\leqslant a<\dfrac{81}{8}$，所以实数 $a$ 的取值范围是 $\left[7,\dfrac{81}{8}\right)$.}
\solend

\fi

\ansspace[8]

\item 是否存在实数 $k$，使 $x=3$ 在不等式 $\dfrac{x+2}{k}>1+\dfrac{x^2-3}{k^2}$ 的解集中？
若存在，求 $k$ 的取值范围；若不存在，说明理由.

\ans{$(2,3)$}

\ansspace[6]

\item 关于 $x$ 的不等式 $x^2-(a+1)x+a<0$ 的解集为 $A$，$x^2-5x+4<0$ 的解集为 $B$，
若 $A\sub B$，求实数 $a$ 的取值范围.

\ans{$[1,4]$}

\ansspace[6]

\item 汽车在行驶中，由于惯性作用，刹车后还要继续向前滑行一段距离才能停住，我们称这段
距离为“刹车距离”. 刹车距离是分析事故的一个重要因素. 在一个限速 $40km/h$ 以内的弯道上，
甲、乙两辆汽车相向而行，发现情况不对，同时刹车，但还是相撞了，事后现场测得甲车的刹车
距离略超过 $12m$，乙车的刹车距离略超过 $10m$，又知甲、乙两种车型的刹车距离 $s(m)$ 与
车速 $x(km/h)$ 之间有如下关系：$s_{\text{甲}}=0.1x+0.01x^2$，$s_{\text{乙}}=0.05x+0.005x^2$. 问：
两车相撞的主要责任是谁？

\ifanswer
\solbegin
\step{由题意列出不等式组 $\begin{cases}0.1x+0.01x^2>12\\ 0.05x+0.005x^2>10\end{cases}$，
得 $\begin{cases}x<-40\text{ 或 }x>30\\ x<-50\text{ 或 }x>40\end{cases}$，}
\step{由于 $x>0$，从而得 $x_{\text{甲}}>30km/h$，$x_{\text{乙}}>40km/h$.}
\step{经比较得乙车超过限速，应负主要责任.}
\solend

\fi

\ansspace*[8]

\end{enumerate}

\end{document}
"
%
% 原始材料是微信公众号图文（公众号“魔都数学加油站”连载“27学年高一 40”，署名“破晓”，
% 2026-09-26 21:29 发布，原文标题“27学年高一40——2026-2027年上海市大同中学高一上周末作业4.2”，
% 原文链接 https://mp.weixin.qq.com/s/RrBSqDF_nJVbbx-bmmdtkQ ），正文为 4 张 PNG 页（4762×6735）。原文本身就是一份**讲评版**——题下直接印【答案】或【解析】，
% 第 11、13、14 题的答案还直接印在题干末尾的括号里。试题文字、答案与解析均照原卷录入，
% 未改内容；卷面的粉色斜水印（“魔都数学加油站”字样）属水印，未录入。
%
% 卷首标题：原卷页首印作“2026-2027 年大同中学高一上周末作业 4.2”（**卷面标题行即带校名**，
% 与高一16／高一18 等拍照卷须由本稿补校名的情形不同），本稿照卷面，年份区间按全稿惯例
% 排 en dash，前缀照原卷保留“年”。
%
% 篇幅与题号：本篇**只有 4 页**，卷面自“一、填空题”的**第 3 题**起编，终于“三、解答题”的
% **第 19 题**，共 17 题（填空 3—10、选择 11—14、解答 15—19）。第 1、2 题未在该文中发布——
% 这是该公众号连载讲评版的通例（高一02—高一09 与 高一36、高一38、高一39 等均自第 3 题起编，
% README“说明”已记），**不是截断**，故本稿题号亦自 3 起。它与 高一39（大同中学 周末作业4.1）
% **不构成同一份试卷的两半**：4.1 是另一次作业（11 页，题号**也自第 3 题**起、一直编到**第 24 题**，
% 以集合／常用逻辑用语为主、兼及不等式），两份各自独立起编（都从第 3 题开始），只是同为“大同中学高一上
% 周末作业 4”系列。
%
% 与原卷的排版差异（均已在 README“说明”中交代）：
%   ① 原文自**第 3 题**起（第 1、2 题未在该文中发布），故本稿题号亦自 3 起；
%   ② 原卷本就印有“一、填空题”“二、选择题”“三、解答题”三节标题，本稿照原卷分节，
%      只把标题改排成全稿统一的“大节标题”样式；
%   ③ 原卷为讲评版：第 4、5、6、7、8、17、18 题只印【答案】行，第 3、9、10、12、15、16、19 题
%      印【解析】（无单独的【答案】行），第 11、13、14 题的答案直接印在题干末尾的括号内；
%      本稿统一改为试卷版隐去、答案版以 \ans{}／\ifanswer 块给出，两版彻底分离（照 高一36 口径：
%      只印【解析】的题不另给 \ans{}，只印【答案】的题仅给 \ans{}）；
%   ④ 数集记号原卷作斜体的 $R$（第 5、10、13 题），本稿按全稿惯例统一写作 $\R$；
%   ⑤ 补集符号原卷作上横线（第 5 题的 $\overline{M}$），本稿照录，未改为 $\complement_\R$；
%   ⑥ 包含符号原卷两套并用：第 4 题作**无下划线**的 $\subset$（真包含），第 18 题作**带下划线**的
%      $\sub$（即 $\subseteq$），本稿照录、不归一；
%   ⑦ 原卷填空的横线约两个字宽，本稿按全稿惯例统一排 \blankline[4.4em]；
%   ⑧ 原卷未印副标题，本稿按**本卷实际内容**补出副标题“不等式”——全卷 17 题里约 15 题为
%      不等式（解一元二次／无理／分式不等式、恒成立与根的分布、含参不等式、应用题），
%      集合仅第 4、5 两题（第 18 题只是用集合关系表述解集），常用逻辑用语 0 题；
%   ⑨ 本卷**无插图**（全卷检索确认无“如图”）；
%   ⑩ 标点与单位：第 6 题的“$a,b$”本稿按全稿惯例排顿号“$a$、$b$”，第 10 题的三例
%      “$[\pi]=3,[-1.2]=-2,[\tfrac12]=0$”之间排顿号，第 5 题原卷“全集 $U=R$ 集合 $M$”之间
%      漏排标点、本稿补逗号；车速单位原卷印作斜体 $km/h$（第 19 题），本稿照录；半角括号、逗号、
%      问号一律排全角（第 17、19 题原卷的两个问号本就作全角，本稿照录）。
%
% 原卷存疑两处（照抄原卷、未改，见源码中该处上方注释）：
%   ① 第 13 题的 D 选项与 A 选项文字完全重复（均作“$a<0$，且 $b^2-4ac>0$”），原卷如此。
%   ② 第 10 题题干末作“…的取值范围____”，**漏排“是”**（本卷其余填空均作“…是____”）——
%      本稿照录、未补字，见源码该题上方 `%` 注。

\documentclass[a4paper,11pt]{article}
\usepackage{common}

\begin{document}

\examtitlest[3]{2026--2027 年大同中学高一上周末作业 4.2}{不等式}

\bigsec{一、填空题}

\begin{enumerate}
\setcounter{enumi}{2}

\item 在平面直角坐标系中，若点 $(2,b)$ 到原点的距离不小于 $5$，则 $b$ 的取值范围是 \blankline[4.4em].

\ifanswer
\solbegin
\step{因为点 $(2,b)$ 到原点的距离不小于 $5$，所以 $\sqrt{(2-0)^2+(b-0)^2}\geqslant5$，}
\step{所以 $4+b^2\geqslant25$，得 $b\geqslant\sqrt{21}$ 或 $b\leqslant-\sqrt{21}$，}
\step{则 $b$ 的取值范围是 $(-\infty,-\sqrt{21}]\cup[\sqrt{21},+\infty)$.}
\solend

\fi

\item 若集合 $A=\{x\mid|x|>3\}$，集合 $B=\{x\mid x\geqslant a\}$，且 $B\subset A$，
则实数 $a$ 的取值范围是 \blankline[4.4em].

\ans{$(3,+\infty)$}

% 注：原卷“全集 $U=R$”与“集合 $M$”之间**漏排标点**（只有空格），本稿补逗号——
%     照 高一b 第 13 题“原卷漏排、本稿补上”的先例（高一10、高一24 同例），见文件头⑩。
\item 若全集 $U=\R$，集合 $M=\{x\mid y=\sqrt{9-x^2}\}$，$S=\{x\mid x^2-3x+2\geqslant0\}$，
则 $\overline{M}\cup S=$ \blankline[4.4em].

\ans{$(-\infty,1]\cup[2,+\infty)$}

\item 设 $a$、$b$ 为实数，若不等式 $x^2-ax-b<0$ 的解集为 $\{x\mid2<x<8\}$，
则 $bx^2-ax-1>0$ 的解集为 \blankline[4.4em].

\ans{$\left(-\dfrac12,-\dfrac18\right)$}

\item 不等式 $4x-13+\sqrt{x+5}\geqslant8x-1+\sqrt{x+5}$ 的解集是 \blankline[4.4em].

\ans{$[-5,-3]$}

\item 若关于 $x$ 的不等式 $ax^2-2ax+3\leqslant0$ 的解集为 $\varnothing$，
则实数 $a$ 的取值范围是 \blankline[4.4em].

\ans{$[0,3)$}

\item 若不等式 $x^2-ax-2<0$ 在 $x\in(1,2)$ 内恒成立，则实数 $a$ 的取值范围是 \blankline[4.4em].

\ifanswer
\solbegin
\step{由不等式 $x^2-ax-2<0$ 在 $x\in(1,2)$ 内恒成立，}
\step{得 $ax>x^2-2$，即 $a>x-\dfrac2x$ 在 $x\in(1,2)$ 内恒成立，}
\step{令 $t(x)=x-\dfrac2x$，$x\in(1,2)$，该函数为严格增函数，则 $t(x)<t(2)=1$，得 $a\geqslant1$，}
\step{所以 $a$ 的取值范围是 $[1,+\infty)$.}
\solend

\fi

% 注：原卷此处作“的取值范围\_\_\_\_”，**漏排“是”**（本卷其余填空均作“…是\_\_\_\_”），
%     本稿照录、未补字，见文件头“原卷存疑”②。
\item 设 $x\in\R$，$[x]$ 表示不大于 $x$ 的最大整数，如 $[\pi]=3$、$[-1.2]=-2$、$\left[\dfrac12\right]=0$，
则使 $[|x-1|]=3$ 成立的 $x$ 的取值范围 \blankline[4.4em].

\ifanswer
\solbegin
\step{由题意 $[|x-1|]=3$，则 $3\leqslant|x-1|<4$，所以 $3\leqslant x-1<4$ 或 $-4\leqslant x-1<-3$，}
\step{解得 $4\leqslant x<5$ 或 $-3<x\leqslant-2$，}
\step{所以使 $[|x-1|]=3$ 成立的 $x$ 的取值范围是 $(-3,-2]\cup[4,5)$.}
\solend

\fi

\end{enumerate}

\bigsec{二、选择题}

\begin{enumerate}
\setcounter{enumi}{10}

\item 若关于 $x$ 的方程 $k(x+2)-3(k-1)=x+1$ 有负数解，则实数 $k$ 的取值范围是\mbox{（\quad）}.

\keepwithprev
A. $(1,2)$\qquad B. $(-\infty,1)\cup(2,+\infty)$\qquad C. $(-\infty,2)$\qquad D. $(-\infty,1)$

\ans{$A$}

\item 与 $x^2>4$ 同解的不等式是\mbox{（\quad）}.

\keepwithprev
A. $x^2+\dfrac{1}{x-4}>4+\dfrac{1}{x-4}$\qquad B. $x^2+\dfrac1x>4+\dfrac1x$

C. $x^2+\sqrt{x-2}>4+\sqrt{x-2}$\qquad D. $\dfrac{1}{x^2-4}<0$

\ifanswer
\solbegin
\step{不等式 $x^2+\dfrac{1}{x-4}>4+\dfrac{1}{x-4}$ 等价于 $\begin{cases}x^2>4\\ x\neq4\end{cases}$，
与 $x^2>4$ 不同解，故选项 $A$ 错误；}
\step{不等式 $x^2+\dfrac1x>4+\dfrac1x$ 等价于 $\begin{cases}x^2>4\\ x\neq0\end{cases}$，
即 $x^2>4$，所以与 $x^2>4$ 同解，}
\step{故选项 $B$ 正确；}
\step{不等式 $x^2+\sqrt{x-2}>4+\sqrt{x-2}$ 等价于 $\begin{cases}x^2>4\\ x-2\geqslant0\end{cases}$，
与 $x^2>4$ 不同解，}
\step{故选项 $C$ 错误；}
\step{不等式 $\dfrac{1}{x^2-4}<0$ 等价于 $x^2<4$，与 $x^2>4$ 不同解，故选项 $D$ 错误.}
\step{故选 $B$.}
\solend

\fi

\item 若关于 $x$ 的二次不等式 $ax^2+bx+c\geqslant0(a\neq0)$ 的解集是 $\R$，那么\mbox{（\quad）}.

\keepwithprev
A. $a<0$，且 $b^2-4ac>0$\qquad B. $a<0$，且 $b^2-4ac\leqslant0$

% 注：原卷本题 D 选项与 A 选项文字完全重复（均作“$a<0$，且 $b^2-4ac>0$”），
%     照录未改，见文件头“原卷存疑”①。
C. $a>0$，且 $b^2-4ac\leqslant0$\qquad D. $a<0$，且 $b^2-4ac>0$

\ans{$C$}

\item 若不等式 $x^2+mx+1>2x+m$ 对满足 $|m|<2$ 的所有实数 $m$ 恒成立，
则实数 $x$ 的取值范围是\mbox{（\quad）}.

\keepwithprev
A. $-2<x<2$\qquad B. $x\geqslant3$\qquad C. $x\leqslant1$\qquad D. $x\leqslant-1$ 或 $x\geqslant3$

\ans{$D$}

\end{enumerate}

% “三、解答题”原本落在试卷版第 1 页页末、第 15 题翻到次页（切页审计 B 类）。
% 节标题自带的守卫只有 3\baselineskip，够放标题不够放标题＋首题，故在标题前再量一次页高。
\needvspace{4cm}
\bigsec{三、解答题}

\begin{enumerate}
\setcounter{enumi}{14}

\item 设关于 $x$ 的二次方程 $(m-1)x^2+(2m-4)x+m=0$ 有两个正根，求实数 $m$ 的取值范围.

\ifanswer
\solbegin
\step{$\begin{cases}\Delta=(2m-4)^2-4m(m-1)\geqslant0\\ -\dfrac{2m-4}{m-1}>0\\ \dfrac{m}{m-1}>0\end{cases}$，
解得 $1<m\leqslant\dfrac43$.}
\solend

\fi

\ansspace[6]

\item 已知关于 $x$ 的不等式 $2x^2-9x+a<0$ 的解集非空，对于其解集内的每一个 $x$ 的值，
至少能使不等式 $x^2-4x+3<0$ 或 $x^2-6x+8<0$ 中的一个成立，求实数 $a$ 的取值范围.

\ifanswer
\solbegin
\step{不等式 $x^2-4x+3<0$ 的解集为 $(1,3)$；$x^2-6x+8<0$ 的解集为 $(2,4)$.}
\step{关于 $x$ 的不等式 $2x^2-9x+a<0$ 的解集非空，所以 $\Delta=81-8a>0$，解得 $a<\dfrac{81}{8}$.}
\step{所以 $\dfrac{9-\sqrt{81-8a}}{4}<x<\dfrac{9+\sqrt{81-8a}}{4}$.}
\step{由于关于 $x$ 的不等式 $2x^2-9x+a<0$（解集非空）的每一个 $x$ 的值}
\step{至少满足不等式 $x^2-4x+3<0$ 和 $x^2-6x+8<0$ 中的一个，}
\step{所以 $\begin{cases}\dfrac{9-\sqrt{81-8a}}{4}\geqslant1\\ \dfrac{9+\sqrt{81-8a}}{4}\leqslant4\\ a<\dfrac{81}{8}\end{cases}$，
解得 $7\leqslant a<\dfrac{81}{8}$，所以实数 $a$ 的取值范围是 $\left[7,\dfrac{81}{8}\right)$.}
\solend

\fi

\ansspace[8]

\item 是否存在实数 $k$，使 $x=3$ 在不等式 $\dfrac{x+2}{k}>1+\dfrac{x^2-3}{k^2}$ 的解集中？
若存在，求 $k$ 的取值范围；若不存在，说明理由.

\ans{$(2,3)$}

\ansspace[6]

\item 关于 $x$ 的不等式 $x^2-(a+1)x+a<0$ 的解集为 $A$，$x^2-5x+4<0$ 的解集为 $B$，
若 $A\sub B$，求实数 $a$ 的取值范围.

\ans{$[1,4]$}

\ansspace[6]

\item 汽车在行驶中，由于惯性作用，刹车后还要继续向前滑行一段距离才能停住，我们称这段
距离为“刹车距离”. 刹车距离是分析事故的一个重要因素. 在一个限速 $40km/h$ 以内的弯道上，
甲、乙两辆汽车相向而行，发现情况不对，同时刹车，但还是相撞了，事后现场测得甲车的刹车
距离略超过 $12m$，乙车的刹车距离略超过 $10m$，又知甲、乙两种车型的刹车距离 $s(m)$ 与
车速 $x(km/h)$ 之间有如下关系：$s_{\text{甲}}=0.1x+0.01x^2$，$s_{\text{乙}}=0.05x+0.005x^2$. 问：
两车相撞的主要责任是谁？

\ifanswer
\solbegin
\step{由题意列出不等式组 $\begin{cases}0.1x+0.01x^2>12\\ 0.05x+0.005x^2>10\end{cases}$，
得 $\begin{cases}x<-40\text{ 或 }x>30\\ x<-50\text{ 或 }x>40\end{cases}$，}
\step{由于 $x>0$，从而得 $x_{\text{甲}}>30km/h$，$x_{\text{乙}}>40km/h$.}
\step{经比较得乙车超过限速，应负主要责任.}
\solend

\fi

\ansspace*[8]

\end{enumerate}

\end{document}
"
%
% 原始材料是微信公众号图文（公众号“魔都数学加油站”连载“27学年高一 40”，署名“破晓”，
% 2026-09-26 21:29 发布，原文标题“27学年高一40——2026-2027年上海市大同中学高一上周末作业4.2”，
% 原文链接 https://mp.weixin.qq.com/s/RrBSqDF_nJVbbx-bmmdtkQ ），正文为 4 张 PNG 页（4762×6735）。原文本身就是一份**讲评版**——题下直接印【答案】或【解析】，
% 第 11、13、14 题的答案还直接印在题干末尾的括号里。试题文字、答案与解析均照原卷录入，
% 未改内容；卷面的粉色斜水印（“魔都数学加油站”字样）属水印，未录入。
%
% 卷首标题：原卷页首印作“2026-2027 年大同中学高一上周末作业 4.2”（**卷面标题行即带校名**，
% 与高一16／高一18 等拍照卷须由本稿补校名的情形不同），本稿照卷面，年份区间按全稿惯例
% 排 en dash，前缀照原卷保留“年”。
%
% 篇幅与题号：本篇**只有 4 页**，卷面自“一、填空题”的**第 3 题**起编，终于“三、解答题”的
% **第 19 题**，共 17 题（填空 3—10、选择 11—14、解答 15—19）。第 1、2 题未在该文中发布——
% 这是该公众号连载讲评版的通例（高一02—高一09 与 高一36、高一38、高一39 等均自第 3 题起编，
% README“说明”已记），**不是截断**，故本稿题号亦自 3 起。它与 高一39（大同中学 周末作业4.1）
% **不构成同一份试卷的两半**：4.1 是另一次作业（11 页，题号**也自第 3 题**起、一直编到**第 24 题**，
% 以集合／常用逻辑用语为主、兼及不等式），两份各自独立起编（都从第 3 题开始），只是同为“大同中学高一上
% 周末作业 4”系列。
%
% 与原卷的排版差异（均已在 README“说明”中交代）：
%   ① 原文自**第 3 题**起（第 1、2 题未在该文中发布），故本稿题号亦自 3 起；
%   ② 原卷本就印有“一、填空题”“二、选择题”“三、解答题”三节标题，本稿照原卷分节，
%      只把标题改排成全稿统一的“大节标题”样式；
%   ③ 原卷为讲评版：第 4、5、6、7、8、17、18 题只印【答案】行，第 3、9、10、12、15、16、19 题
%      印【解析】（无单独的【答案】行），第 11、13、14 题的答案直接印在题干末尾的括号内；
%      本稿统一改为试卷版隐去、答案版以 \ans{}／\ifanswer 块给出，两版彻底分离（照 高一36 口径：
%      只印【解析】的题不另给 \ans{}，只印【答案】的题仅给 \ans{}）；
%   ④ 数集记号原卷作斜体的 $R$（第 5、10、13 题），本稿按全稿惯例统一写作 $\R$；
%   ⑤ 补集符号原卷作上横线（第 5 题的 $\overline{M}$），本稿照录，未改为 $\complement_\R$；
%   ⑥ 包含符号原卷两套并用：第 4 题作**无下划线**的 $\subset$（真包含），第 18 题作**带下划线**的
%      $\sub$（即 $\subseteq$），本稿照录、不归一；
%   ⑦ 原卷填空的横线约两个字宽，本稿按全稿惯例统一排 \blankline[4.4em]；
%   ⑧ 原卷未印副标题，本稿按**本卷实际内容**补出副标题“不等式”——全卷 17 题里约 15 题为
%      不等式（解一元二次／无理／分式不等式、恒成立与根的分布、含参不等式、应用题），
%      集合仅第 4、5 两题（第 18 题只是用集合关系表述解集），常用逻辑用语 0 题；
%   ⑨ 本卷**无插图**（全卷检索确认无“如图”）；
%   ⑩ 标点与单位：第 6 题的“$a,b$”本稿按全稿惯例排顿号“$a$、$b$”，第 10 题的三例
%      “$[\pi]=3,[-1.2]=-2,[\tfrac12]=0$”之间排顿号，第 5 题原卷“全集 $U=R$ 集合 $M$”之间
%      漏排标点、本稿补逗号；车速单位原卷印作斜体 $km/h$（第 19 题），本稿照录；半角括号、逗号、
%      问号一律排全角（第 17、19 题原卷的两个问号本就作全角，本稿照录）。
%
% 原卷存疑两处（照抄原卷、未改，见源码中该处上方注释）：
%   ① 第 13 题的 D 选项与 A 选项文字完全重复（均作“$a<0$，且 $b^2-4ac>0$”），原卷如此。
%   ② 第 10 题题干末作“…的取值范围____”，**漏排“是”**（本卷其余填空均作“…是____”）——
%      本稿照录、未补字，见源码该题上方 `%` 注。

\documentclass[a4paper,11pt]{article}
\usepackage{common}

\begin{document}

\examtitlest[3]{2026--2027 年大同中学高一上周末作业 4.2}{不等式}

\bigsec{一、填空题}

\begin{enumerate}
\setcounter{enumi}{2}

\item 在平面直角坐标系中，若点 $(2,b)$ 到原点的距离不小于 $5$，则 $b$ 的取值范围是 \blankline[4.4em].

\ifanswer
\solbegin
\step{因为点 $(2,b)$ 到原点的距离不小于 $5$，所以 $\sqrt{(2-0)^2+(b-0)^2}\geqslant5$，}
\step{所以 $4+b^2\geqslant25$，得 $b\geqslant\sqrt{21}$ 或 $b\leqslant-\sqrt{21}$，}
\step{则 $b$ 的取值范围是 $(-\infty,-\sqrt{21}]\cup[\sqrt{21},+\infty)$.}
\solend

\fi

\item 若集合 $A=\{x\mid|x|>3\}$，集合 $B=\{x\mid x\geqslant a\}$，且 $B\subset A$，
则实数 $a$ 的取值范围是 \blankline[4.4em].

\ans{$(3,+\infty)$}

% 注：原卷“全集 $U=R$”与“集合 $M$”之间**漏排标点**（只有空格），本稿补逗号——
%     照 高一b 第 13 题“原卷漏排、本稿补上”的先例（高一10、高一24 同例），见文件头⑩。
\item 若全集 $U=\R$，集合 $M=\{x\mid y=\sqrt{9-x^2}\}$，$S=\{x\mid x^2-3x+2\geqslant0\}$，
则 $\overline{M}\cup S=$ \blankline[4.4em].

\ans{$(-\infty,1]\cup[2,+\infty)$}

\item 设 $a$、$b$ 为实数，若不等式 $x^2-ax-b<0$ 的解集为 $\{x\mid2<x<8\}$，
则 $bx^2-ax-1>0$ 的解集为 \blankline[4.4em].

\ans{$\left(-\dfrac12,-\dfrac18\right)$}

\item 不等式 $4x-13+\sqrt{x+5}\geqslant8x-1+\sqrt{x+5}$ 的解集是 \blankline[4.4em].

\ans{$[-5,-3]$}

\item 若关于 $x$ 的不等式 $ax^2-2ax+3\leqslant0$ 的解集为 $\varnothing$，
则实数 $a$ 的取值范围是 \blankline[4.4em].

\ans{$[0,3)$}

\item 若不等式 $x^2-ax-2<0$ 在 $x\in(1,2)$ 内恒成立，则实数 $a$ 的取值范围是 \blankline[4.4em].

\ifanswer
\solbegin
\step{由不等式 $x^2-ax-2<0$ 在 $x\in(1,2)$ 内恒成立，}
\step{得 $ax>x^2-2$，即 $a>x-\dfrac2x$ 在 $x\in(1,2)$ 内恒成立，}
\step{令 $t(x)=x-\dfrac2x$，$x\in(1,2)$，该函数为严格增函数，则 $t(x)<t(2)=1$，得 $a\geqslant1$，}
\step{所以 $a$ 的取值范围是 $[1,+\infty)$.}
\solend

\fi

% 注：原卷此处作“的取值范围\_\_\_\_”，**漏排“是”**（本卷其余填空均作“…是\_\_\_\_”），
%     本稿照录、未补字，见文件头“原卷存疑”②。
\item 设 $x\in\R$，$[x]$ 表示不大于 $x$ 的最大整数，如 $[\pi]=3$、$[-1.2]=-2$、$\left[\dfrac12\right]=0$，
则使 $[|x-1|]=3$ 成立的 $x$ 的取值范围 \blankline[4.4em].

\ifanswer
\solbegin
\step{由题意 $[|x-1|]=3$，则 $3\leqslant|x-1|<4$，所以 $3\leqslant x-1<4$ 或 $-4\leqslant x-1<-3$，}
\step{解得 $4\leqslant x<5$ 或 $-3<x\leqslant-2$，}
\step{所以使 $[|x-1|]=3$ 成立的 $x$ 的取值范围是 $(-3,-2]\cup[4,5)$.}
\solend

\fi

\end{enumerate}

\bigsec{二、选择题}

\begin{enumerate}
\setcounter{enumi}{10}

\item 若关于 $x$ 的方程 $k(x+2)-3(k-1)=x+1$ 有负数解，则实数 $k$ 的取值范围是\mbox{（\quad）}.

\keepwithprev
A. $(1,2)$\qquad B. $(-\infty,1)\cup(2,+\infty)$\qquad C. $(-\infty,2)$\qquad D. $(-\infty,1)$

\ans{$A$}

\item 与 $x^2>4$ 同解的不等式是\mbox{（\quad）}.

\keepwithprev
A. $x^2+\dfrac{1}{x-4}>4+\dfrac{1}{x-4}$\qquad B. $x^2+\dfrac1x>4+\dfrac1x$

C. $x^2+\sqrt{x-2}>4+\sqrt{x-2}$\qquad D. $\dfrac{1}{x^2-4}<0$

\ifanswer
\solbegin
\step{不等式 $x^2+\dfrac{1}{x-4}>4+\dfrac{1}{x-4}$ 等价于 $\begin{cases}x^2>4\\ x\neq4\end{cases}$，
与 $x^2>4$ 不同解，故选项 $A$ 错误；}
\step{不等式 $x^2+\dfrac1x>4+\dfrac1x$ 等价于 $\begin{cases}x^2>4\\ x\neq0\end{cases}$，
即 $x^2>4$，所以与 $x^2>4$ 同解，}
\step{故选项 $B$ 正确；}
\step{不等式 $x^2+\sqrt{x-2}>4+\sqrt{x-2}$ 等价于 $\begin{cases}x^2>4\\ x-2\geqslant0\end{cases}$，
与 $x^2>4$ 不同解，}
\step{故选项 $C$ 错误；}
\step{不等式 $\dfrac{1}{x^2-4}<0$ 等价于 $x^2<4$，与 $x^2>4$ 不同解，故选项 $D$ 错误.}
\step{故选 $B$.}
\solend

\fi

\item 若关于 $x$ 的二次不等式 $ax^2+bx+c\geqslant0(a\neq0)$ 的解集是 $\R$，那么\mbox{（\quad）}.

\keepwithprev
A. $a<0$，且 $b^2-4ac>0$\qquad B. $a<0$，且 $b^2-4ac\leqslant0$

% 注：原卷本题 D 选项与 A 选项文字完全重复（均作“$a<0$，且 $b^2-4ac>0$”），
%     照录未改，见文件头“原卷存疑”①。
C. $a>0$，且 $b^2-4ac\leqslant0$\qquad D. $a<0$，且 $b^2-4ac>0$

\ans{$C$}

\item 若不等式 $x^2+mx+1>2x+m$ 对满足 $|m|<2$ 的所有实数 $m$ 恒成立，
则实数 $x$ 的取值范围是\mbox{（\quad）}.

\keepwithprev
A. $-2<x<2$\qquad B. $x\geqslant3$\qquad C. $x\leqslant1$\qquad D. $x\leqslant-1$ 或 $x\geqslant3$

\ans{$D$}

\end{enumerate}

% “三、解答题”原本落在试卷版第 1 页页末、第 15 题翻到次页（切页审计 B 类）。
% 节标题自带的守卫只有 3\baselineskip，够放标题不够放标题＋首题，故在标题前再量一次页高。
\needvspace{4cm}
\bigsec{三、解答题}

\begin{enumerate}
\setcounter{enumi}{14}

\item 设关于 $x$ 的二次方程 $(m-1)x^2+(2m-4)x+m=0$ 有两个正根，求实数 $m$ 的取值范围.

\ifanswer
\solbegin
\step{$\begin{cases}\Delta=(2m-4)^2-4m(m-1)\geqslant0\\ -\dfrac{2m-4}{m-1}>0\\ \dfrac{m}{m-1}>0\end{cases}$，
解得 $1<m\leqslant\dfrac43$.}
\solend

\fi

\ansspace[6]

\item 已知关于 $x$ 的不等式 $2x^2-9x+a<0$ 的解集非空，对于其解集内的每一个 $x$ 的值，
至少能使不等式 $x^2-4x+3<0$ 或 $x^2-6x+8<0$ 中的一个成立，求实数 $a$ 的取值范围.

\ifanswer
\solbegin
\step{不等式 $x^2-4x+3<0$ 的解集为 $(1,3)$；$x^2-6x+8<0$ 的解集为 $(2,4)$.}
\step{关于 $x$ 的不等式 $2x^2-9x+a<0$ 的解集非空，所以 $\Delta=81-8a>0$，解得 $a<\dfrac{81}{8}$.}
\step{所以 $\dfrac{9-\sqrt{81-8a}}{4}<x<\dfrac{9+\sqrt{81-8a}}{4}$.}
\step{由于关于 $x$ 的不等式 $2x^2-9x+a<0$（解集非空）的每一个 $x$ 的值}
\step{至少满足不等式 $x^2-4x+3<0$ 和 $x^2-6x+8<0$ 中的一个，}
\step{所以 $\begin{cases}\dfrac{9-\sqrt{81-8a}}{4}\geqslant1\\ \dfrac{9+\sqrt{81-8a}}{4}\leqslant4\\ a<\dfrac{81}{8}\end{cases}$，
解得 $7\leqslant a<\dfrac{81}{8}$，所以实数 $a$ 的取值范围是 $\left[7,\dfrac{81}{8}\right)$.}
\solend

\fi

\ansspace[8]

\item 是否存在实数 $k$，使 $x=3$ 在不等式 $\dfrac{x+2}{k}>1+\dfrac{x^2-3}{k^2}$ 的解集中？
若存在，求 $k$ 的取值范围；若不存在，说明理由.

\ans{$(2,3)$}

\ansspace[6]

\item 关于 $x$ 的不等式 $x^2-(a+1)x+a<0$ 的解集为 $A$，$x^2-5x+4<0$ 的解集为 $B$，
若 $A\sub B$，求实数 $a$ 的取值范围.

\ans{$[1,4]$}

\ansspace[6]

\item 汽车在行驶中，由于惯性作用，刹车后还要继续向前滑行一段距离才能停住，我们称这段
距离为“刹车距离”. 刹车距离是分析事故的一个重要因素. 在一个限速 $40km/h$ 以内的弯道上，
甲、乙两辆汽车相向而行，发现情况不对，同时刹车，但还是相撞了，事后现场测得甲车的刹车
距离略超过 $12m$，乙车的刹车距离略超过 $10m$，又知甲、乙两种车型的刹车距离 $s(m)$ 与
车速 $x(km/h)$ 之间有如下关系：$s_{\text{甲}}=0.1x+0.01x^2$，$s_{\text{乙}}=0.05x+0.005x^2$. 问：
两车相撞的主要责任是谁？

\ifanswer
\solbegin
\step{由题意列出不等式组 $\begin{cases}0.1x+0.01x^2>12\\ 0.05x+0.005x^2>10\end{cases}$，
得 $\begin{cases}x<-40\text{ 或 }x>30\\ x<-50\text{ 或 }x>40\end{cases}$，}
\step{由于 $x>0$，从而得 $x_{\text{甲}}>30km/h$，$x_{\text{乙}}>40km/h$.}
\step{经比较得乙车超过限速，应负主要责任.}
\solend

\fi

\ansspace*[8]

\end{enumerate}

\end{document}
"
% 紧凑版：xelatex "\def\iscompact{1}% 高一40_大同中学_周末作业4.2.tex
%   —— 高一上数学“2026-2027 年大同中学高一上周末作业 4.2”（试卷版/答案版）
% 试卷版：xelatex 高一40_大同中学_周末作业4.2.tex
% 答案版：xelatex "\def\isanswer{1}% 高一40_大同中学_周末作业4.2.tex
%   —— 高一上数学“2026-2027 年大同中学高一上周末作业 4.2”（试卷版/答案版）
% 试卷版：xelatex 高一40_大同中学_周末作业4.2.tex
% 答案版：xelatex "\def\isanswer{1}% 高一40_大同中学_周末作业4.2.tex
%   —— 高一上数学“2026-2027 年大同中学高一上周末作业 4.2”（试卷版/答案版）
% 试卷版：xelatex 高一40_大同中学_周末作业4.2.tex
% 答案版：xelatex "\def\isanswer{1}\input{高一40_大同中学_周末作业4.2.tex}"
% 紧凑版：xelatex "\def\iscompact{1}\input{高一40_大同中学_周末作业4.2.tex}"
%
% 原始材料是微信公众号图文（公众号“魔都数学加油站”连载“27学年高一 40”，署名“破晓”，
% 2026-09-26 21:29 发布，原文标题“27学年高一40——2026-2027年上海市大同中学高一上周末作业4.2”，
% 原文链接 https://mp.weixin.qq.com/s/RrBSqDF_nJVbbx-bmmdtkQ ），正文为 4 张 PNG 页（4762×6735）。原文本身就是一份**讲评版**——题下直接印【答案】或【解析】，
% 第 11、13、14 题的答案还直接印在题干末尾的括号里。试题文字、答案与解析均照原卷录入，
% 未改内容；卷面的粉色斜水印（“魔都数学加油站”字样）属水印，未录入。
%
% 卷首标题：原卷页首印作“2026-2027 年大同中学高一上周末作业 4.2”（**卷面标题行即带校名**，
% 与高一16／高一18 等拍照卷须由本稿补校名的情形不同），本稿照卷面，年份区间按全稿惯例
% 排 en dash，前缀照原卷保留“年”。
%
% 篇幅与题号：本篇**只有 4 页**，卷面自“一、填空题”的**第 3 题**起编，终于“三、解答题”的
% **第 19 题**，共 17 题（填空 3—10、选择 11—14、解答 15—19）。第 1、2 题未在该文中发布——
% 这是该公众号连载讲评版的通例（高一02—高一09 与 高一36、高一38、高一39 等均自第 3 题起编，
% README“说明”已记），**不是截断**，故本稿题号亦自 3 起。它与 高一39（大同中学 周末作业4.1）
% **不构成同一份试卷的两半**：4.1 是另一次作业（11 页，题号**也自第 3 题**起、一直编到**第 24 题**，
% 以集合／常用逻辑用语为主、兼及不等式），两份各自独立起编（都从第 3 题开始），只是同为“大同中学高一上
% 周末作业 4”系列。
%
% 与原卷的排版差异（均已在 README“说明”中交代）：
%   ① 原文自**第 3 题**起（第 1、2 题未在该文中发布），故本稿题号亦自 3 起；
%   ② 原卷本就印有“一、填空题”“二、选择题”“三、解答题”三节标题，本稿照原卷分节，
%      只把标题改排成全稿统一的“大节标题”样式；
%   ③ 原卷为讲评版：第 4、5、6、7、8、17、18 题只印【答案】行，第 3、9、10、12、15、16、19 题
%      印【解析】（无单独的【答案】行），第 11、13、14 题的答案直接印在题干末尾的括号内；
%      本稿统一改为试卷版隐去、答案版以 \ans{}／\ifanswer 块给出，两版彻底分离（照 高一36 口径：
%      只印【解析】的题不另给 \ans{}，只印【答案】的题仅给 \ans{}）；
%   ④ 数集记号原卷作斜体的 $R$（第 5、10、13 题），本稿按全稿惯例统一写作 $\R$；
%   ⑤ 补集符号原卷作上横线（第 5 题的 $\overline{M}$），本稿照录，未改为 $\complement_\R$；
%   ⑥ 包含符号原卷两套并用：第 4 题作**无下划线**的 $\subset$（真包含），第 18 题作**带下划线**的
%      $\sub$（即 $\subseteq$），本稿照录、不归一；
%   ⑦ 原卷填空的横线约两个字宽，本稿按全稿惯例统一排 \blankline[4.4em]；
%   ⑧ 原卷未印副标题，本稿按**本卷实际内容**补出副标题“不等式”——全卷 17 题里约 15 题为
%      不等式（解一元二次／无理／分式不等式、恒成立与根的分布、含参不等式、应用题），
%      集合仅第 4、5 两题（第 18 题只是用集合关系表述解集），常用逻辑用语 0 题；
%   ⑨ 本卷**无插图**（全卷检索确认无“如图”）；
%   ⑩ 标点与单位：第 6 题的“$a,b$”本稿按全稿惯例排顿号“$a$、$b$”，第 10 题的三例
%      “$[\pi]=3,[-1.2]=-2,[\tfrac12]=0$”之间排顿号，第 5 题原卷“全集 $U=R$ 集合 $M$”之间
%      漏排标点、本稿补逗号；车速单位原卷印作斜体 $km/h$（第 19 题），本稿照录；半角括号、逗号、
%      问号一律排全角（第 17、19 题原卷的两个问号本就作全角，本稿照录）。
%
% 原卷存疑两处（照抄原卷、未改，见源码中该处上方注释）：
%   ① 第 13 题的 D 选项与 A 选项文字完全重复（均作“$a<0$，且 $b^2-4ac>0$”），原卷如此。
%   ② 第 10 题题干末作“…的取值范围____”，**漏排“是”**（本卷其余填空均作“…是____”）——
%      本稿照录、未补字，见源码该题上方 `%` 注。

\documentclass[a4paper,11pt]{article}
\usepackage{common}

\begin{document}

\examtitlest[3]{2026--2027 年大同中学高一上周末作业 4.2}{不等式}

\bigsec{一、填空题}

\begin{enumerate}
\setcounter{enumi}{2}

\item 在平面直角坐标系中，若点 $(2,b)$ 到原点的距离不小于 $5$，则 $b$ 的取值范围是 \blankline[4.4em].

\ifanswer
\solbegin
\step{因为点 $(2,b)$ 到原点的距离不小于 $5$，所以 $\sqrt{(2-0)^2+(b-0)^2}\geqslant5$，}
\step{所以 $4+b^2\geqslant25$，得 $b\geqslant\sqrt{21}$ 或 $b\leqslant-\sqrt{21}$，}
\step{则 $b$ 的取值范围是 $(-\infty,-\sqrt{21}]\cup[\sqrt{21},+\infty)$.}
\solend

\fi

\item 若集合 $A=\{x\mid|x|>3\}$，集合 $B=\{x\mid x\geqslant a\}$，且 $B\subset A$，
则实数 $a$ 的取值范围是 \blankline[4.4em].

\ans{$(3,+\infty)$}

% 注：原卷“全集 $U=R$”与“集合 $M$”之间**漏排标点**（只有空格），本稿补逗号——
%     照 高一b 第 13 题“原卷漏排、本稿补上”的先例（高一10、高一24 同例），见文件头⑩。
\item 若全集 $U=\R$，集合 $M=\{x\mid y=\sqrt{9-x^2}\}$，$S=\{x\mid x^2-3x+2\geqslant0\}$，
则 $\overline{M}\cup S=$ \blankline[4.4em].

\ans{$(-\infty,1]\cup[2,+\infty)$}

\item 设 $a$、$b$ 为实数，若不等式 $x^2-ax-b<0$ 的解集为 $\{x\mid2<x<8\}$，
则 $bx^2-ax-1>0$ 的解集为 \blankline[4.4em].

\ans{$\left(-\dfrac12,-\dfrac18\right)$}

\item 不等式 $4x-13+\sqrt{x+5}\geqslant8x-1+\sqrt{x+5}$ 的解集是 \blankline[4.4em].

\ans{$[-5,-3]$}

\item 若关于 $x$ 的不等式 $ax^2-2ax+3\leqslant0$ 的解集为 $\varnothing$，
则实数 $a$ 的取值范围是 \blankline[4.4em].

\ans{$[0,3)$}

\item 若不等式 $x^2-ax-2<0$ 在 $x\in(1,2)$ 内恒成立，则实数 $a$ 的取值范围是 \blankline[4.4em].

\ifanswer
\solbegin
\step{由不等式 $x^2-ax-2<0$ 在 $x\in(1,2)$ 内恒成立，}
\step{得 $ax>x^2-2$，即 $a>x-\dfrac2x$ 在 $x\in(1,2)$ 内恒成立，}
\step{令 $t(x)=x-\dfrac2x$，$x\in(1,2)$，该函数为严格增函数，则 $t(x)<t(2)=1$，得 $a\geqslant1$，}
\step{所以 $a$ 的取值范围是 $[1,+\infty)$.}
\solend

\fi

% 注：原卷此处作“的取值范围\_\_\_\_”，**漏排“是”**（本卷其余填空均作“…是\_\_\_\_”），
%     本稿照录、未补字，见文件头“原卷存疑”②。
\item 设 $x\in\R$，$[x]$ 表示不大于 $x$ 的最大整数，如 $[\pi]=3$、$[-1.2]=-2$、$\left[\dfrac12\right]=0$，
则使 $[|x-1|]=3$ 成立的 $x$ 的取值范围 \blankline[4.4em].

\ifanswer
\solbegin
\step{由题意 $[|x-1|]=3$，则 $3\leqslant|x-1|<4$，所以 $3\leqslant x-1<4$ 或 $-4\leqslant x-1<-3$，}
\step{解得 $4\leqslant x<5$ 或 $-3<x\leqslant-2$，}
\step{所以使 $[|x-1|]=3$ 成立的 $x$ 的取值范围是 $(-3,-2]\cup[4,5)$.}
\solend

\fi

\end{enumerate}

\bigsec{二、选择题}

\begin{enumerate}
\setcounter{enumi}{10}

\item 若关于 $x$ 的方程 $k(x+2)-3(k-1)=x+1$ 有负数解，则实数 $k$ 的取值范围是\mbox{（\quad）}.

\keepwithprev
A. $(1,2)$\qquad B. $(-\infty,1)\cup(2,+\infty)$\qquad C. $(-\infty,2)$\qquad D. $(-\infty,1)$

\ans{$A$}

\item 与 $x^2>4$ 同解的不等式是\mbox{（\quad）}.

\keepwithprev
A. $x^2+\dfrac{1}{x-4}>4+\dfrac{1}{x-4}$\qquad B. $x^2+\dfrac1x>4+\dfrac1x$

C. $x^2+\sqrt{x-2}>4+\sqrt{x-2}$\qquad D. $\dfrac{1}{x^2-4}<0$

\ifanswer
\solbegin
\step{不等式 $x^2+\dfrac{1}{x-4}>4+\dfrac{1}{x-4}$ 等价于 $\begin{cases}x^2>4\\ x\neq4\end{cases}$，
与 $x^2>4$ 不同解，故选项 $A$ 错误；}
\step{不等式 $x^2+\dfrac1x>4+\dfrac1x$ 等价于 $\begin{cases}x^2>4\\ x\neq0\end{cases}$，
即 $x^2>4$，所以与 $x^2>4$ 同解，}
\step{故选项 $B$ 正确；}
\step{不等式 $x^2+\sqrt{x-2}>4+\sqrt{x-2}$ 等价于 $\begin{cases}x^2>4\\ x-2\geqslant0\end{cases}$，
与 $x^2>4$ 不同解，}
\step{故选项 $C$ 错误；}
\step{不等式 $\dfrac{1}{x^2-4}<0$ 等价于 $x^2<4$，与 $x^2>4$ 不同解，故选项 $D$ 错误.}
\step{故选 $B$.}
\solend

\fi

\item 若关于 $x$ 的二次不等式 $ax^2+bx+c\geqslant0(a\neq0)$ 的解集是 $\R$，那么\mbox{（\quad）}.

\keepwithprev
A. $a<0$，且 $b^2-4ac>0$\qquad B. $a<0$，且 $b^2-4ac\leqslant0$

% 注：原卷本题 D 选项与 A 选项文字完全重复（均作“$a<0$，且 $b^2-4ac>0$”），
%     照录未改，见文件头“原卷存疑”①。
C. $a>0$，且 $b^2-4ac\leqslant0$\qquad D. $a<0$，且 $b^2-4ac>0$

\ans{$C$}

\item 若不等式 $x^2+mx+1>2x+m$ 对满足 $|m|<2$ 的所有实数 $m$ 恒成立，
则实数 $x$ 的取值范围是\mbox{（\quad）}.

\keepwithprev
A. $-2<x<2$\qquad B. $x\geqslant3$\qquad C. $x\leqslant1$\qquad D. $x\leqslant-1$ 或 $x\geqslant3$

\ans{$D$}

\end{enumerate}

% “三、解答题”原本落在试卷版第 1 页页末、第 15 题翻到次页（切页审计 B 类）。
% 节标题自带的守卫只有 3\baselineskip，够放标题不够放标题＋首题，故在标题前再量一次页高。
\needvspace{4cm}
\bigsec{三、解答题}

\begin{enumerate}
\setcounter{enumi}{14}

\item 设关于 $x$ 的二次方程 $(m-1)x^2+(2m-4)x+m=0$ 有两个正根，求实数 $m$ 的取值范围.

\ifanswer
\solbegin
\step{$\begin{cases}\Delta=(2m-4)^2-4m(m-1)\geqslant0\\ -\dfrac{2m-4}{m-1}>0\\ \dfrac{m}{m-1}>0\end{cases}$，
解得 $1<m\leqslant\dfrac43$.}
\solend

\fi

\ansspace[6]

\item 已知关于 $x$ 的不等式 $2x^2-9x+a<0$ 的解集非空，对于其解集内的每一个 $x$ 的值，
至少能使不等式 $x^2-4x+3<0$ 或 $x^2-6x+8<0$ 中的一个成立，求实数 $a$ 的取值范围.

\ifanswer
\solbegin
\step{不等式 $x^2-4x+3<0$ 的解集为 $(1,3)$；$x^2-6x+8<0$ 的解集为 $(2,4)$.}
\step{关于 $x$ 的不等式 $2x^2-9x+a<0$ 的解集非空，所以 $\Delta=81-8a>0$，解得 $a<\dfrac{81}{8}$.}
\step{所以 $\dfrac{9-\sqrt{81-8a}}{4}<x<\dfrac{9+\sqrt{81-8a}}{4}$.}
\step{由于关于 $x$ 的不等式 $2x^2-9x+a<0$（解集非空）的每一个 $x$ 的值}
\step{至少满足不等式 $x^2-4x+3<0$ 和 $x^2-6x+8<0$ 中的一个，}
\step{所以 $\begin{cases}\dfrac{9-\sqrt{81-8a}}{4}\geqslant1\\ \dfrac{9+\sqrt{81-8a}}{4}\leqslant4\\ a<\dfrac{81}{8}\end{cases}$，
解得 $7\leqslant a<\dfrac{81}{8}$，所以实数 $a$ 的取值范围是 $\left[7,\dfrac{81}{8}\right)$.}
\solend

\fi

\ansspace[8]

\item 是否存在实数 $k$，使 $x=3$ 在不等式 $\dfrac{x+2}{k}>1+\dfrac{x^2-3}{k^2}$ 的解集中？
若存在，求 $k$ 的取值范围；若不存在，说明理由.

\ans{$(2,3)$}

\ansspace[6]

\item 关于 $x$ 的不等式 $x^2-(a+1)x+a<0$ 的解集为 $A$，$x^2-5x+4<0$ 的解集为 $B$，
若 $A\sub B$，求实数 $a$ 的取值范围.

\ans{$[1,4]$}

\ansspace[6]

\item 汽车在行驶中，由于惯性作用，刹车后还要继续向前滑行一段距离才能停住，我们称这段
距离为“刹车距离”. 刹车距离是分析事故的一个重要因素. 在一个限速 $40km/h$ 以内的弯道上，
甲、乙两辆汽车相向而行，发现情况不对，同时刹车，但还是相撞了，事后现场测得甲车的刹车
距离略超过 $12m$，乙车的刹车距离略超过 $10m$，又知甲、乙两种车型的刹车距离 $s(m)$ 与
车速 $x(km/h)$ 之间有如下关系：$s_{\text{甲}}=0.1x+0.01x^2$，$s_{\text{乙}}=0.05x+0.005x^2$. 问：
两车相撞的主要责任是谁？

\ifanswer
\solbegin
\step{由题意列出不等式组 $\begin{cases}0.1x+0.01x^2>12\\ 0.05x+0.005x^2>10\end{cases}$，
得 $\begin{cases}x<-40\text{ 或 }x>30\\ x<-50\text{ 或 }x>40\end{cases}$，}
\step{由于 $x>0$，从而得 $x_{\text{甲}}>30km/h$，$x_{\text{乙}}>40km/h$.}
\step{经比较得乙车超过限速，应负主要责任.}
\solend

\fi

\ansspace*[8]

\end{enumerate}

\end{document}
"
% 紧凑版：xelatex "\def\iscompact{1}% 高一40_大同中学_周末作业4.2.tex
%   —— 高一上数学“2026-2027 年大同中学高一上周末作业 4.2”（试卷版/答案版）
% 试卷版：xelatex 高一40_大同中学_周末作业4.2.tex
% 答案版：xelatex "\def\isanswer{1}\input{高一40_大同中学_周末作业4.2.tex}"
% 紧凑版：xelatex "\def\iscompact{1}\input{高一40_大同中学_周末作业4.2.tex}"
%
% 原始材料是微信公众号图文（公众号“魔都数学加油站”连载“27学年高一 40”，署名“破晓”，
% 2026-09-26 21:29 发布，原文标题“27学年高一40——2026-2027年上海市大同中学高一上周末作业4.2”，
% 原文链接 https://mp.weixin.qq.com/s/RrBSqDF_nJVbbx-bmmdtkQ ），正文为 4 张 PNG 页（4762×6735）。原文本身就是一份**讲评版**——题下直接印【答案】或【解析】，
% 第 11、13、14 题的答案还直接印在题干末尾的括号里。试题文字、答案与解析均照原卷录入，
% 未改内容；卷面的粉色斜水印（“魔都数学加油站”字样）属水印，未录入。
%
% 卷首标题：原卷页首印作“2026-2027 年大同中学高一上周末作业 4.2”（**卷面标题行即带校名**，
% 与高一16／高一18 等拍照卷须由本稿补校名的情形不同），本稿照卷面，年份区间按全稿惯例
% 排 en dash，前缀照原卷保留“年”。
%
% 篇幅与题号：本篇**只有 4 页**，卷面自“一、填空题”的**第 3 题**起编，终于“三、解答题”的
% **第 19 题**，共 17 题（填空 3—10、选择 11—14、解答 15—19）。第 1、2 题未在该文中发布——
% 这是该公众号连载讲评版的通例（高一02—高一09 与 高一36、高一38、高一39 等均自第 3 题起编，
% README“说明”已记），**不是截断**，故本稿题号亦自 3 起。它与 高一39（大同中学 周末作业4.1）
% **不构成同一份试卷的两半**：4.1 是另一次作业（11 页，题号**也自第 3 题**起、一直编到**第 24 题**，
% 以集合／常用逻辑用语为主、兼及不等式），两份各自独立起编（都从第 3 题开始），只是同为“大同中学高一上
% 周末作业 4”系列。
%
% 与原卷的排版差异（均已在 README“说明”中交代）：
%   ① 原文自**第 3 题**起（第 1、2 题未在该文中发布），故本稿题号亦自 3 起；
%   ② 原卷本就印有“一、填空题”“二、选择题”“三、解答题”三节标题，本稿照原卷分节，
%      只把标题改排成全稿统一的“大节标题”样式；
%   ③ 原卷为讲评版：第 4、5、6、7、8、17、18 题只印【答案】行，第 3、9、10、12、15、16、19 题
%      印【解析】（无单独的【答案】行），第 11、13、14 题的答案直接印在题干末尾的括号内；
%      本稿统一改为试卷版隐去、答案版以 \ans{}／\ifanswer 块给出，两版彻底分离（照 高一36 口径：
%      只印【解析】的题不另给 \ans{}，只印【答案】的题仅给 \ans{}）；
%   ④ 数集记号原卷作斜体的 $R$（第 5、10、13 题），本稿按全稿惯例统一写作 $\R$；
%   ⑤ 补集符号原卷作上横线（第 5 题的 $\overline{M}$），本稿照录，未改为 $\complement_\R$；
%   ⑥ 包含符号原卷两套并用：第 4 题作**无下划线**的 $\subset$（真包含），第 18 题作**带下划线**的
%      $\sub$（即 $\subseteq$），本稿照录、不归一；
%   ⑦ 原卷填空的横线约两个字宽，本稿按全稿惯例统一排 \blankline[4.4em]；
%   ⑧ 原卷未印副标题，本稿按**本卷实际内容**补出副标题“不等式”——全卷 17 题里约 15 题为
%      不等式（解一元二次／无理／分式不等式、恒成立与根的分布、含参不等式、应用题），
%      集合仅第 4、5 两题（第 18 题只是用集合关系表述解集），常用逻辑用语 0 题；
%   ⑨ 本卷**无插图**（全卷检索确认无“如图”）；
%   ⑩ 标点与单位：第 6 题的“$a,b$”本稿按全稿惯例排顿号“$a$、$b$”，第 10 题的三例
%      “$[\pi]=3,[-1.2]=-2,[\tfrac12]=0$”之间排顿号，第 5 题原卷“全集 $U=R$ 集合 $M$”之间
%      漏排标点、本稿补逗号；车速单位原卷印作斜体 $km/h$（第 19 题），本稿照录；半角括号、逗号、
%      问号一律排全角（第 17、19 题原卷的两个问号本就作全角，本稿照录）。
%
% 原卷存疑两处（照抄原卷、未改，见源码中该处上方注释）：
%   ① 第 13 题的 D 选项与 A 选项文字完全重复（均作“$a<0$，且 $b^2-4ac>0$”），原卷如此。
%   ② 第 10 题题干末作“…的取值范围____”，**漏排“是”**（本卷其余填空均作“…是____”）——
%      本稿照录、未补字，见源码该题上方 `%` 注。

\documentclass[a4paper,11pt]{article}
\usepackage{common}

\begin{document}

\examtitlest[3]{2026--2027 年大同中学高一上周末作业 4.2}{不等式}

\bigsec{一、填空题}

\begin{enumerate}
\setcounter{enumi}{2}

\item 在平面直角坐标系中，若点 $(2,b)$ 到原点的距离不小于 $5$，则 $b$ 的取值范围是 \blankline[4.4em].

\ifanswer
\solbegin
\step{因为点 $(2,b)$ 到原点的距离不小于 $5$，所以 $\sqrt{(2-0)^2+(b-0)^2}\geqslant5$，}
\step{所以 $4+b^2\geqslant25$，得 $b\geqslant\sqrt{21}$ 或 $b\leqslant-\sqrt{21}$，}
\step{则 $b$ 的取值范围是 $(-\infty,-\sqrt{21}]\cup[\sqrt{21},+\infty)$.}
\solend

\fi

\item 若集合 $A=\{x\mid|x|>3\}$，集合 $B=\{x\mid x\geqslant a\}$，且 $B\subset A$，
则实数 $a$ 的取值范围是 \blankline[4.4em].

\ans{$(3,+\infty)$}

% 注：原卷“全集 $U=R$”与“集合 $M$”之间**漏排标点**（只有空格），本稿补逗号——
%     照 高一b 第 13 题“原卷漏排、本稿补上”的先例（高一10、高一24 同例），见文件头⑩。
\item 若全集 $U=\R$，集合 $M=\{x\mid y=\sqrt{9-x^2}\}$，$S=\{x\mid x^2-3x+2\geqslant0\}$，
则 $\overline{M}\cup S=$ \blankline[4.4em].

\ans{$(-\infty,1]\cup[2,+\infty)$}

\item 设 $a$、$b$ 为实数，若不等式 $x^2-ax-b<0$ 的解集为 $\{x\mid2<x<8\}$，
则 $bx^2-ax-1>0$ 的解集为 \blankline[4.4em].

\ans{$\left(-\dfrac12,-\dfrac18\right)$}

\item 不等式 $4x-13+\sqrt{x+5}\geqslant8x-1+\sqrt{x+5}$ 的解集是 \blankline[4.4em].

\ans{$[-5,-3]$}

\item 若关于 $x$ 的不等式 $ax^2-2ax+3\leqslant0$ 的解集为 $\varnothing$，
则实数 $a$ 的取值范围是 \blankline[4.4em].

\ans{$[0,3)$}

\item 若不等式 $x^2-ax-2<0$ 在 $x\in(1,2)$ 内恒成立，则实数 $a$ 的取值范围是 \blankline[4.4em].

\ifanswer
\solbegin
\step{由不等式 $x^2-ax-2<0$ 在 $x\in(1,2)$ 内恒成立，}
\step{得 $ax>x^2-2$，即 $a>x-\dfrac2x$ 在 $x\in(1,2)$ 内恒成立，}
\step{令 $t(x)=x-\dfrac2x$，$x\in(1,2)$，该函数为严格增函数，则 $t(x)<t(2)=1$，得 $a\geqslant1$，}
\step{所以 $a$ 的取值范围是 $[1,+\infty)$.}
\solend

\fi

% 注：原卷此处作“的取值范围\_\_\_\_”，**漏排“是”**（本卷其余填空均作“…是\_\_\_\_”），
%     本稿照录、未补字，见文件头“原卷存疑”②。
\item 设 $x\in\R$，$[x]$ 表示不大于 $x$ 的最大整数，如 $[\pi]=3$、$[-1.2]=-2$、$\left[\dfrac12\right]=0$，
则使 $[|x-1|]=3$ 成立的 $x$ 的取值范围 \blankline[4.4em].

\ifanswer
\solbegin
\step{由题意 $[|x-1|]=3$，则 $3\leqslant|x-1|<4$，所以 $3\leqslant x-1<4$ 或 $-4\leqslant x-1<-3$，}
\step{解得 $4\leqslant x<5$ 或 $-3<x\leqslant-2$，}
\step{所以使 $[|x-1|]=3$ 成立的 $x$ 的取值范围是 $(-3,-2]\cup[4,5)$.}
\solend

\fi

\end{enumerate}

\bigsec{二、选择题}

\begin{enumerate}
\setcounter{enumi}{10}

\item 若关于 $x$ 的方程 $k(x+2)-3(k-1)=x+1$ 有负数解，则实数 $k$ 的取值范围是\mbox{（\quad）}.

\keepwithprev
A. $(1,2)$\qquad B. $(-\infty,1)\cup(2,+\infty)$\qquad C. $(-\infty,2)$\qquad D. $(-\infty,1)$

\ans{$A$}

\item 与 $x^2>4$ 同解的不等式是\mbox{（\quad）}.

\keepwithprev
A. $x^2+\dfrac{1}{x-4}>4+\dfrac{1}{x-4}$\qquad B. $x^2+\dfrac1x>4+\dfrac1x$

C. $x^2+\sqrt{x-2}>4+\sqrt{x-2}$\qquad D. $\dfrac{1}{x^2-4}<0$

\ifanswer
\solbegin
\step{不等式 $x^2+\dfrac{1}{x-4}>4+\dfrac{1}{x-4}$ 等价于 $\begin{cases}x^2>4\\ x\neq4\end{cases}$，
与 $x^2>4$ 不同解，故选项 $A$ 错误；}
\step{不等式 $x^2+\dfrac1x>4+\dfrac1x$ 等价于 $\begin{cases}x^2>4\\ x\neq0\end{cases}$，
即 $x^2>4$，所以与 $x^2>4$ 同解，}
\step{故选项 $B$ 正确；}
\step{不等式 $x^2+\sqrt{x-2}>4+\sqrt{x-2}$ 等价于 $\begin{cases}x^2>4\\ x-2\geqslant0\end{cases}$，
与 $x^2>4$ 不同解，}
\step{故选项 $C$ 错误；}
\step{不等式 $\dfrac{1}{x^2-4}<0$ 等价于 $x^2<4$，与 $x^2>4$ 不同解，故选项 $D$ 错误.}
\step{故选 $B$.}
\solend

\fi

\item 若关于 $x$ 的二次不等式 $ax^2+bx+c\geqslant0(a\neq0)$ 的解集是 $\R$，那么\mbox{（\quad）}.

\keepwithprev
A. $a<0$，且 $b^2-4ac>0$\qquad B. $a<0$，且 $b^2-4ac\leqslant0$

% 注：原卷本题 D 选项与 A 选项文字完全重复（均作“$a<0$，且 $b^2-4ac>0$”），
%     照录未改，见文件头“原卷存疑”①。
C. $a>0$，且 $b^2-4ac\leqslant0$\qquad D. $a<0$，且 $b^2-4ac>0$

\ans{$C$}

\item 若不等式 $x^2+mx+1>2x+m$ 对满足 $|m|<2$ 的所有实数 $m$ 恒成立，
则实数 $x$ 的取值范围是\mbox{（\quad）}.

\keepwithprev
A. $-2<x<2$\qquad B. $x\geqslant3$\qquad C. $x\leqslant1$\qquad D. $x\leqslant-1$ 或 $x\geqslant3$

\ans{$D$}

\end{enumerate}

% “三、解答题”原本落在试卷版第 1 页页末、第 15 题翻到次页（切页审计 B 类）。
% 节标题自带的守卫只有 3\baselineskip，够放标题不够放标题＋首题，故在标题前再量一次页高。
\needvspace{4cm}
\bigsec{三、解答题}

\begin{enumerate}
\setcounter{enumi}{14}

\item 设关于 $x$ 的二次方程 $(m-1)x^2+(2m-4)x+m=0$ 有两个正根，求实数 $m$ 的取值范围.

\ifanswer
\solbegin
\step{$\begin{cases}\Delta=(2m-4)^2-4m(m-1)\geqslant0\\ -\dfrac{2m-4}{m-1}>0\\ \dfrac{m}{m-1}>0\end{cases}$，
解得 $1<m\leqslant\dfrac43$.}
\solend

\fi

\ansspace[6]

\item 已知关于 $x$ 的不等式 $2x^2-9x+a<0$ 的解集非空，对于其解集内的每一个 $x$ 的值，
至少能使不等式 $x^2-4x+3<0$ 或 $x^2-6x+8<0$ 中的一个成立，求实数 $a$ 的取值范围.

\ifanswer
\solbegin
\step{不等式 $x^2-4x+3<0$ 的解集为 $(1,3)$；$x^2-6x+8<0$ 的解集为 $(2,4)$.}
\step{关于 $x$ 的不等式 $2x^2-9x+a<0$ 的解集非空，所以 $\Delta=81-8a>0$，解得 $a<\dfrac{81}{8}$.}
\step{所以 $\dfrac{9-\sqrt{81-8a}}{4}<x<\dfrac{9+\sqrt{81-8a}}{4}$.}
\step{由于关于 $x$ 的不等式 $2x^2-9x+a<0$（解集非空）的每一个 $x$ 的值}
\step{至少满足不等式 $x^2-4x+3<0$ 和 $x^2-6x+8<0$ 中的一个，}
\step{所以 $\begin{cases}\dfrac{9-\sqrt{81-8a}}{4}\geqslant1\\ \dfrac{9+\sqrt{81-8a}}{4}\leqslant4\\ a<\dfrac{81}{8}\end{cases}$，
解得 $7\leqslant a<\dfrac{81}{8}$，所以实数 $a$ 的取值范围是 $\left[7,\dfrac{81}{8}\right)$.}
\solend

\fi

\ansspace[8]

\item 是否存在实数 $k$，使 $x=3$ 在不等式 $\dfrac{x+2}{k}>1+\dfrac{x^2-3}{k^2}$ 的解集中？
若存在，求 $k$ 的取值范围；若不存在，说明理由.

\ans{$(2,3)$}

\ansspace[6]

\item 关于 $x$ 的不等式 $x^2-(a+1)x+a<0$ 的解集为 $A$，$x^2-5x+4<0$ 的解集为 $B$，
若 $A\sub B$，求实数 $a$ 的取值范围.

\ans{$[1,4]$}

\ansspace[6]

\item 汽车在行驶中，由于惯性作用，刹车后还要继续向前滑行一段距离才能停住，我们称这段
距离为“刹车距离”. 刹车距离是分析事故的一个重要因素. 在一个限速 $40km/h$ 以内的弯道上，
甲、乙两辆汽车相向而行，发现情况不对，同时刹车，但还是相撞了，事后现场测得甲车的刹车
距离略超过 $12m$，乙车的刹车距离略超过 $10m$，又知甲、乙两种车型的刹车距离 $s(m)$ 与
车速 $x(km/h)$ 之间有如下关系：$s_{\text{甲}}=0.1x+0.01x^2$，$s_{\text{乙}}=0.05x+0.005x^2$. 问：
两车相撞的主要责任是谁？

\ifanswer
\solbegin
\step{由题意列出不等式组 $\begin{cases}0.1x+0.01x^2>12\\ 0.05x+0.005x^2>10\end{cases}$，
得 $\begin{cases}x<-40\text{ 或 }x>30\\ x<-50\text{ 或 }x>40\end{cases}$，}
\step{由于 $x>0$，从而得 $x_{\text{甲}}>30km/h$，$x_{\text{乙}}>40km/h$.}
\step{经比较得乙车超过限速，应负主要责任.}
\solend

\fi

\ansspace*[8]

\end{enumerate}

\end{document}
"
%
% 原始材料是微信公众号图文（公众号“魔都数学加油站”连载“27学年高一 40”，署名“破晓”，
% 2026-09-26 21:29 发布，原文标题“27学年高一40——2026-2027年上海市大同中学高一上周末作业4.2”，
% 原文链接 https://mp.weixin.qq.com/s/RrBSqDF_nJVbbx-bmmdtkQ ），正文为 4 张 PNG 页（4762×6735）。原文本身就是一份**讲评版**——题下直接印【答案】或【解析】，
% 第 11、13、14 题的答案还直接印在题干末尾的括号里。试题文字、答案与解析均照原卷录入，
% 未改内容；卷面的粉色斜水印（“魔都数学加油站”字样）属水印，未录入。
%
% 卷首标题：原卷页首印作“2026-2027 年大同中学高一上周末作业 4.2”（**卷面标题行即带校名**，
% 与高一16／高一18 等拍照卷须由本稿补校名的情形不同），本稿照卷面，年份区间按全稿惯例
% 排 en dash，前缀照原卷保留“年”。
%
% 篇幅与题号：本篇**只有 4 页**，卷面自“一、填空题”的**第 3 题**起编，终于“三、解答题”的
% **第 19 题**，共 17 题（填空 3—10、选择 11—14、解答 15—19）。第 1、2 题未在该文中发布——
% 这是该公众号连载讲评版的通例（高一02—高一09 与 高一36、高一38、高一39 等均自第 3 题起编，
% README“说明”已记），**不是截断**，故本稿题号亦自 3 起。它与 高一39（大同中学 周末作业4.1）
% **不构成同一份试卷的两半**：4.1 是另一次作业（11 页，题号**也自第 3 题**起、一直编到**第 24 题**，
% 以集合／常用逻辑用语为主、兼及不等式），两份各自独立起编（都从第 3 题开始），只是同为“大同中学高一上
% 周末作业 4”系列。
%
% 与原卷的排版差异（均已在 README“说明”中交代）：
%   ① 原文自**第 3 题**起（第 1、2 题未在该文中发布），故本稿题号亦自 3 起；
%   ② 原卷本就印有“一、填空题”“二、选择题”“三、解答题”三节标题，本稿照原卷分节，
%      只把标题改排成全稿统一的“大节标题”样式；
%   ③ 原卷为讲评版：第 4、5、6、7、8、17、18 题只印【答案】行，第 3、9、10、12、15、16、19 题
%      印【解析】（无单独的【答案】行），第 11、13、14 题的答案直接印在题干末尾的括号内；
%      本稿统一改为试卷版隐去、答案版以 \ans{}／\ifanswer 块给出，两版彻底分离（照 高一36 口径：
%      只印【解析】的题不另给 \ans{}，只印【答案】的题仅给 \ans{}）；
%   ④ 数集记号原卷作斜体的 $R$（第 5、10、13 题），本稿按全稿惯例统一写作 $\R$；
%   ⑤ 补集符号原卷作上横线（第 5 题的 $\overline{M}$），本稿照录，未改为 $\complement_\R$；
%   ⑥ 包含符号原卷两套并用：第 4 题作**无下划线**的 $\subset$（真包含），第 18 题作**带下划线**的
%      $\sub$（即 $\subseteq$），本稿照录、不归一；
%   ⑦ 原卷填空的横线约两个字宽，本稿按全稿惯例统一排 \blankline[4.4em]；
%   ⑧ 原卷未印副标题，本稿按**本卷实际内容**补出副标题“不等式”——全卷 17 题里约 15 题为
%      不等式（解一元二次／无理／分式不等式、恒成立与根的分布、含参不等式、应用题），
%      集合仅第 4、5 两题（第 18 题只是用集合关系表述解集），常用逻辑用语 0 题；
%   ⑨ 本卷**无插图**（全卷检索确认无“如图”）；
%   ⑩ 标点与单位：第 6 题的“$a,b$”本稿按全稿惯例排顿号“$a$、$b$”，第 10 题的三例
%      “$[\pi]=3,[-1.2]=-2,[\tfrac12]=0$”之间排顿号，第 5 题原卷“全集 $U=R$ 集合 $M$”之间
%      漏排标点、本稿补逗号；车速单位原卷印作斜体 $km/h$（第 19 题），本稿照录；半角括号、逗号、
%      问号一律排全角（第 17、19 题原卷的两个问号本就作全角，本稿照录）。
%
% 原卷存疑两处（照抄原卷、未改，见源码中该处上方注释）：
%   ① 第 13 题的 D 选项与 A 选项文字完全重复（均作“$a<0$，且 $b^2-4ac>0$”），原卷如此。
%   ② 第 10 题题干末作“…的取值范围____”，**漏排“是”**（本卷其余填空均作“…是____”）——
%      本稿照录、未补字，见源码该题上方 `%` 注。

\documentclass[a4paper,11pt]{article}
\usepackage{common}

\begin{document}

\examtitlest[3]{2026--2027 年大同中学高一上周末作业 4.2}{不等式}

\bigsec{一、填空题}

\begin{enumerate}
\setcounter{enumi}{2}

\item 在平面直角坐标系中，若点 $(2,b)$ 到原点的距离不小于 $5$，则 $b$ 的取值范围是 \blankline[4.4em].

\ifanswer
\solbegin
\step{因为点 $(2,b)$ 到原点的距离不小于 $5$，所以 $\sqrt{(2-0)^2+(b-0)^2}\geqslant5$，}
\step{所以 $4+b^2\geqslant25$，得 $b\geqslant\sqrt{21}$ 或 $b\leqslant-\sqrt{21}$，}
\step{则 $b$ 的取值范围是 $(-\infty,-\sqrt{21}]\cup[\sqrt{21},+\infty)$.}
\solend

\fi

\item 若集合 $A=\{x\mid|x|>3\}$，集合 $B=\{x\mid x\geqslant a\}$，且 $B\subset A$，
则实数 $a$ 的取值范围是 \blankline[4.4em].

\ans{$(3,+\infty)$}

% 注：原卷“全集 $U=R$”与“集合 $M$”之间**漏排标点**（只有空格），本稿补逗号——
%     照 高一b 第 13 题“原卷漏排、本稿补上”的先例（高一10、高一24 同例），见文件头⑩。
\item 若全集 $U=\R$，集合 $M=\{x\mid y=\sqrt{9-x^2}\}$，$S=\{x\mid x^2-3x+2\geqslant0\}$，
则 $\overline{M}\cup S=$ \blankline[4.4em].

\ans{$(-\infty,1]\cup[2,+\infty)$}

\item 设 $a$、$b$ 为实数，若不等式 $x^2-ax-b<0$ 的解集为 $\{x\mid2<x<8\}$，
则 $bx^2-ax-1>0$ 的解集为 \blankline[4.4em].

\ans{$\left(-\dfrac12,-\dfrac18\right)$}

\item 不等式 $4x-13+\sqrt{x+5}\geqslant8x-1+\sqrt{x+5}$ 的解集是 \blankline[4.4em].

\ans{$[-5,-3]$}

\item 若关于 $x$ 的不等式 $ax^2-2ax+3\leqslant0$ 的解集为 $\varnothing$，
则实数 $a$ 的取值范围是 \blankline[4.4em].

\ans{$[0,3)$}

\item 若不等式 $x^2-ax-2<0$ 在 $x\in(1,2)$ 内恒成立，则实数 $a$ 的取值范围是 \blankline[4.4em].

\ifanswer
\solbegin
\step{由不等式 $x^2-ax-2<0$ 在 $x\in(1,2)$ 内恒成立，}
\step{得 $ax>x^2-2$，即 $a>x-\dfrac2x$ 在 $x\in(1,2)$ 内恒成立，}
\step{令 $t(x)=x-\dfrac2x$，$x\in(1,2)$，该函数为严格增函数，则 $t(x)<t(2)=1$，得 $a\geqslant1$，}
\step{所以 $a$ 的取值范围是 $[1,+\infty)$.}
\solend

\fi

% 注：原卷此处作“的取值范围\_\_\_\_”，**漏排“是”**（本卷其余填空均作“…是\_\_\_\_”），
%     本稿照录、未补字，见文件头“原卷存疑”②。
\item 设 $x\in\R$，$[x]$ 表示不大于 $x$ 的最大整数，如 $[\pi]=3$、$[-1.2]=-2$、$\left[\dfrac12\right]=0$，
则使 $[|x-1|]=3$ 成立的 $x$ 的取值范围 \blankline[4.4em].

\ifanswer
\solbegin
\step{由题意 $[|x-1|]=3$，则 $3\leqslant|x-1|<4$，所以 $3\leqslant x-1<4$ 或 $-4\leqslant x-1<-3$，}
\step{解得 $4\leqslant x<5$ 或 $-3<x\leqslant-2$，}
\step{所以使 $[|x-1|]=3$ 成立的 $x$ 的取值范围是 $(-3,-2]\cup[4,5)$.}
\solend

\fi

\end{enumerate}

\bigsec{二、选择题}

\begin{enumerate}
\setcounter{enumi}{10}

\item 若关于 $x$ 的方程 $k(x+2)-3(k-1)=x+1$ 有负数解，则实数 $k$ 的取值范围是\mbox{（\quad）}.

\keepwithprev
A. $(1,2)$\qquad B. $(-\infty,1)\cup(2,+\infty)$\qquad C. $(-\infty,2)$\qquad D. $(-\infty,1)$

\ans{$A$}

\item 与 $x^2>4$ 同解的不等式是\mbox{（\quad）}.

\keepwithprev
A. $x^2+\dfrac{1}{x-4}>4+\dfrac{1}{x-4}$\qquad B. $x^2+\dfrac1x>4+\dfrac1x$

C. $x^2+\sqrt{x-2}>4+\sqrt{x-2}$\qquad D. $\dfrac{1}{x^2-4}<0$

\ifanswer
\solbegin
\step{不等式 $x^2+\dfrac{1}{x-4}>4+\dfrac{1}{x-4}$ 等价于 $\begin{cases}x^2>4\\ x\neq4\end{cases}$，
与 $x^2>4$ 不同解，故选项 $A$ 错误；}
\step{不等式 $x^2+\dfrac1x>4+\dfrac1x$ 等价于 $\begin{cases}x^2>4\\ x\neq0\end{cases}$，
即 $x^2>4$，所以与 $x^2>4$ 同解，}
\step{故选项 $B$ 正确；}
\step{不等式 $x^2+\sqrt{x-2}>4+\sqrt{x-2}$ 等价于 $\begin{cases}x^2>4\\ x-2\geqslant0\end{cases}$，
与 $x^2>4$ 不同解，}
\step{故选项 $C$ 错误；}
\step{不等式 $\dfrac{1}{x^2-4}<0$ 等价于 $x^2<4$，与 $x^2>4$ 不同解，故选项 $D$ 错误.}
\step{故选 $B$.}
\solend

\fi

\item 若关于 $x$ 的二次不等式 $ax^2+bx+c\geqslant0(a\neq0)$ 的解集是 $\R$，那么\mbox{（\quad）}.

\keepwithprev
A. $a<0$，且 $b^2-4ac>0$\qquad B. $a<0$，且 $b^2-4ac\leqslant0$

% 注：原卷本题 D 选项与 A 选项文字完全重复（均作“$a<0$，且 $b^2-4ac>0$”），
%     照录未改，见文件头“原卷存疑”①。
C. $a>0$，且 $b^2-4ac\leqslant0$\qquad D. $a<0$，且 $b^2-4ac>0$

\ans{$C$}

\item 若不等式 $x^2+mx+1>2x+m$ 对满足 $|m|<2$ 的所有实数 $m$ 恒成立，
则实数 $x$ 的取值范围是\mbox{（\quad）}.

\keepwithprev
A. $-2<x<2$\qquad B. $x\geqslant3$\qquad C. $x\leqslant1$\qquad D. $x\leqslant-1$ 或 $x\geqslant3$

\ans{$D$}

\end{enumerate}

% “三、解答题”原本落在试卷版第 1 页页末、第 15 题翻到次页（切页审计 B 类）。
% 节标题自带的守卫只有 3\baselineskip，够放标题不够放标题＋首题，故在标题前再量一次页高。
\needvspace{4cm}
\bigsec{三、解答题}

\begin{enumerate}
\setcounter{enumi}{14}

\item 设关于 $x$ 的二次方程 $(m-1)x^2+(2m-4)x+m=0$ 有两个正根，求实数 $m$ 的取值范围.

\ifanswer
\solbegin
\step{$\begin{cases}\Delta=(2m-4)^2-4m(m-1)\geqslant0\\ -\dfrac{2m-4}{m-1}>0\\ \dfrac{m}{m-1}>0\end{cases}$，
解得 $1<m\leqslant\dfrac43$.}
\solend

\fi

\ansspace[6]

\item 已知关于 $x$ 的不等式 $2x^2-9x+a<0$ 的解集非空，对于其解集内的每一个 $x$ 的值，
至少能使不等式 $x^2-4x+3<0$ 或 $x^2-6x+8<0$ 中的一个成立，求实数 $a$ 的取值范围.

\ifanswer
\solbegin
\step{不等式 $x^2-4x+3<0$ 的解集为 $(1,3)$；$x^2-6x+8<0$ 的解集为 $(2,4)$.}
\step{关于 $x$ 的不等式 $2x^2-9x+a<0$ 的解集非空，所以 $\Delta=81-8a>0$，解得 $a<\dfrac{81}{8}$.}
\step{所以 $\dfrac{9-\sqrt{81-8a}}{4}<x<\dfrac{9+\sqrt{81-8a}}{4}$.}
\step{由于关于 $x$ 的不等式 $2x^2-9x+a<0$（解集非空）的每一个 $x$ 的值}
\step{至少满足不等式 $x^2-4x+3<0$ 和 $x^2-6x+8<0$ 中的一个，}
\step{所以 $\begin{cases}\dfrac{9-\sqrt{81-8a}}{4}\geqslant1\\ \dfrac{9+\sqrt{81-8a}}{4}\leqslant4\\ a<\dfrac{81}{8}\end{cases}$，
解得 $7\leqslant a<\dfrac{81}{8}$，所以实数 $a$ 的取值范围是 $\left[7,\dfrac{81}{8}\right)$.}
\solend

\fi

\ansspace[8]

\item 是否存在实数 $k$，使 $x=3$ 在不等式 $\dfrac{x+2}{k}>1+\dfrac{x^2-3}{k^2}$ 的解集中？
若存在，求 $k$ 的取值范围；若不存在，说明理由.

\ans{$(2,3)$}

\ansspace[6]

\item 关于 $x$ 的不等式 $x^2-(a+1)x+a<0$ 的解集为 $A$，$x^2-5x+4<0$ 的解集为 $B$，
若 $A\sub B$，求实数 $a$ 的取值范围.

\ans{$[1,4]$}

\ansspace[6]

\item 汽车在行驶中，由于惯性作用，刹车后还要继续向前滑行一段距离才能停住，我们称这段
距离为“刹车距离”. 刹车距离是分析事故的一个重要因素. 在一个限速 $40km/h$ 以内的弯道上，
甲、乙两辆汽车相向而行，发现情况不对，同时刹车，但还是相撞了，事后现场测得甲车的刹车
距离略超过 $12m$，乙车的刹车距离略超过 $10m$，又知甲、乙两种车型的刹车距离 $s(m)$ 与
车速 $x(km/h)$ 之间有如下关系：$s_{\text{甲}}=0.1x+0.01x^2$，$s_{\text{乙}}=0.05x+0.005x^2$. 问：
两车相撞的主要责任是谁？

\ifanswer
\solbegin
\step{由题意列出不等式组 $\begin{cases}0.1x+0.01x^2>12\\ 0.05x+0.005x^2>10\end{cases}$，
得 $\begin{cases}x<-40\text{ 或 }x>30\\ x<-50\text{ 或 }x>40\end{cases}$，}
\step{由于 $x>0$，从而得 $x_{\text{甲}}>30km/h$，$x_{\text{乙}}>40km/h$.}
\step{经比较得乙车超过限速，应负主要责任.}
\solend

\fi

\ansspace*[8]

\end{enumerate}

\end{document}
"
% 紧凑版：xelatex "\def\iscompact{1}% 高一40_大同中学_周末作业4.2.tex
%   —— 高一上数学“2026-2027 年大同中学高一上周末作业 4.2”（试卷版/答案版）
% 试卷版：xelatex 高一40_大同中学_周末作业4.2.tex
% 答案版：xelatex "\def\isanswer{1}% 高一40_大同中学_周末作业4.2.tex
%   —— 高一上数学“2026-2027 年大同中学高一上周末作业 4.2”（试卷版/答案版）
% 试卷版：xelatex 高一40_大同中学_周末作业4.2.tex
% 答案版：xelatex "\def\isanswer{1}\input{高一40_大同中学_周末作业4.2.tex}"
% 紧凑版：xelatex "\def\iscompact{1}\input{高一40_大同中学_周末作业4.2.tex}"
%
% 原始材料是微信公众号图文（公众号“魔都数学加油站”连载“27学年高一 40”，署名“破晓”，
% 2026-09-26 21:29 发布，原文标题“27学年高一40——2026-2027年上海市大同中学高一上周末作业4.2”，
% 原文链接 https://mp.weixin.qq.com/s/RrBSqDF_nJVbbx-bmmdtkQ ），正文为 4 张 PNG 页（4762×6735）。原文本身就是一份**讲评版**——题下直接印【答案】或【解析】，
% 第 11、13、14 题的答案还直接印在题干末尾的括号里。试题文字、答案与解析均照原卷录入，
% 未改内容；卷面的粉色斜水印（“魔都数学加油站”字样）属水印，未录入。
%
% 卷首标题：原卷页首印作“2026-2027 年大同中学高一上周末作业 4.2”（**卷面标题行即带校名**，
% 与高一16／高一18 等拍照卷须由本稿补校名的情形不同），本稿照卷面，年份区间按全稿惯例
% 排 en dash，前缀照原卷保留“年”。
%
% 篇幅与题号：本篇**只有 4 页**，卷面自“一、填空题”的**第 3 题**起编，终于“三、解答题”的
% **第 19 题**，共 17 题（填空 3—10、选择 11—14、解答 15—19）。第 1、2 题未在该文中发布——
% 这是该公众号连载讲评版的通例（高一02—高一09 与 高一36、高一38、高一39 等均自第 3 题起编，
% README“说明”已记），**不是截断**，故本稿题号亦自 3 起。它与 高一39（大同中学 周末作业4.1）
% **不构成同一份试卷的两半**：4.1 是另一次作业（11 页，题号**也自第 3 题**起、一直编到**第 24 题**，
% 以集合／常用逻辑用语为主、兼及不等式），两份各自独立起编（都从第 3 题开始），只是同为“大同中学高一上
% 周末作业 4”系列。
%
% 与原卷的排版差异（均已在 README“说明”中交代）：
%   ① 原文自**第 3 题**起（第 1、2 题未在该文中发布），故本稿题号亦自 3 起；
%   ② 原卷本就印有“一、填空题”“二、选择题”“三、解答题”三节标题，本稿照原卷分节，
%      只把标题改排成全稿统一的“大节标题”样式；
%   ③ 原卷为讲评版：第 4、5、6、7、8、17、18 题只印【答案】行，第 3、9、10、12、15、16、19 题
%      印【解析】（无单独的【答案】行），第 11、13、14 题的答案直接印在题干末尾的括号内；
%      本稿统一改为试卷版隐去、答案版以 \ans{}／\ifanswer 块给出，两版彻底分离（照 高一36 口径：
%      只印【解析】的题不另给 \ans{}，只印【答案】的题仅给 \ans{}）；
%   ④ 数集记号原卷作斜体的 $R$（第 5、10、13 题），本稿按全稿惯例统一写作 $\R$；
%   ⑤ 补集符号原卷作上横线（第 5 题的 $\overline{M}$），本稿照录，未改为 $\complement_\R$；
%   ⑥ 包含符号原卷两套并用：第 4 题作**无下划线**的 $\subset$（真包含），第 18 题作**带下划线**的
%      $\sub$（即 $\subseteq$），本稿照录、不归一；
%   ⑦ 原卷填空的横线约两个字宽，本稿按全稿惯例统一排 \blankline[4.4em]；
%   ⑧ 原卷未印副标题，本稿按**本卷实际内容**补出副标题“不等式”——全卷 17 题里约 15 题为
%      不等式（解一元二次／无理／分式不等式、恒成立与根的分布、含参不等式、应用题），
%      集合仅第 4、5 两题（第 18 题只是用集合关系表述解集），常用逻辑用语 0 题；
%   ⑨ 本卷**无插图**（全卷检索确认无“如图”）；
%   ⑩ 标点与单位：第 6 题的“$a,b$”本稿按全稿惯例排顿号“$a$、$b$”，第 10 题的三例
%      “$[\pi]=3,[-1.2]=-2,[\tfrac12]=0$”之间排顿号，第 5 题原卷“全集 $U=R$ 集合 $M$”之间
%      漏排标点、本稿补逗号；车速单位原卷印作斜体 $km/h$（第 19 题），本稿照录；半角括号、逗号、
%      问号一律排全角（第 17、19 题原卷的两个问号本就作全角，本稿照录）。
%
% 原卷存疑两处（照抄原卷、未改，见源码中该处上方注释）：
%   ① 第 13 题的 D 选项与 A 选项文字完全重复（均作“$a<0$，且 $b^2-4ac>0$”），原卷如此。
%   ② 第 10 题题干末作“…的取值范围____”，**漏排“是”**（本卷其余填空均作“…是____”）——
%      本稿照录、未补字，见源码该题上方 `%` 注。

\documentclass[a4paper,11pt]{article}
\usepackage{common}

\begin{document}

\examtitlest[3]{2026--2027 年大同中学高一上周末作业 4.2}{不等式}

\bigsec{一、填空题}

\begin{enumerate}
\setcounter{enumi}{2}

\item 在平面直角坐标系中，若点 $(2,b)$ 到原点的距离不小于 $5$，则 $b$ 的取值范围是 \blankline[4.4em].

\ifanswer
\solbegin
\step{因为点 $(2,b)$ 到原点的距离不小于 $5$，所以 $\sqrt{(2-0)^2+(b-0)^2}\geqslant5$，}
\step{所以 $4+b^2\geqslant25$，得 $b\geqslant\sqrt{21}$ 或 $b\leqslant-\sqrt{21}$，}
\step{则 $b$ 的取值范围是 $(-\infty,-\sqrt{21}]\cup[\sqrt{21},+\infty)$.}
\solend

\fi

\item 若集合 $A=\{x\mid|x|>3\}$，集合 $B=\{x\mid x\geqslant a\}$，且 $B\subset A$，
则实数 $a$ 的取值范围是 \blankline[4.4em].

\ans{$(3,+\infty)$}

% 注：原卷“全集 $U=R$”与“集合 $M$”之间**漏排标点**（只有空格），本稿补逗号——
%     照 高一b 第 13 题“原卷漏排、本稿补上”的先例（高一10、高一24 同例），见文件头⑩。
\item 若全集 $U=\R$，集合 $M=\{x\mid y=\sqrt{9-x^2}\}$，$S=\{x\mid x^2-3x+2\geqslant0\}$，
则 $\overline{M}\cup S=$ \blankline[4.4em].

\ans{$(-\infty,1]\cup[2,+\infty)$}

\item 设 $a$、$b$ 为实数，若不等式 $x^2-ax-b<0$ 的解集为 $\{x\mid2<x<8\}$，
则 $bx^2-ax-1>0$ 的解集为 \blankline[4.4em].

\ans{$\left(-\dfrac12,-\dfrac18\right)$}

\item 不等式 $4x-13+\sqrt{x+5}\geqslant8x-1+\sqrt{x+5}$ 的解集是 \blankline[4.4em].

\ans{$[-5,-3]$}

\item 若关于 $x$ 的不等式 $ax^2-2ax+3\leqslant0$ 的解集为 $\varnothing$，
则实数 $a$ 的取值范围是 \blankline[4.4em].

\ans{$[0,3)$}

\item 若不等式 $x^2-ax-2<0$ 在 $x\in(1,2)$ 内恒成立，则实数 $a$ 的取值范围是 \blankline[4.4em].

\ifanswer
\solbegin
\step{由不等式 $x^2-ax-2<0$ 在 $x\in(1,2)$ 内恒成立，}
\step{得 $ax>x^2-2$，即 $a>x-\dfrac2x$ 在 $x\in(1,2)$ 内恒成立，}
\step{令 $t(x)=x-\dfrac2x$，$x\in(1,2)$，该函数为严格增函数，则 $t(x)<t(2)=1$，得 $a\geqslant1$，}
\step{所以 $a$ 的取值范围是 $[1,+\infty)$.}
\solend

\fi

% 注：原卷此处作“的取值范围\_\_\_\_”，**漏排“是”**（本卷其余填空均作“…是\_\_\_\_”），
%     本稿照录、未补字，见文件头“原卷存疑”②。
\item 设 $x\in\R$，$[x]$ 表示不大于 $x$ 的最大整数，如 $[\pi]=3$、$[-1.2]=-2$、$\left[\dfrac12\right]=0$，
则使 $[|x-1|]=3$ 成立的 $x$ 的取值范围 \blankline[4.4em].

\ifanswer
\solbegin
\step{由题意 $[|x-1|]=3$，则 $3\leqslant|x-1|<4$，所以 $3\leqslant x-1<4$ 或 $-4\leqslant x-1<-3$，}
\step{解得 $4\leqslant x<5$ 或 $-3<x\leqslant-2$，}
\step{所以使 $[|x-1|]=3$ 成立的 $x$ 的取值范围是 $(-3,-2]\cup[4,5)$.}
\solend

\fi

\end{enumerate}

\bigsec{二、选择题}

\begin{enumerate}
\setcounter{enumi}{10}

\item 若关于 $x$ 的方程 $k(x+2)-3(k-1)=x+1$ 有负数解，则实数 $k$ 的取值范围是\mbox{（\quad）}.

\keepwithprev
A. $(1,2)$\qquad B. $(-\infty,1)\cup(2,+\infty)$\qquad C. $(-\infty,2)$\qquad D. $(-\infty,1)$

\ans{$A$}

\item 与 $x^2>4$ 同解的不等式是\mbox{（\quad）}.

\keepwithprev
A. $x^2+\dfrac{1}{x-4}>4+\dfrac{1}{x-4}$\qquad B. $x^2+\dfrac1x>4+\dfrac1x$

C. $x^2+\sqrt{x-2}>4+\sqrt{x-2}$\qquad D. $\dfrac{1}{x^2-4}<0$

\ifanswer
\solbegin
\step{不等式 $x^2+\dfrac{1}{x-4}>4+\dfrac{1}{x-4}$ 等价于 $\begin{cases}x^2>4\\ x\neq4\end{cases}$，
与 $x^2>4$ 不同解，故选项 $A$ 错误；}
\step{不等式 $x^2+\dfrac1x>4+\dfrac1x$ 等价于 $\begin{cases}x^2>4\\ x\neq0\end{cases}$，
即 $x^2>4$，所以与 $x^2>4$ 同解，}
\step{故选项 $B$ 正确；}
\step{不等式 $x^2+\sqrt{x-2}>4+\sqrt{x-2}$ 等价于 $\begin{cases}x^2>4\\ x-2\geqslant0\end{cases}$，
与 $x^2>4$ 不同解，}
\step{故选项 $C$ 错误；}
\step{不等式 $\dfrac{1}{x^2-4}<0$ 等价于 $x^2<4$，与 $x^2>4$ 不同解，故选项 $D$ 错误.}
\step{故选 $B$.}
\solend

\fi

\item 若关于 $x$ 的二次不等式 $ax^2+bx+c\geqslant0(a\neq0)$ 的解集是 $\R$，那么\mbox{（\quad）}.

\keepwithprev
A. $a<0$，且 $b^2-4ac>0$\qquad B. $a<0$，且 $b^2-4ac\leqslant0$

% 注：原卷本题 D 选项与 A 选项文字完全重复（均作“$a<0$，且 $b^2-4ac>0$”），
%     照录未改，见文件头“原卷存疑”①。
C. $a>0$，且 $b^2-4ac\leqslant0$\qquad D. $a<0$，且 $b^2-4ac>0$

\ans{$C$}

\item 若不等式 $x^2+mx+1>2x+m$ 对满足 $|m|<2$ 的所有实数 $m$ 恒成立，
则实数 $x$ 的取值范围是\mbox{（\quad）}.

\keepwithprev
A. $-2<x<2$\qquad B. $x\geqslant3$\qquad C. $x\leqslant1$\qquad D. $x\leqslant-1$ 或 $x\geqslant3$

\ans{$D$}

\end{enumerate}

% “三、解答题”原本落在试卷版第 1 页页末、第 15 题翻到次页（切页审计 B 类）。
% 节标题自带的守卫只有 3\baselineskip，够放标题不够放标题＋首题，故在标题前再量一次页高。
\needvspace{4cm}
\bigsec{三、解答题}

\begin{enumerate}
\setcounter{enumi}{14}

\item 设关于 $x$ 的二次方程 $(m-1)x^2+(2m-4)x+m=0$ 有两个正根，求实数 $m$ 的取值范围.

\ifanswer
\solbegin
\step{$\begin{cases}\Delta=(2m-4)^2-4m(m-1)\geqslant0\\ -\dfrac{2m-4}{m-1}>0\\ \dfrac{m}{m-1}>0\end{cases}$，
解得 $1<m\leqslant\dfrac43$.}
\solend

\fi

\ansspace[6]

\item 已知关于 $x$ 的不等式 $2x^2-9x+a<0$ 的解集非空，对于其解集内的每一个 $x$ 的值，
至少能使不等式 $x^2-4x+3<0$ 或 $x^2-6x+8<0$ 中的一个成立，求实数 $a$ 的取值范围.

\ifanswer
\solbegin
\step{不等式 $x^2-4x+3<0$ 的解集为 $(1,3)$；$x^2-6x+8<0$ 的解集为 $(2,4)$.}
\step{关于 $x$ 的不等式 $2x^2-9x+a<0$ 的解集非空，所以 $\Delta=81-8a>0$，解得 $a<\dfrac{81}{8}$.}
\step{所以 $\dfrac{9-\sqrt{81-8a}}{4}<x<\dfrac{9+\sqrt{81-8a}}{4}$.}
\step{由于关于 $x$ 的不等式 $2x^2-9x+a<0$（解集非空）的每一个 $x$ 的值}
\step{至少满足不等式 $x^2-4x+3<0$ 和 $x^2-6x+8<0$ 中的一个，}
\step{所以 $\begin{cases}\dfrac{9-\sqrt{81-8a}}{4}\geqslant1\\ \dfrac{9+\sqrt{81-8a}}{4}\leqslant4\\ a<\dfrac{81}{8}\end{cases}$，
解得 $7\leqslant a<\dfrac{81}{8}$，所以实数 $a$ 的取值范围是 $\left[7,\dfrac{81}{8}\right)$.}
\solend

\fi

\ansspace[8]

\item 是否存在实数 $k$，使 $x=3$ 在不等式 $\dfrac{x+2}{k}>1+\dfrac{x^2-3}{k^2}$ 的解集中？
若存在，求 $k$ 的取值范围；若不存在，说明理由.

\ans{$(2,3)$}

\ansspace[6]

\item 关于 $x$ 的不等式 $x^2-(a+1)x+a<0$ 的解集为 $A$，$x^2-5x+4<0$ 的解集为 $B$，
若 $A\sub B$，求实数 $a$ 的取值范围.

\ans{$[1,4]$}

\ansspace[6]

\item 汽车在行驶中，由于惯性作用，刹车后还要继续向前滑行一段距离才能停住，我们称这段
距离为“刹车距离”. 刹车距离是分析事故的一个重要因素. 在一个限速 $40km/h$ 以内的弯道上，
甲、乙两辆汽车相向而行，发现情况不对，同时刹车，但还是相撞了，事后现场测得甲车的刹车
距离略超过 $12m$，乙车的刹车距离略超过 $10m$，又知甲、乙两种车型的刹车距离 $s(m)$ 与
车速 $x(km/h)$ 之间有如下关系：$s_{\text{甲}}=0.1x+0.01x^2$，$s_{\text{乙}}=0.05x+0.005x^2$. 问：
两车相撞的主要责任是谁？

\ifanswer
\solbegin
\step{由题意列出不等式组 $\begin{cases}0.1x+0.01x^2>12\\ 0.05x+0.005x^2>10\end{cases}$，
得 $\begin{cases}x<-40\text{ 或 }x>30\\ x<-50\text{ 或 }x>40\end{cases}$，}
\step{由于 $x>0$，从而得 $x_{\text{甲}}>30km/h$，$x_{\text{乙}}>40km/h$.}
\step{经比较得乙车超过限速，应负主要责任.}
\solend

\fi

\ansspace*[8]

\end{enumerate}

\end{document}
"
% 紧凑版：xelatex "\def\iscompact{1}% 高一40_大同中学_周末作业4.2.tex
%   —— 高一上数学“2026-2027 年大同中学高一上周末作业 4.2”（试卷版/答案版）
% 试卷版：xelatex 高一40_大同中学_周末作业4.2.tex
% 答案版：xelatex "\def\isanswer{1}\input{高一40_大同中学_周末作业4.2.tex}"
% 紧凑版：xelatex "\def\iscompact{1}\input{高一40_大同中学_周末作业4.2.tex}"
%
% 原始材料是微信公众号图文（公众号“魔都数学加油站”连载“27学年高一 40”，署名“破晓”，
% 2026-09-26 21:29 发布，原文标题“27学年高一40——2026-2027年上海市大同中学高一上周末作业4.2”，
% 原文链接 https://mp.weixin.qq.com/s/RrBSqDF_nJVbbx-bmmdtkQ ），正文为 4 张 PNG 页（4762×6735）。原文本身就是一份**讲评版**——题下直接印【答案】或【解析】，
% 第 11、13、14 题的答案还直接印在题干末尾的括号里。试题文字、答案与解析均照原卷录入，
% 未改内容；卷面的粉色斜水印（“魔都数学加油站”字样）属水印，未录入。
%
% 卷首标题：原卷页首印作“2026-2027 年大同中学高一上周末作业 4.2”（**卷面标题行即带校名**，
% 与高一16／高一18 等拍照卷须由本稿补校名的情形不同），本稿照卷面，年份区间按全稿惯例
% 排 en dash，前缀照原卷保留“年”。
%
% 篇幅与题号：本篇**只有 4 页**，卷面自“一、填空题”的**第 3 题**起编，终于“三、解答题”的
% **第 19 题**，共 17 题（填空 3—10、选择 11—14、解答 15—19）。第 1、2 题未在该文中发布——
% 这是该公众号连载讲评版的通例（高一02—高一09 与 高一36、高一38、高一39 等均自第 3 题起编，
% README“说明”已记），**不是截断**，故本稿题号亦自 3 起。它与 高一39（大同中学 周末作业4.1）
% **不构成同一份试卷的两半**：4.1 是另一次作业（11 页，题号**也自第 3 题**起、一直编到**第 24 题**，
% 以集合／常用逻辑用语为主、兼及不等式），两份各自独立起编（都从第 3 题开始），只是同为“大同中学高一上
% 周末作业 4”系列。
%
% 与原卷的排版差异（均已在 README“说明”中交代）：
%   ① 原文自**第 3 题**起（第 1、2 题未在该文中发布），故本稿题号亦自 3 起；
%   ② 原卷本就印有“一、填空题”“二、选择题”“三、解答题”三节标题，本稿照原卷分节，
%      只把标题改排成全稿统一的“大节标题”样式；
%   ③ 原卷为讲评版：第 4、5、6、7、8、17、18 题只印【答案】行，第 3、9、10、12、15、16、19 题
%      印【解析】（无单独的【答案】行），第 11、13、14 题的答案直接印在题干末尾的括号内；
%      本稿统一改为试卷版隐去、答案版以 \ans{}／\ifanswer 块给出，两版彻底分离（照 高一36 口径：
%      只印【解析】的题不另给 \ans{}，只印【答案】的题仅给 \ans{}）；
%   ④ 数集记号原卷作斜体的 $R$（第 5、10、13 题），本稿按全稿惯例统一写作 $\R$；
%   ⑤ 补集符号原卷作上横线（第 5 题的 $\overline{M}$），本稿照录，未改为 $\complement_\R$；
%   ⑥ 包含符号原卷两套并用：第 4 题作**无下划线**的 $\subset$（真包含），第 18 题作**带下划线**的
%      $\sub$（即 $\subseteq$），本稿照录、不归一；
%   ⑦ 原卷填空的横线约两个字宽，本稿按全稿惯例统一排 \blankline[4.4em]；
%   ⑧ 原卷未印副标题，本稿按**本卷实际内容**补出副标题“不等式”——全卷 17 题里约 15 题为
%      不等式（解一元二次／无理／分式不等式、恒成立与根的分布、含参不等式、应用题），
%      集合仅第 4、5 两题（第 18 题只是用集合关系表述解集），常用逻辑用语 0 题；
%   ⑨ 本卷**无插图**（全卷检索确认无“如图”）；
%   ⑩ 标点与单位：第 6 题的“$a,b$”本稿按全稿惯例排顿号“$a$、$b$”，第 10 题的三例
%      “$[\pi]=3,[-1.2]=-2,[\tfrac12]=0$”之间排顿号，第 5 题原卷“全集 $U=R$ 集合 $M$”之间
%      漏排标点、本稿补逗号；车速单位原卷印作斜体 $km/h$（第 19 题），本稿照录；半角括号、逗号、
%      问号一律排全角（第 17、19 题原卷的两个问号本就作全角，本稿照录）。
%
% 原卷存疑两处（照抄原卷、未改，见源码中该处上方注释）：
%   ① 第 13 题的 D 选项与 A 选项文字完全重复（均作“$a<0$，且 $b^2-4ac>0$”），原卷如此。
%   ② 第 10 题题干末作“…的取值范围____”，**漏排“是”**（本卷其余填空均作“…是____”）——
%      本稿照录、未补字，见源码该题上方 `%` 注。

\documentclass[a4paper,11pt]{article}
\usepackage{common}

\begin{document}

\examtitlest[3]{2026--2027 年大同中学高一上周末作业 4.2}{不等式}

\bigsec{一、填空题}

\begin{enumerate}
\setcounter{enumi}{2}

\item 在平面直角坐标系中，若点 $(2,b)$ 到原点的距离不小于 $5$，则 $b$ 的取值范围是 \blankline[4.4em].

\ifanswer
\solbegin
\step{因为点 $(2,b)$ 到原点的距离不小于 $5$，所以 $\sqrt{(2-0)^2+(b-0)^2}\geqslant5$，}
\step{所以 $4+b^2\geqslant25$，得 $b\geqslant\sqrt{21}$ 或 $b\leqslant-\sqrt{21}$，}
\step{则 $b$ 的取值范围是 $(-\infty,-\sqrt{21}]\cup[\sqrt{21},+\infty)$.}
\solend

\fi

\item 若集合 $A=\{x\mid|x|>3\}$，集合 $B=\{x\mid x\geqslant a\}$，且 $B\subset A$，
则实数 $a$ 的取值范围是 \blankline[4.4em].

\ans{$(3,+\infty)$}

% 注：原卷“全集 $U=R$”与“集合 $M$”之间**漏排标点**（只有空格），本稿补逗号——
%     照 高一b 第 13 题“原卷漏排、本稿补上”的先例（高一10、高一24 同例），见文件头⑩。
\item 若全集 $U=\R$，集合 $M=\{x\mid y=\sqrt{9-x^2}\}$，$S=\{x\mid x^2-3x+2\geqslant0\}$，
则 $\overline{M}\cup S=$ \blankline[4.4em].

\ans{$(-\infty,1]\cup[2,+\infty)$}

\item 设 $a$、$b$ 为实数，若不等式 $x^2-ax-b<0$ 的解集为 $\{x\mid2<x<8\}$，
则 $bx^2-ax-1>0$ 的解集为 \blankline[4.4em].

\ans{$\left(-\dfrac12,-\dfrac18\right)$}

\item 不等式 $4x-13+\sqrt{x+5}\geqslant8x-1+\sqrt{x+5}$ 的解集是 \blankline[4.4em].

\ans{$[-5,-3]$}

\item 若关于 $x$ 的不等式 $ax^2-2ax+3\leqslant0$ 的解集为 $\varnothing$，
则实数 $a$ 的取值范围是 \blankline[4.4em].

\ans{$[0,3)$}

\item 若不等式 $x^2-ax-2<0$ 在 $x\in(1,2)$ 内恒成立，则实数 $a$ 的取值范围是 \blankline[4.4em].

\ifanswer
\solbegin
\step{由不等式 $x^2-ax-2<0$ 在 $x\in(1,2)$ 内恒成立，}
\step{得 $ax>x^2-2$，即 $a>x-\dfrac2x$ 在 $x\in(1,2)$ 内恒成立，}
\step{令 $t(x)=x-\dfrac2x$，$x\in(1,2)$，该函数为严格增函数，则 $t(x)<t(2)=1$，得 $a\geqslant1$，}
\step{所以 $a$ 的取值范围是 $[1,+\infty)$.}
\solend

\fi

% 注：原卷此处作“的取值范围\_\_\_\_”，**漏排“是”**（本卷其余填空均作“…是\_\_\_\_”），
%     本稿照录、未补字，见文件头“原卷存疑”②。
\item 设 $x\in\R$，$[x]$ 表示不大于 $x$ 的最大整数，如 $[\pi]=3$、$[-1.2]=-2$、$\left[\dfrac12\right]=0$，
则使 $[|x-1|]=3$ 成立的 $x$ 的取值范围 \blankline[4.4em].

\ifanswer
\solbegin
\step{由题意 $[|x-1|]=3$，则 $3\leqslant|x-1|<4$，所以 $3\leqslant x-1<4$ 或 $-4\leqslant x-1<-3$，}
\step{解得 $4\leqslant x<5$ 或 $-3<x\leqslant-2$，}
\step{所以使 $[|x-1|]=3$ 成立的 $x$ 的取值范围是 $(-3,-2]\cup[4,5)$.}
\solend

\fi

\end{enumerate}

\bigsec{二、选择题}

\begin{enumerate}
\setcounter{enumi}{10}

\item 若关于 $x$ 的方程 $k(x+2)-3(k-1)=x+1$ 有负数解，则实数 $k$ 的取值范围是\mbox{（\quad）}.

\keepwithprev
A. $(1,2)$\qquad B. $(-\infty,1)\cup(2,+\infty)$\qquad C. $(-\infty,2)$\qquad D. $(-\infty,1)$

\ans{$A$}

\item 与 $x^2>4$ 同解的不等式是\mbox{（\quad）}.

\keepwithprev
A. $x^2+\dfrac{1}{x-4}>4+\dfrac{1}{x-4}$\qquad B. $x^2+\dfrac1x>4+\dfrac1x$

C. $x^2+\sqrt{x-2}>4+\sqrt{x-2}$\qquad D. $\dfrac{1}{x^2-4}<0$

\ifanswer
\solbegin
\step{不等式 $x^2+\dfrac{1}{x-4}>4+\dfrac{1}{x-4}$ 等价于 $\begin{cases}x^2>4\\ x\neq4\end{cases}$，
与 $x^2>4$ 不同解，故选项 $A$ 错误；}
\step{不等式 $x^2+\dfrac1x>4+\dfrac1x$ 等价于 $\begin{cases}x^2>4\\ x\neq0\end{cases}$，
即 $x^2>4$，所以与 $x^2>4$ 同解，}
\step{故选项 $B$ 正确；}
\step{不等式 $x^2+\sqrt{x-2}>4+\sqrt{x-2}$ 等价于 $\begin{cases}x^2>4\\ x-2\geqslant0\end{cases}$，
与 $x^2>4$ 不同解，}
\step{故选项 $C$ 错误；}
\step{不等式 $\dfrac{1}{x^2-4}<0$ 等价于 $x^2<4$，与 $x^2>4$ 不同解，故选项 $D$ 错误.}
\step{故选 $B$.}
\solend

\fi

\item 若关于 $x$ 的二次不等式 $ax^2+bx+c\geqslant0(a\neq0)$ 的解集是 $\R$，那么\mbox{（\quad）}.

\keepwithprev
A. $a<0$，且 $b^2-4ac>0$\qquad B. $a<0$，且 $b^2-4ac\leqslant0$

% 注：原卷本题 D 选项与 A 选项文字完全重复（均作“$a<0$，且 $b^2-4ac>0$”），
%     照录未改，见文件头“原卷存疑”①。
C. $a>0$，且 $b^2-4ac\leqslant0$\qquad D. $a<0$，且 $b^2-4ac>0$

\ans{$C$}

\item 若不等式 $x^2+mx+1>2x+m$ 对满足 $|m|<2$ 的所有实数 $m$ 恒成立，
则实数 $x$ 的取值范围是\mbox{（\quad）}.

\keepwithprev
A. $-2<x<2$\qquad B. $x\geqslant3$\qquad C. $x\leqslant1$\qquad D. $x\leqslant-1$ 或 $x\geqslant3$

\ans{$D$}

\end{enumerate}

% “三、解答题”原本落在试卷版第 1 页页末、第 15 题翻到次页（切页审计 B 类）。
% 节标题自带的守卫只有 3\baselineskip，够放标题不够放标题＋首题，故在标题前再量一次页高。
\needvspace{4cm}
\bigsec{三、解答题}

\begin{enumerate}
\setcounter{enumi}{14}

\item 设关于 $x$ 的二次方程 $(m-1)x^2+(2m-4)x+m=0$ 有两个正根，求实数 $m$ 的取值范围.

\ifanswer
\solbegin
\step{$\begin{cases}\Delta=(2m-4)^2-4m(m-1)\geqslant0\\ -\dfrac{2m-4}{m-1}>0\\ \dfrac{m}{m-1}>0\end{cases}$，
解得 $1<m\leqslant\dfrac43$.}
\solend

\fi

\ansspace[6]

\item 已知关于 $x$ 的不等式 $2x^2-9x+a<0$ 的解集非空，对于其解集内的每一个 $x$ 的值，
至少能使不等式 $x^2-4x+3<0$ 或 $x^2-6x+8<0$ 中的一个成立，求实数 $a$ 的取值范围.

\ifanswer
\solbegin
\step{不等式 $x^2-4x+3<0$ 的解集为 $(1,3)$；$x^2-6x+8<0$ 的解集为 $(2,4)$.}
\step{关于 $x$ 的不等式 $2x^2-9x+a<0$ 的解集非空，所以 $\Delta=81-8a>0$，解得 $a<\dfrac{81}{8}$.}
\step{所以 $\dfrac{9-\sqrt{81-8a}}{4}<x<\dfrac{9+\sqrt{81-8a}}{4}$.}
\step{由于关于 $x$ 的不等式 $2x^2-9x+a<0$（解集非空）的每一个 $x$ 的值}
\step{至少满足不等式 $x^2-4x+3<0$ 和 $x^2-6x+8<0$ 中的一个，}
\step{所以 $\begin{cases}\dfrac{9-\sqrt{81-8a}}{4}\geqslant1\\ \dfrac{9+\sqrt{81-8a}}{4}\leqslant4\\ a<\dfrac{81}{8}\end{cases}$，
解得 $7\leqslant a<\dfrac{81}{8}$，所以实数 $a$ 的取值范围是 $\left[7,\dfrac{81}{8}\right)$.}
\solend

\fi

\ansspace[8]

\item 是否存在实数 $k$，使 $x=3$ 在不等式 $\dfrac{x+2}{k}>1+\dfrac{x^2-3}{k^2}$ 的解集中？
若存在，求 $k$ 的取值范围；若不存在，说明理由.

\ans{$(2,3)$}

\ansspace[6]

\item 关于 $x$ 的不等式 $x^2-(a+1)x+a<0$ 的解集为 $A$，$x^2-5x+4<0$ 的解集为 $B$，
若 $A\sub B$，求实数 $a$ 的取值范围.

\ans{$[1,4]$}

\ansspace[6]

\item 汽车在行驶中，由于惯性作用，刹车后还要继续向前滑行一段距离才能停住，我们称这段
距离为“刹车距离”. 刹车距离是分析事故的一个重要因素. 在一个限速 $40km/h$ 以内的弯道上，
甲、乙两辆汽车相向而行，发现情况不对，同时刹车，但还是相撞了，事后现场测得甲车的刹车
距离略超过 $12m$，乙车的刹车距离略超过 $10m$，又知甲、乙两种车型的刹车距离 $s(m)$ 与
车速 $x(km/h)$ 之间有如下关系：$s_{\text{甲}}=0.1x+0.01x^2$，$s_{\text{乙}}=0.05x+0.005x^2$. 问：
两车相撞的主要责任是谁？

\ifanswer
\solbegin
\step{由题意列出不等式组 $\begin{cases}0.1x+0.01x^2>12\\ 0.05x+0.005x^2>10\end{cases}$，
得 $\begin{cases}x<-40\text{ 或 }x>30\\ x<-50\text{ 或 }x>40\end{cases}$，}
\step{由于 $x>0$，从而得 $x_{\text{甲}}>30km/h$，$x_{\text{乙}}>40km/h$.}
\step{经比较得乙车超过限速，应负主要责任.}
\solend

\fi

\ansspace*[8]

\end{enumerate}

\end{document}
"
%
% 原始材料是微信公众号图文（公众号“魔都数学加油站”连载“27学年高一 40”，署名“破晓”，
% 2026-09-26 21:29 发布，原文标题“27学年高一40——2026-2027年上海市大同中学高一上周末作业4.2”，
% 原文链接 https://mp.weixin.qq.com/s/RrBSqDF_nJVbbx-bmmdtkQ ），正文为 4 张 PNG 页（4762×6735）。原文本身就是一份**讲评版**——题下直接印【答案】或【解析】，
% 第 11、13、14 题的答案还直接印在题干末尾的括号里。试题文字、答案与解析均照原卷录入，
% 未改内容；卷面的粉色斜水印（“魔都数学加油站”字样）属水印，未录入。
%
% 卷首标题：原卷页首印作“2026-2027 年大同中学高一上周末作业 4.2”（**卷面标题行即带校名**，
% 与高一16／高一18 等拍照卷须由本稿补校名的情形不同），本稿照卷面，年份区间按全稿惯例
% 排 en dash，前缀照原卷保留“年”。
%
% 篇幅与题号：本篇**只有 4 页**，卷面自“一、填空题”的**第 3 题**起编，终于“三、解答题”的
% **第 19 题**，共 17 题（填空 3—10、选择 11—14、解答 15—19）。第 1、2 题未在该文中发布——
% 这是该公众号连载讲评版的通例（高一02—高一09 与 高一36、高一38、高一39 等均自第 3 题起编，
% README“说明”已记），**不是截断**，故本稿题号亦自 3 起。它与 高一39（大同中学 周末作业4.1）
% **不构成同一份试卷的两半**：4.1 是另一次作业（11 页，题号**也自第 3 题**起、一直编到**第 24 题**，
% 以集合／常用逻辑用语为主、兼及不等式），两份各自独立起编（都从第 3 题开始），只是同为“大同中学高一上
% 周末作业 4”系列。
%
% 与原卷的排版差异（均已在 README“说明”中交代）：
%   ① 原文自**第 3 题**起（第 1、2 题未在该文中发布），故本稿题号亦自 3 起；
%   ② 原卷本就印有“一、填空题”“二、选择题”“三、解答题”三节标题，本稿照原卷分节，
%      只把标题改排成全稿统一的“大节标题”样式；
%   ③ 原卷为讲评版：第 4、5、6、7、8、17、18 题只印【答案】行，第 3、9、10、12、15、16、19 题
%      印【解析】（无单独的【答案】行），第 11、13、14 题的答案直接印在题干末尾的括号内；
%      本稿统一改为试卷版隐去、答案版以 \ans{}／\ifanswer 块给出，两版彻底分离（照 高一36 口径：
%      只印【解析】的题不另给 \ans{}，只印【答案】的题仅给 \ans{}）；
%   ④ 数集记号原卷作斜体的 $R$（第 5、10、13 题），本稿按全稿惯例统一写作 $\R$；
%   ⑤ 补集符号原卷作上横线（第 5 题的 $\overline{M}$），本稿照录，未改为 $\complement_\R$；
%   ⑥ 包含符号原卷两套并用：第 4 题作**无下划线**的 $\subset$（真包含），第 18 题作**带下划线**的
%      $\sub$（即 $\subseteq$），本稿照录、不归一；
%   ⑦ 原卷填空的横线约两个字宽，本稿按全稿惯例统一排 \blankline[4.4em]；
%   ⑧ 原卷未印副标题，本稿按**本卷实际内容**补出副标题“不等式”——全卷 17 题里约 15 题为
%      不等式（解一元二次／无理／分式不等式、恒成立与根的分布、含参不等式、应用题），
%      集合仅第 4、5 两题（第 18 题只是用集合关系表述解集），常用逻辑用语 0 题；
%   ⑨ 本卷**无插图**（全卷检索确认无“如图”）；
%   ⑩ 标点与单位：第 6 题的“$a,b$”本稿按全稿惯例排顿号“$a$、$b$”，第 10 题的三例
%      “$[\pi]=3,[-1.2]=-2,[\tfrac12]=0$”之间排顿号，第 5 题原卷“全集 $U=R$ 集合 $M$”之间
%      漏排标点、本稿补逗号；车速单位原卷印作斜体 $km/h$（第 19 题），本稿照录；半角括号、逗号、
%      问号一律排全角（第 17、19 题原卷的两个问号本就作全角，本稿照录）。
%
% 原卷存疑两处（照抄原卷、未改，见源码中该处上方注释）：
%   ① 第 13 题的 D 选项与 A 选项文字完全重复（均作“$a<0$，且 $b^2-4ac>0$”），原卷如此。
%   ② 第 10 题题干末作“…的取值范围____”，**漏排“是”**（本卷其余填空均作“…是____”）——
%      本稿照录、未补字，见源码该题上方 `%` 注。

\documentclass[a4paper,11pt]{article}
\usepackage{common}

\begin{document}

\examtitlest[3]{2026--2027 年大同中学高一上周末作业 4.2}{不等式}

\bigsec{一、填空题}

\begin{enumerate}
\setcounter{enumi}{2}

\item 在平面直角坐标系中，若点 $(2,b)$ 到原点的距离不小于 $5$，则 $b$ 的取值范围是 \blankline[4.4em].

\ifanswer
\solbegin
\step{因为点 $(2,b)$ 到原点的距离不小于 $5$，所以 $\sqrt{(2-0)^2+(b-0)^2}\geqslant5$，}
\step{所以 $4+b^2\geqslant25$，得 $b\geqslant\sqrt{21}$ 或 $b\leqslant-\sqrt{21}$，}
\step{则 $b$ 的取值范围是 $(-\infty,-\sqrt{21}]\cup[\sqrt{21},+\infty)$.}
\solend

\fi

\item 若集合 $A=\{x\mid|x|>3\}$，集合 $B=\{x\mid x\geqslant a\}$，且 $B\subset A$，
则实数 $a$ 的取值范围是 \blankline[4.4em].

\ans{$(3,+\infty)$}

% 注：原卷“全集 $U=R$”与“集合 $M$”之间**漏排标点**（只有空格），本稿补逗号——
%     照 高一b 第 13 题“原卷漏排、本稿补上”的先例（高一10、高一24 同例），见文件头⑩。
\item 若全集 $U=\R$，集合 $M=\{x\mid y=\sqrt{9-x^2}\}$，$S=\{x\mid x^2-3x+2\geqslant0\}$，
则 $\overline{M}\cup S=$ \blankline[4.4em].

\ans{$(-\infty,1]\cup[2,+\infty)$}

\item 设 $a$、$b$ 为实数，若不等式 $x^2-ax-b<0$ 的解集为 $\{x\mid2<x<8\}$，
则 $bx^2-ax-1>0$ 的解集为 \blankline[4.4em].

\ans{$\left(-\dfrac12,-\dfrac18\right)$}

\item 不等式 $4x-13+\sqrt{x+5}\geqslant8x-1+\sqrt{x+5}$ 的解集是 \blankline[4.4em].

\ans{$[-5,-3]$}

\item 若关于 $x$ 的不等式 $ax^2-2ax+3\leqslant0$ 的解集为 $\varnothing$，
则实数 $a$ 的取值范围是 \blankline[4.4em].

\ans{$[0,3)$}

\item 若不等式 $x^2-ax-2<0$ 在 $x\in(1,2)$ 内恒成立，则实数 $a$ 的取值范围是 \blankline[4.4em].

\ifanswer
\solbegin
\step{由不等式 $x^2-ax-2<0$ 在 $x\in(1,2)$ 内恒成立，}
\step{得 $ax>x^2-2$，即 $a>x-\dfrac2x$ 在 $x\in(1,2)$ 内恒成立，}
\step{令 $t(x)=x-\dfrac2x$，$x\in(1,2)$，该函数为严格增函数，则 $t(x)<t(2)=1$，得 $a\geqslant1$，}
\step{所以 $a$ 的取值范围是 $[1,+\infty)$.}
\solend

\fi

% 注：原卷此处作“的取值范围\_\_\_\_”，**漏排“是”**（本卷其余填空均作“…是\_\_\_\_”），
%     本稿照录、未补字，见文件头“原卷存疑”②。
\item 设 $x\in\R$，$[x]$ 表示不大于 $x$ 的最大整数，如 $[\pi]=3$、$[-1.2]=-2$、$\left[\dfrac12\right]=0$，
则使 $[|x-1|]=3$ 成立的 $x$ 的取值范围 \blankline[4.4em].

\ifanswer
\solbegin
\step{由题意 $[|x-1|]=3$，则 $3\leqslant|x-1|<4$，所以 $3\leqslant x-1<4$ 或 $-4\leqslant x-1<-3$，}
\step{解得 $4\leqslant x<5$ 或 $-3<x\leqslant-2$，}
\step{所以使 $[|x-1|]=3$ 成立的 $x$ 的取值范围是 $(-3,-2]\cup[4,5)$.}
\solend

\fi

\end{enumerate}

\bigsec{二、选择题}

\begin{enumerate}
\setcounter{enumi}{10}

\item 若关于 $x$ 的方程 $k(x+2)-3(k-1)=x+1$ 有负数解，则实数 $k$ 的取值范围是\mbox{（\quad）}.

\keepwithprev
A. $(1,2)$\qquad B. $(-\infty,1)\cup(2,+\infty)$\qquad C. $(-\infty,2)$\qquad D. $(-\infty,1)$

\ans{$A$}

\item 与 $x^2>4$ 同解的不等式是\mbox{（\quad）}.

\keepwithprev
A. $x^2+\dfrac{1}{x-4}>4+\dfrac{1}{x-4}$\qquad B. $x^2+\dfrac1x>4+\dfrac1x$

C. $x^2+\sqrt{x-2}>4+\sqrt{x-2}$\qquad D. $\dfrac{1}{x^2-4}<0$

\ifanswer
\solbegin
\step{不等式 $x^2+\dfrac{1}{x-4}>4+\dfrac{1}{x-4}$ 等价于 $\begin{cases}x^2>4\\ x\neq4\end{cases}$，
与 $x^2>4$ 不同解，故选项 $A$ 错误；}
\step{不等式 $x^2+\dfrac1x>4+\dfrac1x$ 等价于 $\begin{cases}x^2>4\\ x\neq0\end{cases}$，
即 $x^2>4$，所以与 $x^2>4$ 同解，}
\step{故选项 $B$ 正确；}
\step{不等式 $x^2+\sqrt{x-2}>4+\sqrt{x-2}$ 等价于 $\begin{cases}x^2>4\\ x-2\geqslant0\end{cases}$，
与 $x^2>4$ 不同解，}
\step{故选项 $C$ 错误；}
\step{不等式 $\dfrac{1}{x^2-4}<0$ 等价于 $x^2<4$，与 $x^2>4$ 不同解，故选项 $D$ 错误.}
\step{故选 $B$.}
\solend

\fi

\item 若关于 $x$ 的二次不等式 $ax^2+bx+c\geqslant0(a\neq0)$ 的解集是 $\R$，那么\mbox{（\quad）}.

\keepwithprev
A. $a<0$，且 $b^2-4ac>0$\qquad B. $a<0$，且 $b^2-4ac\leqslant0$

% 注：原卷本题 D 选项与 A 选项文字完全重复（均作“$a<0$，且 $b^2-4ac>0$”），
%     照录未改，见文件头“原卷存疑”①。
C. $a>0$，且 $b^2-4ac\leqslant0$\qquad D. $a<0$，且 $b^2-4ac>0$

\ans{$C$}

\item 若不等式 $x^2+mx+1>2x+m$ 对满足 $|m|<2$ 的所有实数 $m$ 恒成立，
则实数 $x$ 的取值范围是\mbox{（\quad）}.

\keepwithprev
A. $-2<x<2$\qquad B. $x\geqslant3$\qquad C. $x\leqslant1$\qquad D. $x\leqslant-1$ 或 $x\geqslant3$

\ans{$D$}

\end{enumerate}

% “三、解答题”原本落在试卷版第 1 页页末、第 15 题翻到次页（切页审计 B 类）。
% 节标题自带的守卫只有 3\baselineskip，够放标题不够放标题＋首题，故在标题前再量一次页高。
\needvspace{4cm}
\bigsec{三、解答题}

\begin{enumerate}
\setcounter{enumi}{14}

\item 设关于 $x$ 的二次方程 $(m-1)x^2+(2m-4)x+m=0$ 有两个正根，求实数 $m$ 的取值范围.

\ifanswer
\solbegin
\step{$\begin{cases}\Delta=(2m-4)^2-4m(m-1)\geqslant0\\ -\dfrac{2m-4}{m-1}>0\\ \dfrac{m}{m-1}>0\end{cases}$，
解得 $1<m\leqslant\dfrac43$.}
\solend

\fi

\ansspace[6]

\item 已知关于 $x$ 的不等式 $2x^2-9x+a<0$ 的解集非空，对于其解集内的每一个 $x$ 的值，
至少能使不等式 $x^2-4x+3<0$ 或 $x^2-6x+8<0$ 中的一个成立，求实数 $a$ 的取值范围.

\ifanswer
\solbegin
\step{不等式 $x^2-4x+3<0$ 的解集为 $(1,3)$；$x^2-6x+8<0$ 的解集为 $(2,4)$.}
\step{关于 $x$ 的不等式 $2x^2-9x+a<0$ 的解集非空，所以 $\Delta=81-8a>0$，解得 $a<\dfrac{81}{8}$.}
\step{所以 $\dfrac{9-\sqrt{81-8a}}{4}<x<\dfrac{9+\sqrt{81-8a}}{4}$.}
\step{由于关于 $x$ 的不等式 $2x^2-9x+a<0$（解集非空）的每一个 $x$ 的值}
\step{至少满足不等式 $x^2-4x+3<0$ 和 $x^2-6x+8<0$ 中的一个，}
\step{所以 $\begin{cases}\dfrac{9-\sqrt{81-8a}}{4}\geqslant1\\ \dfrac{9+\sqrt{81-8a}}{4}\leqslant4\\ a<\dfrac{81}{8}\end{cases}$，
解得 $7\leqslant a<\dfrac{81}{8}$，所以实数 $a$ 的取值范围是 $\left[7,\dfrac{81}{8}\right)$.}
\solend

\fi

\ansspace[8]

\item 是否存在实数 $k$，使 $x=3$ 在不等式 $\dfrac{x+2}{k}>1+\dfrac{x^2-3}{k^2}$ 的解集中？
若存在，求 $k$ 的取值范围；若不存在，说明理由.

\ans{$(2,3)$}

\ansspace[6]

\item 关于 $x$ 的不等式 $x^2-(a+1)x+a<0$ 的解集为 $A$，$x^2-5x+4<0$ 的解集为 $B$，
若 $A\sub B$，求实数 $a$ 的取值范围.

\ans{$[1,4]$}

\ansspace[6]

\item 汽车在行驶中，由于惯性作用，刹车后还要继续向前滑行一段距离才能停住，我们称这段
距离为“刹车距离”. 刹车距离是分析事故的一个重要因素. 在一个限速 $40km/h$ 以内的弯道上，
甲、乙两辆汽车相向而行，发现情况不对，同时刹车，但还是相撞了，事后现场测得甲车的刹车
距离略超过 $12m$，乙车的刹车距离略超过 $10m$，又知甲、乙两种车型的刹车距离 $s(m)$ 与
车速 $x(km/h)$ 之间有如下关系：$s_{\text{甲}}=0.1x+0.01x^2$，$s_{\text{乙}}=0.05x+0.005x^2$. 问：
两车相撞的主要责任是谁？

\ifanswer
\solbegin
\step{由题意列出不等式组 $\begin{cases}0.1x+0.01x^2>12\\ 0.05x+0.005x^2>10\end{cases}$，
得 $\begin{cases}x<-40\text{ 或 }x>30\\ x<-50\text{ 或 }x>40\end{cases}$，}
\step{由于 $x>0$，从而得 $x_{\text{甲}}>30km/h$，$x_{\text{乙}}>40km/h$.}
\step{经比较得乙车超过限速，应负主要责任.}
\solend

\fi

\ansspace*[8]

\end{enumerate}

\end{document}
"
%
% 原始材料是微信公众号图文（公众号“魔都数学加油站”连载“27学年高一 40”，署名“破晓”，
% 2026-09-26 21:29 发布，原文标题“27学年高一40——2026-2027年上海市大同中学高一上周末作业4.2”，
% 原文链接 https://mp.weixin.qq.com/s/RrBSqDF_nJVbbx-bmmdtkQ ），正文为 4 张 PNG 页（4762×6735）。原文本身就是一份**讲评版**——题下直接印【答案】或【解析】，
% 第 11、13、14 题的答案还直接印在题干末尾的括号里。试题文字、答案与解析均照原卷录入，
% 未改内容；卷面的粉色斜水印（“魔都数学加油站”字样）属水印，未录入。
%
% 卷首标题：原卷页首印作“2026-2027 年大同中学高一上周末作业 4.2”（**卷面标题行即带校名**，
% 与高一16／高一18 等拍照卷须由本稿补校名的情形不同），本稿照卷面，年份区间按全稿惯例
% 排 en dash，前缀照原卷保留“年”。
%
% 篇幅与题号：本篇**只有 4 页**，卷面自“一、填空题”的**第 3 题**起编，终于“三、解答题”的
% **第 19 题**，共 17 题（填空 3—10、选择 11—14、解答 15—19）。第 1、2 题未在该文中发布——
% 这是该公众号连载讲评版的通例（高一02—高一09 与 高一36、高一38、高一39 等均自第 3 题起编，
% README“说明”已记），**不是截断**，故本稿题号亦自 3 起。它与 高一39（大同中学 周末作业4.1）
% **不构成同一份试卷的两半**：4.1 是另一次作业（11 页，题号**也自第 3 题**起、一直编到**第 24 题**，
% 以集合／常用逻辑用语为主、兼及不等式），两份各自独立起编（都从第 3 题开始），只是同为“大同中学高一上
% 周末作业 4”系列。
%
% 与原卷的排版差异（均已在 README“说明”中交代）：
%   ① 原文自**第 3 题**起（第 1、2 题未在该文中发布），故本稿题号亦自 3 起；
%   ② 原卷本就印有“一、填空题”“二、选择题”“三、解答题”三节标题，本稿照原卷分节，
%      只把标题改排成全稿统一的“大节标题”样式；
%   ③ 原卷为讲评版：第 4、5、6、7、8、17、18 题只印【答案】行，第 3、9、10、12、15、16、19 题
%      印【解析】（无单独的【答案】行），第 11、13、14 题的答案直接印在题干末尾的括号内；
%      本稿统一改为试卷版隐去、答案版以 \ans{}／\ifanswer 块给出，两版彻底分离（照 高一36 口径：
%      只印【解析】的题不另给 \ans{}，只印【答案】的题仅给 \ans{}）；
%   ④ 数集记号原卷作斜体的 $R$（第 5、10、13 题），本稿按全稿惯例统一写作 $\R$；
%   ⑤ 补集符号原卷作上横线（第 5 题的 $\overline{M}$），本稿照录，未改为 $\complement_\R$；
%   ⑥ 包含符号原卷两套并用：第 4 题作**无下划线**的 $\subset$（真包含），第 18 题作**带下划线**的
%      $\sub$（即 $\subseteq$），本稿照录、不归一；
%   ⑦ 原卷填空的横线约两个字宽，本稿按全稿惯例统一排 \blankline[4.4em]；
%   ⑧ 原卷未印副标题，本稿按**本卷实际内容**补出副标题“不等式”——全卷 17 题里约 15 题为
%      不等式（解一元二次／无理／分式不等式、恒成立与根的分布、含参不等式、应用题），
%      集合仅第 4、5 两题（第 18 题只是用集合关系表述解集），常用逻辑用语 0 题；
%   ⑨ 本卷**无插图**（全卷检索确认无“如图”）；
%   ⑩ 标点与单位：第 6 题的“$a,b$”本稿按全稿惯例排顿号“$a$、$b$”，第 10 题的三例
%      “$[\pi]=3,[-1.2]=-2,[\tfrac12]=0$”之间排顿号，第 5 题原卷“全集 $U=R$ 集合 $M$”之间
%      漏排标点、本稿补逗号；车速单位原卷印作斜体 $km/h$（第 19 题），本稿照录；半角括号、逗号、
%      问号一律排全角（第 17、19 题原卷的两个问号本就作全角，本稿照录）。
%
% 原卷存疑两处（照抄原卷、未改，见源码中该处上方注释）：
%   ① 第 13 题的 D 选项与 A 选项文字完全重复（均作“$a<0$，且 $b^2-4ac>0$”），原卷如此。
%   ② 第 10 题题干末作“…的取值范围____”，**漏排“是”**（本卷其余填空均作“…是____”）——
%      本稿照录、未补字，见源码该题上方 `%` 注。

\documentclass[a4paper,11pt]{article}
\usepackage{common}

\begin{document}

\examtitlest[3]{2026--2027 年大同中学高一上周末作业 4.2}{不等式}

\bigsec{一、填空题}

\begin{enumerate}
\setcounter{enumi}{2}

\item 在平面直角坐标系中，若点 $(2,b)$ 到原点的距离不小于 $5$，则 $b$ 的取值范围是 \blankline[4.4em].

\ifanswer
\solbegin
\step{因为点 $(2,b)$ 到原点的距离不小于 $5$，所以 $\sqrt{(2-0)^2+(b-0)^2}\geqslant5$，}
\step{所以 $4+b^2\geqslant25$，得 $b\geqslant\sqrt{21}$ 或 $b\leqslant-\sqrt{21}$，}
\step{则 $b$ 的取值范围是 $(-\infty,-\sqrt{21}]\cup[\sqrt{21},+\infty)$.}
\solend

\fi

\item 若集合 $A=\{x\mid|x|>3\}$，集合 $B=\{x\mid x\geqslant a\}$，且 $B\subset A$，
则实数 $a$ 的取值范围是 \blankline[4.4em].

\ans{$(3,+\infty)$}

% 注：原卷“全集 $U=R$”与“集合 $M$”之间**漏排标点**（只有空格），本稿补逗号——
%     照 高一b 第 13 题“原卷漏排、本稿补上”的先例（高一10、高一24 同例），见文件头⑩。
\item 若全集 $U=\R$，集合 $M=\{x\mid y=\sqrt{9-x^2}\}$，$S=\{x\mid x^2-3x+2\geqslant0\}$，
则 $\overline{M}\cup S=$ \blankline[4.4em].

\ans{$(-\infty,1]\cup[2,+\infty)$}

\item 设 $a$、$b$ 为实数，若不等式 $x^2-ax-b<0$ 的解集为 $\{x\mid2<x<8\}$，
则 $bx^2-ax-1>0$ 的解集为 \blankline[4.4em].

\ans{$\left(-\dfrac12,-\dfrac18\right)$}

\item 不等式 $4x-13+\sqrt{x+5}\geqslant8x-1+\sqrt{x+5}$ 的解集是 \blankline[4.4em].

\ans{$[-5,-3]$}

\item 若关于 $x$ 的不等式 $ax^2-2ax+3\leqslant0$ 的解集为 $\varnothing$，
则实数 $a$ 的取值范围是 \blankline[4.4em].

\ans{$[0,3)$}

\item 若不等式 $x^2-ax-2<0$ 在 $x\in(1,2)$ 内恒成立，则实数 $a$ 的取值范围是 \blankline[4.4em].

\ifanswer
\solbegin
\step{由不等式 $x^2-ax-2<0$ 在 $x\in(1,2)$ 内恒成立，}
\step{得 $ax>x^2-2$，即 $a>x-\dfrac2x$ 在 $x\in(1,2)$ 内恒成立，}
\step{令 $t(x)=x-\dfrac2x$，$x\in(1,2)$，该函数为严格增函数，则 $t(x)<t(2)=1$，得 $a\geqslant1$，}
\step{所以 $a$ 的取值范围是 $[1,+\infty)$.}
\solend

\fi

% 注：原卷此处作“的取值范围\_\_\_\_”，**漏排“是”**（本卷其余填空均作“…是\_\_\_\_”），
%     本稿照录、未补字，见文件头“原卷存疑”②。
\item 设 $x\in\R$，$[x]$ 表示不大于 $x$ 的最大整数，如 $[\pi]=3$、$[-1.2]=-2$、$\left[\dfrac12\right]=0$，
则使 $[|x-1|]=3$ 成立的 $x$ 的取值范围 \blankline[4.4em].

\ifanswer
\solbegin
\step{由题意 $[|x-1|]=3$，则 $3\leqslant|x-1|<4$，所以 $3\leqslant x-1<4$ 或 $-4\leqslant x-1<-3$，}
\step{解得 $4\leqslant x<5$ 或 $-3<x\leqslant-2$，}
\step{所以使 $[|x-1|]=3$ 成立的 $x$ 的取值范围是 $(-3,-2]\cup[4,5)$.}
\solend

\fi

\end{enumerate}

\bigsec{二、选择题}

\begin{enumerate}
\setcounter{enumi}{10}

\item 若关于 $x$ 的方程 $k(x+2)-3(k-1)=x+1$ 有负数解，则实数 $k$ 的取值范围是\mbox{（\quad）}.

\keepwithprev
A. $(1,2)$\qquad B. $(-\infty,1)\cup(2,+\infty)$\qquad C. $(-\infty,2)$\qquad D. $(-\infty,1)$

\ans{$A$}

\item 与 $x^2>4$ 同解的不等式是\mbox{（\quad）}.

\keepwithprev
A. $x^2+\dfrac{1}{x-4}>4+\dfrac{1}{x-4}$\qquad B. $x^2+\dfrac1x>4+\dfrac1x$

C. $x^2+\sqrt{x-2}>4+\sqrt{x-2}$\qquad D. $\dfrac{1}{x^2-4}<0$

\ifanswer
\solbegin
\step{不等式 $x^2+\dfrac{1}{x-4}>4+\dfrac{1}{x-4}$ 等价于 $\begin{cases}x^2>4\\ x\neq4\end{cases}$，
与 $x^2>4$ 不同解，故选项 $A$ 错误；}
\step{不等式 $x^2+\dfrac1x>4+\dfrac1x$ 等价于 $\begin{cases}x^2>4\\ x\neq0\end{cases}$，
即 $x^2>4$，所以与 $x^2>4$ 同解，}
\step{故选项 $B$ 正确；}
\step{不等式 $x^2+\sqrt{x-2}>4+\sqrt{x-2}$ 等价于 $\begin{cases}x^2>4\\ x-2\geqslant0\end{cases}$，
与 $x^2>4$ 不同解，}
\step{故选项 $C$ 错误；}
\step{不等式 $\dfrac{1}{x^2-4}<0$ 等价于 $x^2<4$，与 $x^2>4$ 不同解，故选项 $D$ 错误.}
\step{故选 $B$.}
\solend

\fi

\item 若关于 $x$ 的二次不等式 $ax^2+bx+c\geqslant0(a\neq0)$ 的解集是 $\R$，那么\mbox{（\quad）}.

\keepwithprev
A. $a<0$，且 $b^2-4ac>0$\qquad B. $a<0$，且 $b^2-4ac\leqslant0$

% 注：原卷本题 D 选项与 A 选项文字完全重复（均作“$a<0$，且 $b^2-4ac>0$”），
%     照录未改，见文件头“原卷存疑”①。
C. $a>0$，且 $b^2-4ac\leqslant0$\qquad D. $a<0$，且 $b^2-4ac>0$

\ans{$C$}

\item 若不等式 $x^2+mx+1>2x+m$ 对满足 $|m|<2$ 的所有实数 $m$ 恒成立，
则实数 $x$ 的取值范围是\mbox{（\quad）}.

\keepwithprev
A. $-2<x<2$\qquad B. $x\geqslant3$\qquad C. $x\leqslant1$\qquad D. $x\leqslant-1$ 或 $x\geqslant3$

\ans{$D$}

\end{enumerate}

% “三、解答题”原本落在试卷版第 1 页页末、第 15 题翻到次页（切页审计 B 类）。
% 节标题自带的守卫只有 3\baselineskip，够放标题不够放标题＋首题，故在标题前再量一次页高。
\needvspace{4cm}
\bigsec{三、解答题}

\begin{enumerate}
\setcounter{enumi}{14}

\item 设关于 $x$ 的二次方程 $(m-1)x^2+(2m-4)x+m=0$ 有两个正根，求实数 $m$ 的取值范围.

\ifanswer
\solbegin
\step{$\begin{cases}\Delta=(2m-4)^2-4m(m-1)\geqslant0\\ -\dfrac{2m-4}{m-1}>0\\ \dfrac{m}{m-1}>0\end{cases}$，
解得 $1<m\leqslant\dfrac43$.}
\solend

\fi

\ansspace[6]

\item 已知关于 $x$ 的不等式 $2x^2-9x+a<0$ 的解集非空，对于其解集内的每一个 $x$ 的值，
至少能使不等式 $x^2-4x+3<0$ 或 $x^2-6x+8<0$ 中的一个成立，求实数 $a$ 的取值范围.

\ifanswer
\solbegin
\step{不等式 $x^2-4x+3<0$ 的解集为 $(1,3)$；$x^2-6x+8<0$ 的解集为 $(2,4)$.}
\step{关于 $x$ 的不等式 $2x^2-9x+a<0$ 的解集非空，所以 $\Delta=81-8a>0$，解得 $a<\dfrac{81}{8}$.}
\step{所以 $\dfrac{9-\sqrt{81-8a}}{4}<x<\dfrac{9+\sqrt{81-8a}}{4}$.}
\step{由于关于 $x$ 的不等式 $2x^2-9x+a<0$（解集非空）的每一个 $x$ 的值}
\step{至少满足不等式 $x^2-4x+3<0$ 和 $x^2-6x+8<0$ 中的一个，}
\step{所以 $\begin{cases}\dfrac{9-\sqrt{81-8a}}{4}\geqslant1\\ \dfrac{9+\sqrt{81-8a}}{4}\leqslant4\\ a<\dfrac{81}{8}\end{cases}$，
解得 $7\leqslant a<\dfrac{81}{8}$，所以实数 $a$ 的取值范围是 $\left[7,\dfrac{81}{8}\right)$.}
\solend

\fi

\ansspace[8]

\item 是否存在实数 $k$，使 $x=3$ 在不等式 $\dfrac{x+2}{k}>1+\dfrac{x^2-3}{k^2}$ 的解集中？
若存在，求 $k$ 的取值范围；若不存在，说明理由.

\ans{$(2,3)$}

\ansspace[6]

\item 关于 $x$ 的不等式 $x^2-(a+1)x+a<0$ 的解集为 $A$，$x^2-5x+4<0$ 的解集为 $B$，
若 $A\sub B$，求实数 $a$ 的取值范围.

\ans{$[1,4]$}

\ansspace[6]

\item 汽车在行驶中，由于惯性作用，刹车后还要继续向前滑行一段距离才能停住，我们称这段
距离为“刹车距离”. 刹车距离是分析事故的一个重要因素. 在一个限速 $40km/h$ 以内的弯道上，
甲、乙两辆汽车相向而行，发现情况不对，同时刹车，但还是相撞了，事后现场测得甲车的刹车
距离略超过 $12m$，乙车的刹车距离略超过 $10m$，又知甲、乙两种车型的刹车距离 $s(m)$ 与
车速 $x(km/h)$ 之间有如下关系：$s_{\text{甲}}=0.1x+0.01x^2$，$s_{\text{乙}}=0.05x+0.005x^2$. 问：
两车相撞的主要责任是谁？

\ifanswer
\solbegin
\step{由题意列出不等式组 $\begin{cases}0.1x+0.01x^2>12\\ 0.05x+0.005x^2>10\end{cases}$，
得 $\begin{cases}x<-40\text{ 或 }x>30\\ x<-50\text{ 或 }x>40\end{cases}$，}
\step{由于 $x>0$，从而得 $x_{\text{甲}}>30km/h$，$x_{\text{乙}}>40km/h$.}
\step{经比较得乙车超过限速，应负主要责任.}
\solend

\fi

\ansspace*[8]

\end{enumerate}

\end{document}
"
%
% 原始材料是微信公众号图文（公众号“魔都数学加油站”连载“27学年高一 40”，署名“破晓”，
% 2026-09-26 21:29 发布，原文标题“27学年高一40——2026-2027年上海市大同中学高一上周末作业4.2”，
% 原文链接 https://mp.weixin.qq.com/s/RrBSqDF_nJVbbx-bmmdtkQ ），正文为 4 张 PNG 页（4762×6735）。原文本身就是一份**讲评版**——题下直接印【答案】或【解析】，
% 第 11、13、14 题的答案还直接印在题干末尾的括号里。试题文字、答案与解析均照原卷录入，
% 未改内容；卷面的粉色斜水印（“魔都数学加油站”字样）属水印，未录入。
%
% 卷首标题：原卷页首印作“2026-2027 年大同中学高一上周末作业 4.2”（**卷面标题行即带校名**，
% 与高一16／高一18 等拍照卷须由本稿补校名的情形不同），本稿照卷面，年份区间按全稿惯例
% 排 en dash，前缀照原卷保留“年”。
%
% 篇幅与题号：本篇**只有 4 页**，卷面自“一、填空题”的**第 3 题**起编，终于“三、解答题”的
% **第 19 题**，共 17 题（填空 3—10、选择 11—14、解答 15—19）。第 1、2 题未在该文中发布——
% 这是该公众号连载讲评版的通例（高一02—高一09 与 高一36、高一38、高一39 等均自第 3 题起编，
% README“说明”已记），**不是截断**，故本稿题号亦自 3 起。它与 高一39（大同中学 周末作业4.1）
% **不构成同一份试卷的两半**：4.1 是另一次作业（11 页，题号**也自第 3 题**起、一直编到**第 24 题**，
% 以集合／常用逻辑用语为主、兼及不等式），两份各自独立起编（都从第 3 题开始），只是同为“大同中学高一上
% 周末作业 4”系列。
%
% 与原卷的排版差异（均已在 README“说明”中交代）：
%   ① 原文自**第 3 题**起（第 1、2 题未在该文中发布），故本稿题号亦自 3 起；
%   ② 原卷本就印有“一、填空题”“二、选择题”“三、解答题”三节标题，本稿照原卷分节，
%      只把标题改排成全稿统一的“大节标题”样式；
%   ③ 原卷为讲评版：第 4、5、6、7、8、17、18 题只印【答案】行，第 3、9、10、12、15、16、19 题
%      印【解析】（无单独的【答案】行），第 11、13、14 题的答案直接印在题干末尾的括号内；
%      本稿统一改为试卷版隐去、答案版以 \ans{}／\ifanswer 块给出，两版彻底分离（照 高一36 口径：
%      只印【解析】的题不另给 \ans{}，只印【答案】的题仅给 \ans{}）；
%   ④ 数集记号原卷作斜体的 $R$（第 5、10、13 题），本稿按全稿惯例统一写作 $\R$；
%   ⑤ 补集符号原卷作上横线（第 5 题的 $\overline{M}$），本稿照录，未改为 $\complement_\R$；
%   ⑥ 包含符号原卷两套并用：第 4 题作**无下划线**的 $\subset$（真包含），第 18 题作**带下划线**的
%      $\sub$（即 $\subseteq$），本稿照录、不归一；
%   ⑦ 原卷填空的横线约两个字宽，本稿按全稿惯例统一排 \blankline[4.4em]；
%   ⑧ 原卷未印副标题，本稿按**本卷实际内容**补出副标题“不等式”——全卷 17 题里约 15 题为
%      不等式（解一元二次／无理／分式不等式、恒成立与根的分布、含参不等式、应用题），
%      集合仅第 4、5 两题（第 18 题只是用集合关系表述解集），常用逻辑用语 0 题；
%   ⑨ 本卷**无插图**（全卷检索确认无“如图”）；
%   ⑩ 标点与单位：第 6 题的“$a,b$”本稿按全稿惯例排顿号“$a$、$b$”，第 10 题的三例
%      “$[\pi]=3,[-1.2]=-2,[\tfrac12]=0$”之间排顿号，第 5 题原卷“全集 $U=R$ 集合 $M$”之间
%      漏排标点、本稿补逗号；车速单位原卷印作斜体 $km/h$（第 19 题），本稿照录；半角括号、逗号、
%      问号一律排全角（第 17、19 题原卷的两个问号本就作全角，本稿照录）。
%
% 原卷存疑两处（照抄原卷、未改，见源码中该处上方注释）：
%   ① 第 13 题的 D 选项与 A 选项文字完全重复（均作“$a<0$，且 $b^2-4ac>0$”），原卷如此。
%   ② 第 10 题题干末作“…的取值范围____”，**漏排“是”**（本卷其余填空均作“…是____”）——
%      本稿照录、未补字，见源码该题上方 `%` 注。

\documentclass[a4paper,11pt]{article}
\usepackage{common}

\begin{document}

\examtitlest[3]{2026--2027 年大同中学高一上周末作业 4.2}{不等式}

\bigsec{一、填空题}

\begin{enumerate}
\setcounter{enumi}{2}

\item 在平面直角坐标系中，若点 $(2,b)$ 到原点的距离不小于 $5$，则 $b$ 的取值范围是 \blankline[4.4em].

\ifanswer
\solbegin
\step{因为点 $(2,b)$ 到原点的距离不小于 $5$，所以 $\sqrt{(2-0)^2+(b-0)^2}\geqslant5$，}
\step{所以 $4+b^2\geqslant25$，得 $b\geqslant\sqrt{21}$ 或 $b\leqslant-\sqrt{21}$，}
\step{则 $b$ 的取值范围是 $(-\infty,-\sqrt{21}]\cup[\sqrt{21},+\infty)$.}
\solend

\fi

\item 若集合 $A=\{x\mid|x|>3\}$，集合 $B=\{x\mid x\geqslant a\}$，且 $B\subset A$，
则实数 $a$ 的取值范围是 \blankline[4.4em].

\ans{$(3,+\infty)$}

% 注：原卷“全集 $U=R$”与“集合 $M$”之间**漏排标点**（只有空格），本稿补逗号——
%     照 高一b 第 13 题“原卷漏排、本稿补上”的先例（高一10、高一24 同例），见文件头⑩。
\item 若全集 $U=\R$，集合 $M=\{x\mid y=\sqrt{9-x^2}\}$，$S=\{x\mid x^2-3x+2\geqslant0\}$，
则 $\overline{M}\cup S=$ \blankline[4.4em].

\ans{$(-\infty,1]\cup[2,+\infty)$}

\item 设 $a$、$b$ 为实数，若不等式 $x^2-ax-b<0$ 的解集为 $\{x\mid2<x<8\}$，
则 $bx^2-ax-1>0$ 的解集为 \blankline[4.4em].

\ans{$\left(-\dfrac12,-\dfrac18\right)$}

\item 不等式 $4x-13+\sqrt{x+5}\geqslant8x-1+\sqrt{x+5}$ 的解集是 \blankline[4.4em].

\ans{$[-5,-3]$}

\item 若关于 $x$ 的不等式 $ax^2-2ax+3\leqslant0$ 的解集为 $\varnothing$，
则实数 $a$ 的取值范围是 \blankline[4.4em].

\ans{$[0,3)$}

\item 若不等式 $x^2-ax-2<0$ 在 $x\in(1,2)$ 内恒成立，则实数 $a$ 的取值范围是 \blankline[4.4em].

\ifanswer
\solbegin
\step{由不等式 $x^2-ax-2<0$ 在 $x\in(1,2)$ 内恒成立，}
\step{得 $ax>x^2-2$，即 $a>x-\dfrac2x$ 在 $x\in(1,2)$ 内恒成立，}
\step{令 $t(x)=x-\dfrac2x$，$x\in(1,2)$，该函数为严格增函数，则 $t(x)<t(2)=1$，得 $a\geqslant1$，}
\step{所以 $a$ 的取值范围是 $[1,+\infty)$.}
\solend

\fi

% 注：原卷此处作“的取值范围\_\_\_\_”，**漏排“是”**（本卷其余填空均作“…是\_\_\_\_”），
%     本稿照录、未补字，见文件头“原卷存疑”②。
\item 设 $x\in\R$，$[x]$ 表示不大于 $x$ 的最大整数，如 $[\pi]=3$、$[-1.2]=-2$、$\left[\dfrac12\right]=0$，
则使 $[|x-1|]=3$ 成立的 $x$ 的取值范围 \blankline[4.4em].

\ifanswer
\solbegin
\step{由题意 $[|x-1|]=3$，则 $3\leqslant|x-1|<4$，所以 $3\leqslant x-1<4$ 或 $-4\leqslant x-1<-3$，}
\step{解得 $4\leqslant x<5$ 或 $-3<x\leqslant-2$，}
\step{所以使 $[|x-1|]=3$ 成立的 $x$ 的取值范围是 $(-3,-2]\cup[4,5)$.}
\solend

\fi

\end{enumerate}

\bigsec{二、选择题}

\begin{enumerate}
\setcounter{enumi}{10}

\item 若关于 $x$ 的方程 $k(x+2)-3(k-1)=x+1$ 有负数解，则实数 $k$ 的取值范围是\mbox{（\quad）}.

\keepwithprev
A. $(1,2)$\qquad B. $(-\infty,1)\cup(2,+\infty)$\qquad C. $(-\infty,2)$\qquad D. $(-\infty,1)$

\ans{$A$}

\item 与 $x^2>4$ 同解的不等式是\mbox{（\quad）}.

\keepwithprev
A. $x^2+\dfrac{1}{x-4}>4+\dfrac{1}{x-4}$\qquad B. $x^2+\dfrac1x>4+\dfrac1x$

C. $x^2+\sqrt{x-2}>4+\sqrt{x-2}$\qquad D. $\dfrac{1}{x^2-4}<0$

\ifanswer
\solbegin
\step{不等式 $x^2+\dfrac{1}{x-4}>4+\dfrac{1}{x-4}$ 等价于 $\begin{cases}x^2>4\\ x\neq4\end{cases}$，
与 $x^2>4$ 不同解，故选项 $A$ 错误；}
\step{不等式 $x^2+\dfrac1x>4+\dfrac1x$ 等价于 $\begin{cases}x^2>4\\ x\neq0\end{cases}$，
即 $x^2>4$，所以与 $x^2>4$ 同解，}
\step{故选项 $B$ 正确；}
\step{不等式 $x^2+\sqrt{x-2}>4+\sqrt{x-2}$ 等价于 $\begin{cases}x^2>4\\ x-2\geqslant0\end{cases}$，
与 $x^2>4$ 不同解，}
\step{故选项 $C$ 错误；}
\step{不等式 $\dfrac{1}{x^2-4}<0$ 等价于 $x^2<4$，与 $x^2>4$ 不同解，故选项 $D$ 错误.}
\step{故选 $B$.}
\solend

\fi

\item 若关于 $x$ 的二次不等式 $ax^2+bx+c\geqslant0(a\neq0)$ 的解集是 $\R$，那么\mbox{（\quad）}.

\keepwithprev
A. $a<0$，且 $b^2-4ac>0$\qquad B. $a<0$，且 $b^2-4ac\leqslant0$

% 注：原卷本题 D 选项与 A 选项文字完全重复（均作“$a<0$，且 $b^2-4ac>0$”），
%     照录未改，见文件头“原卷存疑”①。
C. $a>0$，且 $b^2-4ac\leqslant0$\qquad D. $a<0$，且 $b^2-4ac>0$

\ans{$C$}

\item 若不等式 $x^2+mx+1>2x+m$ 对满足 $|m|<2$ 的所有实数 $m$ 恒成立，
则实数 $x$ 的取值范围是\mbox{（\quad）}.

\keepwithprev
A. $-2<x<2$\qquad B. $x\geqslant3$\qquad C. $x\leqslant1$\qquad D. $x\leqslant-1$ 或 $x\geqslant3$

\ans{$D$}

\end{enumerate}

% “三、解答题”原本落在试卷版第 1 页页末、第 15 题翻到次页（切页审计 B 类）。
% 节标题自带的守卫只有 3\baselineskip，够放标题不够放标题＋首题，故在标题前再量一次页高。
\needvspace{4cm}
\bigsec{三、解答题}

\begin{enumerate}
\setcounter{enumi}{14}

\item 设关于 $x$ 的二次方程 $(m-1)x^2+(2m-4)x+m=0$ 有两个正根，求实数 $m$ 的取值范围.

\ifanswer
\solbegin
\step{$\begin{cases}\Delta=(2m-4)^2-4m(m-1)\geqslant0\\ -\dfrac{2m-4}{m-1}>0\\ \dfrac{m}{m-1}>0\end{cases}$，
解得 $1<m\leqslant\dfrac43$.}
\solend

\fi

\ansspace[6]

\item 已知关于 $x$ 的不等式 $2x^2-9x+a<0$ 的解集非空，对于其解集内的每一个 $x$ 的值，
至少能使不等式 $x^2-4x+3<0$ 或 $x^2-6x+8<0$ 中的一个成立，求实数 $a$ 的取值范围.

\ifanswer
\solbegin
\step{不等式 $x^2-4x+3<0$ 的解集为 $(1,3)$；$x^2-6x+8<0$ 的解集为 $(2,4)$.}
\step{关于 $x$ 的不等式 $2x^2-9x+a<0$ 的解集非空，所以 $\Delta=81-8a>0$，解得 $a<\dfrac{81}{8}$.}
\step{所以 $\dfrac{9-\sqrt{81-8a}}{4}<x<\dfrac{9+\sqrt{81-8a}}{4}$.}
\step{由于关于 $x$ 的不等式 $2x^2-9x+a<0$（解集非空）的每一个 $x$ 的值}
\step{至少满足不等式 $x^2-4x+3<0$ 和 $x^2-6x+8<0$ 中的一个，}
\step{所以 $\begin{cases}\dfrac{9-\sqrt{81-8a}}{4}\geqslant1\\ \dfrac{9+\sqrt{81-8a}}{4}\leqslant4\\ a<\dfrac{81}{8}\end{cases}$，
解得 $7\leqslant a<\dfrac{81}{8}$，所以实数 $a$ 的取值范围是 $\left[7,\dfrac{81}{8}\right)$.}
\solend

\fi

\ansspace[8]

\item 是否存在实数 $k$，使 $x=3$ 在不等式 $\dfrac{x+2}{k}>1+\dfrac{x^2-3}{k^2}$ 的解集中？
若存在，求 $k$ 的取值范围；若不存在，说明理由.

\ans{$(2,3)$}

\ansspace[6]

\item 关于 $x$ 的不等式 $x^2-(a+1)x+a<0$ 的解集为 $A$，$x^2-5x+4<0$ 的解集为 $B$，
若 $A\sub B$，求实数 $a$ 的取值范围.

\ans{$[1,4]$}

\ansspace[6]

\item 汽车在行驶中，由于惯性作用，刹车后还要继续向前滑行一段距离才能停住，我们称这段
距离为“刹车距离”. 刹车距离是分析事故的一个重要因素. 在一个限速 $40km/h$ 以内的弯道上，
甲、乙两辆汽车相向而行，发现情况不对，同时刹车，但还是相撞了，事后现场测得甲车的刹车
距离略超过 $12m$，乙车的刹车距离略超过 $10m$，又知甲、乙两种车型的刹车距离 $s(m)$ 与
车速 $x(km/h)$ 之间有如下关系：$s_{\text{甲}}=0.1x+0.01x^2$，$s_{\text{乙}}=0.05x+0.005x^2$. 问：
两车相撞的主要责任是谁？

\ifanswer
\solbegin
\step{由题意列出不等式组 $\begin{cases}0.1x+0.01x^2>12\\ 0.05x+0.005x^2>10\end{cases}$，
得 $\begin{cases}x<-40\text{ 或 }x>30\\ x<-50\text{ 或 }x>40\end{cases}$，}
\step{由于 $x>0$，从而得 $x_{\text{甲}}>30km/h$，$x_{\text{乙}}>40km/h$.}
\step{经比较得乙车超过限速，应负主要责任.}
\solend

\fi

\ansspace*[8]

\end{enumerate}

\end{document}
